\chapter{Composite kernels and the odd two-body kernels}

In the construction every one-loop component of a local operator along a generator $x$ of $\mathfrak{psl}(4|4)$ is read from the defect flow of Lemma \ref{lem:c8-transfer}: it is the sum over the pairs of letters of the coefficients of $\varepsilon^0$ of the defect cycles, cycles through the two letters with one special edge and one vertex $S^{(1)}(x)$, minus the bracket of the operator with the arity-one finite counterterm $\hbar F_1^{(1)}(x)$, which solves the arity-one equation of the action with all its free directions. This master formula (Theorem \ref{thm:c11-classes}) needs an argument: the graphs with the special edge on a bridge have finite limits, which cancel exactly against the special-edge parts of their counterterm partners. The chapter first treats the composite term of Chapter 10, the second evaluation of $D_1$ and $P_1$ through the scaling of the regulator, in which $S^{(1)}(\Delta)$ or $S^{(1)}(P_b)$ sits on the cycle as a line from a ghost leg to a gauge leg (Theorem \ref{thm:c10-D1}, Proposition \ref{prop:c10-P1}); on the letters $A|_V$, $\lambda$, $c$ the renormalised composite two-point functional contains $\log|p|^2$ of the momentum (Remark \ref{rem:c9-Pi}), and the locality of the composite term is proved: $Q$ maps the composite graphs of a cocycle to cycles with one special edge, and the composite term is a sum over the pairs of letters of the composite pair kernels. These kernels are simplex integrals, evaluated on twist two, the $\mathfrak{su}(1|1)$ sector, the fermion and gluon $\mathfrak{sl}(2)$ sectors and two mixed families; with the sector coefficients $-\tfrac13$ units of $\hbar F_1^{(1)}(\Delta)$ computed from the action-level arity-two equations (Theorem \ref{thm:c9-theta2}), $D_1$ is four times the map of Theorem \ref{thm:c10-brst}, with spectrum twice the one-loop anomalous dimensions, and $P_1$ vanishes on cohomology in the translation-covariant scheme; the operator-level arity-two equation reproduces the value $-\tfrac13$ family by family. The last section derives the master formula along every generator, computes the action-level obstructions and solutions along the odd generators and along $K_b$ with their free directions and the arity-two equations that relate them, assembles $Q_1$, $S_1$, $D_1$, $P_1$ and $K_1$ on the scalar bilinears as representatives together with their scheme-independent content, and compares with the one-loop differential of the holomorphic twist at the level of classes. Chapter 12 states the construction for every generator and every arity together (Theorems \ref{thm:c12-master}, \ref{thm:c12-ar2} and \ref{thm:c12-F0}).

The unit is $\frac{\hbar\operatorname{tr}_{\mathfrak g}(\mathrm{ad}_X\mathrm{ad}_Y)}{16\pi^2(X,Y)}$, equal to $-\frac{g^2N}{32\pi^2}$ for $\mathfrak g=\mathfrak{su}(N)$ (Remark \ref{rem:c6-konishi}). Edges carry $P(\varepsilon,T)$ of Definition \ref{def:c8-PT}; the special edge is $K_\varepsilon-K_0=-\int_0^\varepsilon(H\otimes1)K_t\,dt$, the factor $e^{-\varepsilon|p|^2}-1$ times the components of $K_0$ (Proposition \ref{prop:c4-defect1}). Every component below is independent of the endpoint $T\in(0,\infty]$ of the transfer (Theorem \ref{thm:c8-L} (ii)) and is evaluated at $T=\infty$, where the graph sums converge at generic momenta (Proposition \ref{prop:c8-Tsafe}); along $K_b$ it is taken at a finite $T$ (Definition \ref{def:c8-K}). On arity zero $s=Q+\{I,-\}$ (Theorem \ref{thm:c8-obstruction}). Letters, the pair weight, the pair degree and the Koszul signs are those of Definition \ref{def:c10-letters} and \eqref{eq:c10-decomp}. The field-counting functional is $2\int(A,A^*)-4\int(\lambda^*,\lambda)-\int(c^*,c)$, whose bracket multiplies the coordinate functionals of $c$, $A$, $\lambda$, $A^*$, $\lambda^*$, $c^*$ by $1$, $2$, $4$, $-2$, $-4$, $-1$ (Lemma \ref{lem:c6-count}), and the map of Theorem \ref{thm:c10-brst} is the coefficient of $\log\varepsilon$ in $\partial_\delta I^{CT}[\varepsilon](O)$ minus the unit times the bracket of $O$ with it (Section 10.7). Along $\Delta$, $v=2E$ and $v.$ is the operator of Lemma \ref{lem:c9-SDelta}, so that $\{S^{(1)}(\Delta),M\}=v.M$; along $P_b$, $v.$ is $\partial_b$ on every slot (Proposition \ref{prop:c7-cov}); for $x\in\{\Delta,P_b\}$, $v_x.$ is the operator of the vector field of $x$ (Definition \ref{def:c7-vdot}). The two-component form of Definition \ref{def:c8-twocomp} of every main formula is in Remark \ref{rem:c11-twocomp}.

\section{The composite term as a sum over pairs}

\begin{definition}[Square kernel and composite kernels]
\label{def:c11-kernels}
For $0<\varepsilon<T\le\infty$ and $P=P(\varepsilon,T)$, the square kernel is $\big((P\star)^2\otimes1\big)K_0$, and the composite kernel along $x\in\{\Delta,P_b\}$ is
\[
\big(P\star\,v_x.\otimes1+1\otimes P\star\,v_x.\big)P,
\]
the two legs of $S^{(1)}(x)$ joined, in both orders, to two edges $P$ (Definition \ref{def:c9-composite}, Theorem \ref{thm:c9-P1}). The sector parts of $S^{(1)}(P_b)$ are
\[
\int\big(\partial_b(A|_{W_{\mathbb R}^\perp}),A^*|_{W_{\mathbb R}^\perp}\big),\qquad -\int(\partial_b\lambda,\lambda^*),\qquad \int\big(\partial_b(A|_V),A^*|_V\big)-\int(\partial_bc,c^*);
\]
their sum is $S^{(1)}(P_b)$, and each sector part of $S^{(1)}(\Delta)$ of Section 10.8 is $2\int x^\mu$ times the density of the sector part of $S^{(1)}(P_{e_\mu})$ of the same sector, summed over $\mu$, plus a translation-invariant functional, with $e_\mu$ an orthonormal basis of $V$ and $x^\mu$ the coordinates.
\end{definition}

\begin{lemma}[Structure of the composite kernels]
\label{lem:c11-kernels}
\begin{itemize}
\item (i) The composite kernel along $x$ is $2\big(P\star\,v_x.\,P\star\otimes1\big)K_0=(v_x.\otimes1-1\otimes v_x.)$ applied to the square kernel.
\item (ii) $(v_x.\otimes1+1\otimes v_x.)$ applied to the square kernel is $-4\big(\varepsilon\partial_\varepsilon+T\partial_T\big)$ of the square kernel for $x=\Delta$ and $0$ for $x=P_b$. Hence the composite kernel along $\Delta$ is $2(v.\otimes1)$ of the square kernel plus $4(\varepsilon\partial_\varepsilon+T\partial_T)$ of it, and the composite kernel along $P_b$ is $2(\partial_b\otimes1)$ of the square kernel.
\item (iii) The square kernel and the composite kernels have components in $c\otimes A|_V$ and $A|_V\otimes c$ alone: as operators, $(P\star)^2$ and $P\star\,v_x.\,P\star$ map the slot of $A^*|_V$ to the slot of $c$ and the slot of $c^*$ to the slot of $A|_V$, and vanish on the other slots.
\item (iv) In a graph the composite kernel is the vertex $S^{(1)}(x)$ with both legs on edges $P$. Its term $\int(v_x.A|_V,A^*|_V)$ contributes, with the leg $A|_V$ on the component $A\otimes A$ of $P$ and the leg $A^*|_V$ on the component $A^*|_V\otimes c$. Along $\Delta$ this internal vertex is
\[
\int\big((2E+2)A|_V,A^*|_V\big)=2\int x^\mu\big(\partial_\mu(A|_V),A^*|_V\big)+2\int\big(A|_V,A^*|_V\big),
\]
along $P_b$ it is $\int(\partial_b(A|_V),A^*|_V)$. The components of $P$ lie in $A\otimes A$, $\lambda\otimes\lambda$, $c\otimes A^*|_V$ and $A^*|_V\otimes c$ (Proposition \ref{prop:c3-P}), so the terms of $S^{(1)}(x)$ in the sectors of $A|_{W_{\mathbb R}^\perp}$, $\lambda$ and $c$, whose antifield legs are $A^*|_{W_{\mathbb R}^\perp}$, $\lambda^*$, $c^*$, enter a cycle through the special edge.
\item (v) $(Q\otimes1+1\otimes Q)$ annihilates the composite kernel; the operator $\partial$ of the composite kernel is odd and anticommutes with $Q$.
\item (vi) At $T=\infty$ the component of the square kernel is $(Q^{GF})^2$ on the slot of $A^*|_V$, which is $2d^*$ with values in the slot of $c$ (Proposition \ref{prop:c9-composite} (ii)), times $(-,-)^{-1}$ and
\[
\int_\varepsilon^\infty\!\!\int_\varepsilon^\infty e^{-(s'+s'')|p|^2}\,ds'\,ds''=e^{-2\varepsilon|p|^2}|p|^{-4}:
\]
a doubled edge with two Schwinger times $s',s''\in[\varepsilon,\infty)$, of momentum degree $-3$, which obeys the weight rule of Proposition \ref{prop:c3-P} for the slots $c$ (weight $0$) and $A|_V$ (weight $1$).
\end{itemize}
\end{lemma}

\begin{proof}
(i) The first equality is Proposition \ref{prop:c9-composite} (i) along $\Delta$; along $P_b$ the same computation holds, $\partial_b$ preserving the pairing and commuting with $Q^{GF}$ and $H$ (Proposition \ref{prop:c7-cov}). For the second, $2P\star\,v_x.\,P\star=v_x.(P\star)^2+(P\star)^2v_x.$, along $\Delta$ because $[v.,P\star]$ is $Q^{GF}$ times a function of $H$ and commutes with $P\star$ (proof of Proposition \ref{prop:c9-composite} (i)), along $P_b$ because $\partial_b$ commutes with $P\star$. For the even operator $X=(P\star)^2$, $(X\,v_x.\otimes1)K_0=(X\otimes1)(v_x.\otimes1)K_0=-(1\otimes v_x.)(X\otimes1)K_0$ by Lemma \ref{lem:c9-ops}, $v_x.$ being even and pairing preserving.

(ii) Since $(v_x.\otimes1+1\otimes v_x.)K_0=0$, the left side is $([v_x.,X]\otimes1)K_0$. Along $\Delta$, $X=(Q^{GF})^2g(H)^2$ with $g(H)=\int_\varepsilon^Te^{-tH}dt$, $[v.,Q^{GF}]=-4Q^{GF}$ and $[v.,f(H)]=-4Hf'(H)$ (Proposition \ref{prop:c7-dil}). From $Hg'(H)=Te^{-TH}-\varepsilon e^{-\varepsilon H}-g=(\varepsilon\partial_\varepsilon+T\partial_T)g-g$,
\[
[v.,X]=-8X-4(Q^{GF})^2H\partial_H\big(g^2\big)=-8X-4\big((\varepsilon\partial_\varepsilon+T\partial_T)X-2X\big)=-4\big(\varepsilon\partial_\varepsilon+T\partial_T\big)X .
\]
Along $P_b$, $[\partial_b,X]=0$. Insert into (i).

(iii) $(Q^{GF})^2$ vanishes on the slots of $c$, $A$, $\lambda$, $\lambda^*$; on $A^*$ and $c^*$ it is $2d^*(A^*|_V)$ in the slot of $c$ and $-2dc^*$ in the slot of $A|_V$ (Proposition \ref{prop:c9-composite} (ii)); $v_x.$ preserves the slots, and the components $A^*\otimes A$ and $c^*\otimes c$ of $K_0$ go to $c\otimes A|_V$ and $A|_V\otimes c$.

(iv) Applied to $A^*\otimes A$, the kernel $(P\star\,v_x.\,P\star\otimes1)K_0$ is the component $A^*\to A$ of $Q^{GF}$, then $v_x.$ on the slot of $A$, then the component $A\to c$ of $Q^{GF}$: the vertex $\int(v_x.A|_V,A^*|_V)$ between an $A\otimes A$ and an $A^*|_V\otimes c$ component, the factor $2$ counting the two orientations. On $A|_V$, $L_{2E}A|_V=2x^\mu\partial_\mu(A|_V)+2A|_V$ (Definition \ref{def:c7-vdot}).

(v) Proposition \ref{prop:c9-composite} (iii) along $\Delta$; along $P_b$, $QP\star+P\star Q=e^{-\varepsilon H}-e^{-TH}$ commutes with $\partial_b$, which commutes with $Q$, so $Q(P\star\,\partial_b\,P\star)=(P\star\,\partial_b\,P\star)Q$. The identity $[\partial_\phi,Q]=\partial_{(Q\otimes1+1\otimes Q)\phi}$ of Lemma \ref{lem:c9-ops} gives the anticommutation; the composite kernel has degree $-1$.

(vi) Proposition \ref{prop:c9-composite} (iv) and $\big(\int_\varepsilon^\infty e^{-s|p|^2}ds\big)^2=e^{-2\varepsilon|p|^2}|p|^{-4}$.
\end{proof}

\begin{definition}[Composite graphs and the arity-one counterterm of an operator]
\label{def:c11-graphs}
Let $M$ be a local functional, $\delta$ a square-zero parameter of degree $-|M|$, $x\in\{\Delta,P_b\}$ and $P=P(\varepsilon,\infty)$. The composite graphs of $M$ along $x$ are the $\delta$-linear part of
\[
\tfrac\hbar2\,\partial_{\text{composite kernel along }x}\,W(P,I+\delta M)_{\rm tree}:
\]
the composite kernel closes two legs of a tree of vertices $I$ and one vertex $M$. The cycle so formed either passes through two letters $X_a$, $X_b$ of $M$, a composite cycle through the pair $(X_a,X_b)$ with trees of $I$ on its outputs, or passes through vertices of $I$ alone and is joined to $M$ by a tree through one letter, a composite loop of $I$ on a leg of $M$. The arity-one counterterm of $M$ along $x$ is $\hbar\,\mathrm{Sing}_\varepsilon$ of the amputated composite cycles through the pairs, summed over the pairs with the Koszul signs of \eqref{eq:c10-decomp}. The singular part of the composite loops of $I$ is $I^{CT}_{\mathfrak{psl}(4|4)}[\varepsilon]^{(1)}(x)$ (Proposition \ref{prop:c9-structure} (i)); its quadratic part along $\Delta$ is the unit times $\log\varepsilon$ times $Q\big(\int(\iota_xA|_V,c^*)-\tfrac12\int(\rho(x)\lambda^*,\lambda^*)\big)$, with no term in $\varepsilon^{-1}$ (Theorem \ref{thm:c9-quadratic} (i)).
\end{definition}

\begin{proposition}[The image of the composite graphs under $Q$]
\label{prop:c11-master}
For every local functional $M$, $x\in\{\Delta,P_b\}$ and $\varepsilon>0$, at $T=\infty$,
\[
Q\big(\text{composite graphs of }M\text{ along }x\big)=-\tfrac\hbar2\,\partial_{\text{composite kernel along }x}\Big[\text{tree-level flow, }\delta\text{-linear, with one vertex }sM\text{ or one special edge}\Big],
\]
the special edge joining a vertex $I$ to a letter of $M$, from $\{I,M\}_\varepsilon-\{I,M\}$, or two vertices $I$, from $\tfrac12(\{I,I\}_\varepsilon-\{I,I\})$. Since $\Delta_\varepsilon-\Delta_0=-\partial_{K_\varepsilon-K_0}$ (Definition \ref{def:c1-bracket}, Proposition \ref{prop:c4-defect1}), $\{f,g\}_\varepsilon-\{f,g\}$ is minus $f$ and $g$ joined by one edge $K_\varepsilon-K_0$; for $sM=0$ the right side is $\tfrac\hbar2\,\partial_{\text{composite kernel along }x}$ of the trees with exactly one edge $K_\varepsilon-K_0$.
\end{proposition}

\begin{proof}
By Lemma \ref{lem:c11-kernels} (v), $Q$ anticommutes with $\partial$ of the composite kernel. By Proposition \ref{prop:c4-defectflow} at order $\hbar^0$, $Q\,W(P,J)_{\rm tree}$ is the tree-level flow of $QJ+\tfrac12\{J,J\}_\varepsilon$; the term $\hbar\Delta_\varepsilon J$ is of order $\hbar$, and its composite closure has two cycles. For $J=I+\delta M$ the classical master equation $QI+\tfrac12\{I,I\}=0$ (Proposition \ref{prop:c1-cme}) leaves $\tfrac12(\{I,I\}_\varepsilon-\{I,I\})$ in the $\delta$-independent part and $sM+(\{I,M\}_\varepsilon-\{I,M\})$ in the $\delta$-linear part.
\end{proof}

\begin{definition}[Special-edge composite cycles and the composite pair kernel]
\label{def:c11-sp}
For two letters $X_a$, $X_b$ of any slots and $x\in\{\Delta,P_b\}$, the special-edge composite cycles of $(X_a,X_b)$ along $x$ are
\[
\tfrac\hbar2\,\partial_{\text{composite kernel along }x}\,\partial_{K_\varepsilon-K_0}\big(\text{amputated chains from }X_a\text{ to }X_b\big):
\]
the chains of Definition \ref{def:c10-chain} with $k\ge1$ vertices of $I$ in which one edge $e_c$ is the composite kernel along $x$, one other edge $e_\sigma$ is the special edge, with the components of $K_0$ (each field slot with its antifield slot), which may end at a letter of any slot, and every remaining edge is $P(\varepsilon,\infty)$. By Proposition \ref{prop:c11-master} they are minus $\tfrac\hbar2\,\partial$ of the composite kernel applied to the chains with one bracket-defect vertex. A chain is read with the signs of Lemma \ref{lem:c10-sign} along the chain from $X_a$, times
\[
(-1)^{n(<\sigma)+n(<c)+[c<\sigma]},
\]
where $n(<j)$ is the number of odd outputs at the vertices before the edge $e_j$, and $[c<\sigma]$ is $1$ if $e_c$ precedes $e_\sigma$ and $0$ otherwise; the outputs are left open as in Definition \ref{def:c10-pair}. The composite pair kernel of $(X_a,X_b)$ along $x$ is the coefficient of $\varepsilon^0$ of its special-edge composite cycles.
\end{definition}

By Proposition \ref{prop:c11-master}, for $sM=0$ the image under $Q$ of the composite cycles through the pairs and of the composite loops on the legs is the sum of the special-edge composite cycles with trees of $I$ on their outputs, together with the composite closures in which the special edge lies on a tree; Theorem \ref{thm:c11-local} identifies the composite term of a cocycle with the sum over its pairs of letters of the composite pair kernels.

\begin{proposition}[Reachability and census]
\label{prop:c11-census}
\begin{itemize}
\item (i) The composite kernel joins a slot $c$ to a slot $A|_V$. The letters $c$ and $A|_V$ reach it directly; the letters $A|_{W_{\mathbb R}^\perp}$, $\lambda$ and $A^*|_V$ reach it through one vertex, $A|_V(A|_{W_{\mathbb R}^\perp})^2$ entered at $A|_{W_{\mathbb R}^\perp}$, $\lambda|_{S_+\otimes W}\lambda|_{S_-\otimes W^*}A|_V$ entered at $\lambda$, $([A|_V,c],A^*|_V)$ entered at $c$, each left at $A|_V$; the letters $A^*|_{W_{\mathbb R}^\perp}$, $\lambda^*$, $c^*$, on which $P$ and the composite kernel have no component, reach it through the special edge, to a slot $A|_{W_{\mathbb R}^\perp}$, $\lambda$, $c$.
\item (ii) A chain of Definition \ref{def:c11-sp} exists for $45$ of the $55$ unordered pairs of the ten slots; the other ten are the pairs of two letters among $A^*|_{W_{\mathbb R}^\perp}$, $\lambda^*|_{S_+\otimes W^*}$, $\lambda^*|_{S_-\otimes W}$, $c^*$, which would need two special edges. Every chain conserves the degree, the fermion parity and the centre $\mathbb Z/4$ of $SL(4,\mathbb C)$, and its outputs carry the pair degree; the composite kernel has degree $-1$ and the special edge $+1$.
\item (iii) Every chain has $k\ge1$ vertices. A chain with $k=1$ has the composite kernel at one letter and the special edge at the other, so $k=1$ occurs exactly for the pairs with a letter $A|_V$ and for the pairs $(c,A^*|_V)$, $(c,A^*|_{W_{\mathbb R}^\perp})$, $(c,\lambda^*)$, $(c,c^*)$; the other pairs, the field pairs $(A|_{W_{\mathbb R}^\perp},A|_{W_{\mathbb R}^\perp})$, $(A|_{W_{\mathbb R}^\perp},\lambda)$, $(\lambda,\lambda)$ among them, start at $k=2$.
\item (iv) At the pair weight without derivatives the leg sets of the ten pairs of field letters are those of Table \ref{tab:c11-fieldpairs}; the numbers of leg sets with $n\le3$ derivatives on the two letters, for the $21$ pairs of the active slots of Lemma \ref{lem:c10-active}, are those of Table \ref{tab:c11-counts}.
\end{itemize}
\end{proposition}

\begin{proof}
(i) Lemma \ref{lem:c11-kernels} (iii), Lemma \ref{lem:c10-active} and the vertices of Definition \ref{def:c1-vertex}. (ii), (iv) Enumeration of the chains with one composite kernel and one special edge, the vertices of $I$ and the components of $P$ and $K_0$, up to six vertices and output weight four above the pair weight; every chain is checked for degree, parity and the centre. A pair of two letters among $A^*|_{W_{\mathbb R}^\perp}$, $\lambda^*$, $c^*$ needs an edge with a component on each letter; $K_0$ alone has one, and a chain has one special edge. (iii) A chain with $k=0$ would join the two letters by the composite kernel and the special edge at once; the rest is read off from the enumeration.
\end{proof}

\begin{tabular}{l|c|l}
pair & pair weight & leg sets of the special-edge composite cycles, fewest vertices in brackets; ${}^\dagger$: position part along $\Delta$ and part along $P_b$\\
\hline
$(A|_V,A|_V)$ & $2$ & $A|_V$ $(1)$, $A|_V^2$ $(2)$, $(A|_{W_{\mathbb R}^\perp})^2$ $(2)$, $A|_V^3{}^\dagger$ $(2)$, $A|_V(A|_{W_{\mathbb R}^\perp})^2{}^\dagger$ $(2)$, $A^*|_V\,c^\dagger$ $(2)$, $\lambda|_{S_-\otimes W^*}\lambda|_{S_+\otimes W}{}^\dagger$ $(2)$\\
$(A|_V,A|_{W_{\mathbb R}^\perp})$ & $2$ & $A|_{W_{\mathbb R}^\perp}$ $(1)$, $A|_V\,A|_{W_{\mathbb R}^\perp}$ $(2)$, $A|_V^2A|_{W_{\mathbb R}^\perp}{}^\dagger$ $(2)$, $(A|_{W_{\mathbb R}^\perp})^3{}^\dagger$ $(2)$, $A^*|_{W_{\mathbb R}^\perp}c^\dagger$ $(2)$, $(\lambda|_{S_-\otimes W^*})^2{}^\dagger$ $(2)$, $(\lambda|_{S_+\otimes W})^2{}^\dagger$ $(2)$\\
$(A|_V,\lambda|_{S_+\otimes W})$ & $\tfrac52$ & $\lambda|_{S_+\otimes W}$ $(1)$, $A|_V\,\lambda|_{S_+\otimes W}$ $(2)$, $A|_{W_{\mathbb R}^\perp}\lambda|_{S_-\otimes W^*}$ $(2)$, $c\,\lambda^*|_{S_-\otimes W}$ $(2)$, $A|_V^2\lambda|_{S_+\otimes W}{}^\dagger$ $(2)$, $A|_V\,A|_{W_{\mathbb R}^\perp}\lambda|_{S_-\otimes W^*}{}^\dagger$ $(3)$, $A|_V\,c\,\lambda^*|_{S_-\otimes W}{}^\dagger$ $(3)$, $(A|_{W_{\mathbb R}^\perp})^2\lambda|_{S_+\otimes W}{}^\dagger$ $(2)$, $A|_{W_{\mathbb R}^\perp}c\,\lambda^*|_{S_+\otimes W^*}{}^\dagger$ $(3)$\\
$(A|_V,\lambda|_{S_-\otimes W^*})$ & $\tfrac52$ & $\lambda|_{S_-\otimes W^*}$ $(1)$, $A|_V\,\lambda|_{S_-\otimes W^*}$ $(2)$, $A|_{W_{\mathbb R}^\perp}\lambda|_{S_+\otimes W}$ $(2)$, $c\,\lambda^*|_{S_+\otimes W^*}$ $(2)$, $A|_V^2\lambda|_{S_-\otimes W^*}{}^\dagger$ $(2)$, $A|_V\,A|_{W_{\mathbb R}^\perp}\lambda|_{S_+\otimes W}{}^\dagger$ $(3)$, $A|_V\,c\,\lambda^*|_{S_+\otimes W^*}{}^\dagger$ $(3)$, $(A|_{W_{\mathbb R}^\perp})^2\lambda|_{S_-\otimes W^*}{}^\dagger$ $(2)$, $A|_{W_{\mathbb R}^\perp}c\,\lambda^*|_{S_-\otimes W}{}^\dagger$ $(3)$\\
$(A|_{W_{\mathbb R}^\perp},A|_{W_{\mathbb R}^\perp})$ & $2$ & $(A|_{W_{\mathbb R}^\perp})^2$ $(2)$, $A|_V(A|_{W_{\mathbb R}^\perp})^2{}^\dagger$ $(2)$\\
$(A|_{W_{\mathbb R}^\perp},\lambda|_{S_+\otimes W})$ & $\tfrac52$ & $A|_{W_{\mathbb R}^\perp}\lambda|_{S_+\otimes W}$ $(2)$, $A|_V\,A|_{W_{\mathbb R}^\perp}\lambda|_{S_+\otimes W}{}^\dagger$ $(2)$, $(A|_{W_{\mathbb R}^\perp})^2\lambda|_{S_-\otimes W^*}{}^\dagger$ $(3)$, $A|_{W_{\mathbb R}^\perp}c\,\lambda^*|_{S_-\otimes W}{}^\dagger$ $(3)$\\
$(A|_{W_{\mathbb R}^\perp},\lambda|_{S_-\otimes W^*})$ & $\tfrac52$ & $A|_{W_{\mathbb R}^\perp}\lambda|_{S_-\otimes W^*}$ $(2)$, $A|_V\,A|_{W_{\mathbb R}^\perp}\lambda|_{S_-\otimes W^*}{}^\dagger$ $(2)$, $(A|_{W_{\mathbb R}^\perp})^2\lambda|_{S_+\otimes W}{}^\dagger$ $(3)$, $A|_{W_{\mathbb R}^\perp}c\,\lambda^*|_{S_+\otimes W^*}{}^\dagger$ $(3)$\\
$(\lambda|_{S_+\otimes W},\lambda|_{S_+\otimes W})$ & $3$ & $(\lambda|_{S_+\otimes W})^2$ $(2)$, $A|_V(\lambda|_{S_+\otimes W})^2{}^\dagger$ $(3)$, $A|_{W_{\mathbb R}^\perp}\lambda|_{S_-\otimes W^*}\lambda|_{S_+\otimes W}{}^\dagger$ $(3)$, $c\,\lambda|_{S_+\otimes W}\lambda^*|_{S_-\otimes W}{}^\dagger$ $(3)$\\
$(\lambda|_{S_+\otimes W},\lambda|_{S_-\otimes W^*})$ & $3$ & $\lambda|_{S_-\otimes W^*}\lambda|_{S_+\otimes W}$ $(2)$, $A|_V\,\lambda|_{S_-\otimes W^*}\lambda|_{S_+\otimes W}{}^\dagger$ $(3)$, $A|_{W_{\mathbb R}^\perp}(\lambda|_{S_-\otimes W^*})^2{}^\dagger$ $(3)$, $A|_{W_{\mathbb R}^\perp}(\lambda|_{S_+\otimes W})^2{}^\dagger$ $(3)$, $c\,\lambda|_{S_-\otimes W^*}\lambda^*|_{S_-\otimes W}{}^\dagger$ $(3)$, $c\,\lambda|_{S_+\otimes W}\lambda^*|_{S_+\otimes W^*}{}^\dagger$ $(3)$\\
$(\lambda|_{S_-\otimes W^*},\lambda|_{S_-\otimes W^*})$ & $3$ & $(\lambda|_{S_-\otimes W^*})^2$ $(2)$, $A|_V(\lambda|_{S_-\otimes W^*})^2{}^\dagger$ $(3)$, $A|_{W_{\mathbb R}^\perp}\lambda|_{S_-\otimes W^*}\lambda|_{S_+\otimes W}{}^\dagger$ $(3)$, $c\,\lambda|_{S_-\otimes W^*}\lambda^*|_{S_+\otimes W^*}{}^\dagger$ $(3)$\\
\end{tabular}
\tabcaption{The leg sets of the special-edge composite cycles through the ten pairs of field letters at the pair weight without derivatives. The leg sets without ${}^\dagger$ enter the translation-invariant part along $\Delta$ as well.\label{tab:c11-fieldpairs}}

\begin{tabular}{l|c|c|c|c}
pair & $n=0$ & $n=1$ & $n=2$ & $n=3$\\
\hline
$(A|_V,A|_V)$ & $7\,(3)$ & $18\,(7)$ & $34\,(18)$ & $67\,(34)$\\
$(A|_V,A|_{W_{\mathbb R}^\perp})$ & $7\,(2)$ & $16\,(7)$ & $32\,(16)$ & $59\,(32)$\\
$(A|_V,\lambda|_{S_+\otimes W})$ & $9\,(4)$ & $20\,(9)$ & $42\,(20)$ & $79\,(42)$\\
$(A|_V,\lambda|_{S_-\otimes W^*})$ & $9\,(4)$ & $20\,(9)$ & $42\,(20)$ & $79\,(42)$\\
$(A|_V,c)$ & $4\,(2)$ & $8\,(4)$ & $19\,(8)$ & $35\,(19)$\\
$(A|_V,A^*|_V)$ & $15\,(6)$ & $35\,(15)$ & $73\,(35)$ & $142\,(73)$\\
$(A|_{W_{\mathbb R}^\perp},A|_{W_{\mathbb R}^\perp})$ & $2\,(1)$ & $7\,(2)$ & $16\,(7)$ & $38\,(16)$\\
$(A|_{W_{\mathbb R}^\perp},\lambda|_{S_+\otimes W})$ & $4\,(1)$ & $12\,(4)$ & $30\,(12)$ & $62\,(30)$\\
$(A|_{W_{\mathbb R}^\perp},\lambda|_{S_-\otimes W^*})$ & $4\,(1)$ & $12\,(4)$ & $30\,(12)$ & $62\,(30)$\\
$(A|_{W_{\mathbb R}^\perp},c)$ & $2\,(1)$ & $7\,(2)$ & $16\,(7)$ & $32\,(16)$\\
$(A|_{W_{\mathbb R}^\perp},A^*|_V)$ & $10\,(2)$ & $26\,(10)$ & $64\,(26)$ & $127\,(64)$\\
$(\lambda|_{S_+\otimes W},\lambda|_{S_+\otimes W})$ & $4\,(1)$ & $12\,(4)$ & $28\,(12)$ & $63\,(28)$\\
$(\lambda|_{S_+\otimes W},\lambda|_{S_-\otimes W^*})$ & $6\,(1)$ & $15\,(6)$ & $36\,(15)$ & $75\,(36)$\\
$(\lambda|_{S_+\otimes W},c)$ & $4\,(1)$ & $9\,(4)$ & $20\,(9)$ & $42\,(20)$\\
$(\lambda|_{S_+\otimes W},A^*|_V)$ & $17\,(5)$ & $42\,(17)$ & $88\,(42)$ & $173\,(88)$\\
$(\lambda|_{S_-\otimes W^*},\lambda|_{S_-\otimes W^*})$ & $4\,(1)$ & $12\,(4)$ & $28\,(12)$ & $63\,(28)$\\
$(\lambda|_{S_-\otimes W^*},c)$ & $4\,(1)$ & $9\,(4)$ & $20\,(9)$ & $42\,(20)$\\
$(\lambda|_{S_-\otimes W^*},A^*|_V)$ & $17\,(5)$ & $42\,(17)$ & $88\,(42)$ & $173\,(88)$\\
$(c,c)$ & $2\,(1)$ & $4\,(2)$ & $8\,(4)$ & $19\,(8)$\\
$(c,A^*|_V)$ & $18\,(7)$ & $34\,(18)$ & $67\,(34)$ & $121\,(67)$\\
$(A^*|_V,A^*|_V)$ & $18\,(5)$ & $52\,(18)$ & $114\,(52)$ & $235\,(114)$\\
\end{tabular}
\tabcaption{The number of leg sets of the special-edge composite cycles entering the part along $P_b$ and, in brackets, the translation-invariant part along $\Delta$, for $n$ derivatives on the two letters.\label{tab:c11-counts}}

\begin{proposition}[Output bound]
\label{prop:c11-bound}
\begin{itemize}
\item (i) A chain of Definition \ref{def:c11-sp} with outputs of total Weyl weight $w$ has superficial degree $w_{X_a}+w_{X_b}-w$ before the insertion factor of $S^{(1)}(x)$; the factor $\partial_b$ adds one, and the translation-invariant part along $\Delta$ adds none.
\item (ii) The composite pair kernel along $P_b$ is a translation-invariant local functional of Weyl weight $w_{X_a}+w_{X_b}+1$; the composite pair kernel along $\Delta$ is local of Weyl weight $w_{X_a}+w_{X_b}$, the position counted at weight $-1$, at most linear in the position. Both have the pair degree and the parity $|X_a|+|X_b|$.
\item (iii) Chains with $w\le w_{X_a}+w_{X_b}+1$ contribute to the kernel along $P_b$ and to the part along $\Delta$ linear in the position, chains with $w\le w_{X_a}+w_{X_b}$ to the translation-invariant part along $\Delta$; a chain above these bounds has a vanishing limit as $\varepsilon\to0$.
\item (iv) The number of output fields is at most $w_{X_a}+w_{X_b}+1$ plus the pair degree for both kernels; the coefficients of $\varepsilon^{-m}$ of the special-edge composite cycles have Weyl weight $w_{X_a}+w_{X_b}-2m$ along $\Delta$ and $w_{X_a}+w_{X_b}+1-2m$ along $P_b$.
\end{itemize}
The same bounds hold for the composite cycles through the pair and for the arity-one counterterm, of degree one less than the pair degree: the coefficient of $\log\varepsilon$ of the arity-one counterterm of $X_aX_b$ has Weyl weight $w_{X_a}+w_{X_b}$ along $\Delta$ and $w_{X_a}+w_{X_b}+1$ along $P_b$.
\end{proposition}

\begin{proof}
(i) Every edge obeys the weight rule $w_1+w_2-4$: $P$ (Proposition \ref{prop:c3-P}), $K_0$ (a field and its antifield have the weights $w$ and $4-w$, momentum degree $0$) and the composite kernel (Lemma \ref{lem:c11-kernels} (vi)); the count of Lemma \ref{lem:c10-degree} applies. (ii) The coefficient of $\varepsilon^0$ is a polynomial in the output momenta of degree equal to the superficial degree (Theorem \ref{thm:c11-simplex}); along $\Delta$ the position enters once, through the factor $2x^\mu$ of Lemma \ref{lem:c11-kernels} (iv). Degree and parity: Proposition \ref{prop:c11-census} (ii). (iii) A polynomial of negative degree vanishes, and so does the limit of a chain of negative degree with a special edge (Theorem \ref{thm:c11-simplex}). (iv) Weyl weight plus degree is at least $1$ for every output (Proposition \ref{prop:c10-bound}), so the number of outputs is at most $w$ plus the pair degree. The statements on the arity-one counterterm follow in the same way from Lemma \ref{lem:c4-local}.
\end{proof}

\section{Closed forms and locality}

\begin{definition}[Position part and translation-invariant part]
\label{def:c11-canonical}
A local functional is in canonical form if it is written $\int\omega\,D$ with $\omega\in\Omega^4_c(\mathbb R^4)$ and $D$ a polynomial in letters, every derivative on a field. A map $M$ of local functionals, linear in $\omega$, of the form
\[
M\Big(\int\omega\,D\Big)=\int\omega\,\big(x^\mu M_\mu(D)+M_0(D)\big),
\]
with $M_\mu$, $M_0$ maps of densities at one point and $x^\mu$ the coordinates of $V$ in an orthonormal basis, has the position part $b^\mu M_\mu$ along $b\in V$ and the translation-invariant part $M_0$. A local functional $G=\int\big(x^\mu g_\mu+g_0\big)$ at most linear in the position has the position part $\int b^\mu g_\mu$ along $b$ and the translation-invariant part $\int g_0$. On plane waves, with the momenta of the fields independent and $\omega$ absorbing the total momentum, a graph is $\alpha\,\delta(\cdot)+\beta^\mu\partial_\mu\delta(\cdot)$, and $M_0$ corresponds to $\alpha$, $M_\mu$ to $i\beta^\mu$.
\end{definition}

\begin{lemma}[Translations and the canonical form]
\label{lem:c11-canonical}
$M$ is well defined, independently of integrations by parts against $\omega$, exactly when every $M_\mu$ commutes with the derivatives and $[M_0,\partial_\mu]=M_\mu$. Then $[M,\{S^{(1)}(P_b),-\}]$ is the map $\int\omega\,D\mapsto\int\omega\,b^\mu M_\mu(D)$. For $M=v.=\{S^{(1)}(\Delta),-\}$ the position part is $2\partial_b$ and the translation-invariant part is the multiplication by $2w_D$; $\{S^{(1)}(P_b),-\}$ has position part $0$.
\end{lemma}

\begin{proof}
$M(\int\omega\,\partial_\mu D)$ and $M(-\int(\partial_\mu\omega)D)=\int\omega\big(M_\mu(D)+x^\nu\partial_\mu M_\nu(D)+\partial_\mu M_0(D)\big)$ agree for all $\omega$ exactly under the two conditions. Then $M(\int\omega\,\partial_bD)-\{S^{(1)}(P_b),M(\int\omega\,D)\}=\int\omega\,\big(M_0(\partial_bD)-\partial_bM_0(D)\big)=\int\omega\,b^\mu M_\mu(D)$, since $\{S^{(1)}(P_b),-\}$ acts by $\partial_b$ on the fields and leaves the position and $\omega$ untouched. Lemma \ref{lem:c9-SDelta} (i) gives $v.\int\omega\,D=\int\omega\,2(x\cdot\partial+w_D)D$.
\end{proof}

\begin{proposition}[Position decomposition]
\label{prop:c11-position}
\begin{itemize}
\item (i) For every pair of letters, the position part along $b$ of the composite pair kernel along $\Delta$ is twice the composite pair kernel along $P_b$. For every local functional the position part of its arity-one counterterm along $\Delta$ is twice its arity-one counterterm along $P_b$; on cocycles the position part of the composite term along $\Delta$ is twice the composite term along $P_b$; and the position part of $I^{CT}_{\mathfrak{psl}(4|4)}[\varepsilon]^{(1)}(\Delta)$ along $b$ is $2I^{CT}_{\mathfrak{psl}(4|4)}[\varepsilon]^{(1)}(P_b)$. In particular the quadratic part of the coefficient of $\log\varepsilon$ in $I^{CT}_{\mathfrak{psl}(4|4)}[\varepsilon]^{(1)}(P_b)$ is $\tfrac12$ unit times $Q\big(\int(\iota_bA|_V,c^*)-\tfrac12\int(\rho(b)\lambda^*,\lambda^*)\big)$, and the quadratic part of the arity-one obstruction along $P_b$ vanishes.
\item (ii) The translation-invariant part of the composite pair kernel along $\Delta$ is computed with the internal vertex $\int\big(2(x-x_O)^\mu\partial_\mu(A|_V)+2A|_V,A^*|_V\big)$, $x_O$ the point of the density, and pair by pair its commutator with $\{S^{(1)}(P_b),-\}$ is twice the composite pair kernel along $P_b$.
\end{itemize}
\end{proposition}

\begin{proof}
(i) By Lemma \ref{lem:c11-kernels} (iv) the internal vertex along $\Delta$ is $2\int x^\mu s_\mu+\int s_0$, with $\int s_\mu$ the internal vertex along $P_{e_\mu}$ and $s_0=2(A|_V,A^*|_V)$; insert $x^\mu=x_O^\mu+(x-x_O)^\mu$. On plane waves the factor $x^\mu$ is $-i\partial/\partial r^\mu$ of $e^{i(r,x)}$, with the momentum $r$ injected at the internal vertex and absorbed by $\omega$; its action on the momentum delta function is the position part, $-i$ times the graph with $\int s_\mu$ at $r=0$, which is $2\cdot(-i)$ times the graph along $P_{e_\mu}$; its action on the integrand is translation invariant. The same holds graph by graph for the composite cycles, their singular parts, the counterterm trees and the composite loops of $I$. Theorem \ref{thm:c9-quadratic} (i) gives the quadratic part along $\Delta$, which is $Q$ of $\int x^\mu$ times a density with $Q$ mapping letters to letters; Theorem \ref{thm:c9-quadratic} (ii) gives a quadratic obstruction along $\Delta$ without explicit position. (ii) The first statement is the decomposition of (i); the second is Lemma \ref{lem:c11-canonical} applied to the composite pair kernel along $\Delta$, whose position part is given by (i).
\end{proof}

\begin{theorem}[Special-edge composite cycles as simplex integrals]
\label{thm:c11-simplex}
For a chain of Definition \ref{def:c11-sp} with $k$ vertices of $I$ and edges $e_0,\dots,e_k$, $e_0$ at $X_a$, let $e_c$ be the composite kernel and $e_\sigma$, $\sigma\ne c$, the special edge. The cycle has $k+2$ Schwinger times: $t_j\in[\varepsilon,\infty)$ for $j\notin\{c,\sigma\}$, $t_\sigma\in[0,\varepsilon]$, and $t'_c,t''_c\in[\varepsilon,\infty)$ for the two segments of the composite kernel (Lemma \ref{lem:c11-kernels} (vi)), which meet at the internal vertex of $S^{(1)}(x)$. The edge $e_j$ carries $l+q_j$ as in Definition \ref{def:c10-cycle}, both segments carry $l+q_c$, and along $\Delta$ a momentum $r$ is injected at the internal vertex and carried by the second segment and by $e_{c+1},\dots,e_k$ to $X_b$ and $\omega$. The numerator is the product of the vertex functions, the edge factors of Definition \ref{def:c10-cycle}, the special-edge factor $-|l+q_\sigma|^2$ times the fibre factor of $K_0$, the segment factors $2(-,-)^{-1}$ times the inverse metric and $-i(l+q_c)(-,-)^{-1}$, and the insertion: $i(l+q_c)_b$ along $P_b$; $2$ for $2\int(A|_V,A^*|_V)$ and $2i(l+q_c)_\mu$ for $2\int x^\mu(\partial_\mu(A|_V),A^*|_V)$ along $\Delta$. With the simplex of dimension $k+1$ of the coordinates $\alpha=(\alpha_j)_{j\ne c}\cup(\alpha'_c,\alpha''_c)$, its centre, momentum form and $m$-th simplex integrand as in Definition \ref{def:c10-cycle} with $k+1$ in place of $k$, at generic output momenta and $T=\infty$,
\begin{equation}
\begin{aligned}
16\pi^2\big(\text{coefficient of }\varepsilon^{-m}\big)&=\frac1m\int_{\alpha_\sigma\text{ smallest}}\Big[\big(\min_{f\ne\sigma}\alpha_f\big)^m-\alpha_\sigma^m\Big]\big(m\text{-th simplex integrand}\big)\,d\alpha\qquad(m\ge1),\\
16\pi^2\big(\text{coefficient of }\varepsilon^0\big)&=\int_{\alpha_\sigma\text{ smallest}}\big(0\text{-th simplex integrand}\big)\,\log\frac{\min_{f\ne\sigma}\alpha_f}{\alpha_\sigma}\,d\alpha,
\end{aligned}
\label{eq:c11-simplex}
\end{equation}
the coefficient of $\log\varepsilon$ is $0$, and the remainder tends to $0$. Each coefficient is a polynomial in the output momenta and in $r$, of degree the superficial degree minus $2m$, and the limit vanishes for negative superficial degree. The coefficients of $\varepsilon^0$ are integrals of polynomials times logarithms of ratios of affine functions over rational subsimplices; for $k=1$ they lie in $\mathbb Q+\mathbb Q\log2+\mathbb Q\log3$ (Proposition \ref{prop:c11-k1}).
\end{theorem}

\begin{proof}
The domain $\{t_\sigma\le\varepsilon,\ t_f\ge\varepsilon\ (f\ne\sigma)\}$ is the difference of the domains with the cutoff shapes $1/\min_{f\ne\sigma}\alpha_f$ and $1/\min_f\alpha_f$, which differ on the region where $\alpha_\sigma$ is the smallest coordinate. Proposition \ref{prop:c10-shape} applies to each with $k+2$ edges, the Schwinger measure being $\tau^{k+1}d\tau\,d\alpha$, and the difference is taken as in Lemma \ref{lem:c10-special}: the two terms in $\log\varepsilon$ cancel, and the finite parts differ by $-\int F\log(\alpha_\sigma/\min_{f\ne\sigma}\alpha_f)$ over that region, $F$ the $0$-th simplex integrand. The region $\tau\ge1$ converges at $T=\infty$: the composite kernel is of order $|p|^{-3}$ at small momentum (Proposition \ref{prop:c9-composite} (iv)) and the other edges carry distinct generic momenta (proof of Theorem \ref{thm:c10-decomp} (ii)), so the integrand is integrable and the measure of $\{\alpha_\sigma\le\varepsilon/\tau\}$ is of order $\varepsilon$; the injected momentum enters polynomially at each order in $r$. The degree statement is Proposition \ref{prop:c11-bound} (i).
\end{proof}

\begin{proposition}[Generating functions in the point splitting]
\label{prop:c11-genfun}
With the letters at $x+y_a$ and $x+y_b$, $y=y_b-y_a$, the $m$-th simplex integrand for the numerator $(\text{chain numerator})\,e^{i(y,l)}$ is the chain numerator with $l$ replaced by $-i\partial_y$, applied to
\[
(-1)^{m+k}e^{i(y,\text{centre})}\sum_{i\ge0}\frac{(|y|^2/4)^{i+m+k}(\text{momentum form})^i}{i!\,(i+m+k)!}=(-1)^{m+k}e^{i(y,\text{centre})}\Big(\frac{|y|}{2\sqrt{\text{momentum form}}}\Big)^{m+k}I_{m+k}\big(|y|\sqrt{\text{momentum form}}\big),
\]
and, summed over the chains with the signs of Definition \ref{def:c11-sp},
\[
\sum_{\alpha,\beta}\frac{y_a^\alpha y_b^\beta}{\alpha!\,\beta!}\big(\text{composite pair kernel of }(\partial^\alpha\varphi_a,\partial^\beta\varphi_b)\text{ along }x\big)=\frac{\hbar}{16\pi^2}\,e^{i(y_b,q_k)}\int_{\alpha_\sigma\text{ smallest}}\big(0\text{-th simplex integrand}\big)\log\frac{\min_{f\ne\sigma}\alpha_f}{\alpha_\sigma}\,d\alpha,
\]
with $q_k$ including $r$ along $\Delta$; the coefficients of $\varepsilon^{-m}$ have the same generating function with the weights of \eqref{eq:c11-simplex}. In position space $iq_j$ is the derivative of the product of the outputs at $x$.
\end{proposition}

\begin{proof}
As Proposition \ref{prop:c10-genfun}, with $k+2$ edges: the power $\tau^{-1-m}$ of the Gaussian-integrated integrand times $\tau^{k+1}$ fixes the index shift $m+k$, and Theorem \ref{thm:c11-simplex} replaces the weight $(\min_e\alpha_e)^m$ by the weights of the region where $\alpha_\sigma$ is the smallest.
\end{proof}

\begin{proposition}[One vertex in closed form]
\label{prop:c11-k1}
Let $k=1$, $\beta:=\alpha_\sigma$, and let $\gamma'$ and $\gamma''=1-\beta-\gamma'$ be the coordinates of the two segments. At $r=0$ the centre $-\beta q_\sigma-(1-\beta)q_c$ and the momentum form $\beta(1-\beta)|q_c-q_\sigma|^2$ depend on $\beta$ alone, the region where $\beta$ is the smallest coordinate requires $\beta\le\tfrac13$, and
\[
\begin{aligned}
16\pi^2\big(\text{composite pair kernel}\big)&=\int_0^{1/3}\big(0\text{-th simplex integrand}\big)(\beta)\Big[(1-\beta)\log\frac{1-\beta}{2\beta}-(1-3\beta)\Big]d\beta,\\
16\pi^2\big(\text{coefficient of }\varepsilon^{-m}\big)&=\int_0^{1/3}\big(m\text{-th simplex integrand}\big)(\beta)\,\frac2m\Big[\frac{\big(\tfrac{1-\beta}2\big)^{m+1}-\beta^{m+1}}{m+1}-\beta^m\Big(\frac{1-\beta}2-\beta\Big)\Big]d\beta .
\end{aligned}
\]
The part linear in $r$ carries one factor $\gamma''$ in the simplex integrand, and $\int\gamma''(\cdots)\,d\gamma'=\tfrac{1-\beta}2\int(\cdots)\,d\gamma'$ by the symmetry $\gamma'\leftrightarrow\gamma''$ of the region and of the weights. The coefficients are integrals of $\beta^a(1-\beta)^b$ times these weights over $[0,\tfrac13]$, in $\mathbb Q+\mathbb Q\log2+\mathbb Q\log3$.
\end{proposition}

\begin{proof}
On the region, $\min_{f\ne\sigma}\alpha_f=\min(\gamma',\gamma'')$ and $\gamma'\in[\beta,1-2\beta]$. By the symmetry about $\gamma'=\tfrac{1-\beta}2$,
\[
\int\log\frac{\min(\gamma',\gamma'')}\beta\,d\gamma'=2\int_\beta^{(1-\beta)/2}\log\frac{\gamma'}\beta\,d\gamma'=(1-\beta)\log\frac{1-\beta}{2\beta}-(1-3\beta),
\]
and $\frac1m\int\big(\min(\gamma',\gamma'')^m-\beta^m\big)d\gamma'$ is the second weight. At $r=0$ both segments carry $l+q_c$, so the Gaussian data depend on $\beta$ alone; along $\Delta$ the derivative in $r$ of $\sum_jt_j|l+q_j|^2$ is $2t''_c(l+q_c)\cdot$, linear in $\gamma''$. The integrals $\int_0^{1/3}\beta^a(1-\beta)^b\log\frac{1-\beta}{2\beta}\,d\beta$ are rational combinations of $1$, $\log2$, $\log3$, the endpoint $\beta=\tfrac13$ giving $\log\tfrac23$.
\end{proof}

\begin{definition}[Composite loops of the action]
\label{def:c11-Xi}
For $x\in\{\Delta,P_b\}$ and $P=P(\varepsilon,T)$, the composite loops of $I$ along $x$, written $\Xi_x(\varepsilon)$, are the amputated one-particle-irreducible cycles of vertices of $I$ and edges $P$ in which one edge is the composite kernel along $x$, all other legs open: $\tfrac\hbar2\,\partial_{\text{composite kernel along }x}W(P,I)_{\rm tree}$ is the tree-level flow of $\Xi_x(\varepsilon)$. $\Xi^\sigma_x(\varepsilon)$ is the sum of the same loops in which one edge of the cycle other than the composite kernel is the special edge, with the sign of Proposition \ref{prop:c11-master}, minus the loops with one edge $K_\varepsilon-K_0$. The renormalised composite functional is $\Xi^{\rm ren}_x:=\lim_{\varepsilon\to0}\big(\Xi_x(\varepsilon)-I^{CT}_{\mathfrak{psl}(4|4)}[\varepsilon]^{(1)}(x)\big)$, a functional of degree $-1$ at generic momenta, whose quadratic part along $\Delta$ is the renormalised composite two-point functional of Remark \ref{rem:c9-Pi}. The composite graphs of a local functional $O$ (Definition \ref{def:c11-graphs}) are of the first type when the cycle passes through two letters of $O$ and of the second type when it passes through vertices of $I$ alone; the second type is the tree insertion of $\Xi_x(\varepsilon)$ into $O$, written $\mathrm{tr}(\Xi_x,O)$.
\end{definition}

\begin{lemma}[Images of the composite loops]
\label{lem:c11-Cid}
For every $\varepsilon>0$, as functionals at generic momenta, $s\,\Xi_x(\varepsilon)=\Xi^\sigma_x(\varepsilon)$. Consequently $\Xi^\sigma_x$ has no term in $\log\varepsilon$ and no poles, $I^{CT}_{\mathfrak{psl}(4|4)}[\varepsilon]^{(1)}(x)=\log\varepsilon\,L^{(1)}(x)$ with $L^{(1)}(x)$ an $s$-closed local functional, and the composite part along $x$ is $s\,\Xi^{\rm ren}_x$, the coefficient of $\varepsilon^0$ of $\Xi^\sigma_x$, a local functional.
\end{lemma}

\begin{proof}
Compute $Q$ of the tree-level flow of $\Xi_x$ in two ways. First, $\tfrac\hbar2\,\partial_{\text{composite kernel}}W(P,I)_{\rm tree}$ is that flow, and $\partial$ of the composite kernel anticommutes with $Q$ (Lemma \ref{lem:c11-kernels} (v)); by Proposition \ref{prop:c4-defectflow} at order $\hbar^0$ and the classical master equation, $Q$ of it is $-\tfrac\hbar2\,\partial$ of the composite kernel applied to the tree-level flow of $\tfrac12\{I,I\}_\varepsilon-\tfrac12\{I,I\}$, the composite closures of the trees of $I$ with one special edge (Proposition \ref{prop:c11-master} with $M=0$). Sort them by the position of the special edge: on the cycle closed by the composite kernel, which gives the tree-level flow of $\Xi^\sigma_x$, or on a tree, between two vertices of $I$ or between $\Xi_x$ and a vertex of $I$. Second, with $\Xi_x$ as a vertex, Proposition \ref{prop:c4-defectflow} applied to $I+\delta\,\Xi_x$, $\delta$ of degree $+1$, gives the tree-level flow of $s\,\Xi_x$ plus the tree-level flows with the special-edge parts of $\{I,\Xi_x\}_\varepsilon$ and of $\tfrac12\{I,I\}_\varepsilon$ on tree edges. The graphs with the special edge on a tree are the same in both evaluations, with the same signs, the sign rule of Proposition \ref{prop:c11-master} being the Koszul rule of the second evaluation. Hence the tree-level flows of $s\,\Xi_x$ and of $\Xi^\sigma_x$ agree, and the tree-level flow is injective on functionals graded by the number of legs (the argument of Theorem \ref{thm:c8-L} (i)). $\Xi^\sigma_x$ is a sum of cycles with one special edge and the composite kernel, which by the argument of Theorem \ref{thm:c11-simplex} are $\sum_m\varepsilon^{-m}a_m+a_0+o(1)$ without $\log\varepsilon$. The poles of $I^{CT}_{\mathfrak{psl}(4|4)}[\varepsilon]^{(1)}(x)$ vanish by the weight count of Proposition \ref{prop:c9-structure} (i), and along $P_b$ by the same count at Weyl weight $5-2m$; so $I^{CT}_{\mathfrak{psl}(4|4)}[\varepsilon]^{(1)}(x)=\log\varepsilon\,L^{(1)}(x)$, and the coefficient of $\log\varepsilon$ in $s\,\Xi_x=s\,I^{CT}_{\mathfrak{psl}(4|4)}[\varepsilon]^{(1)}(x)+s\,\Xi^{\rm ren}_x+o(1)$ gives $s\,L^{(1)}(x)=0$, while the poles of $\Xi^\sigma_x$ vanish. The coefficient of $\varepsilon^0$ gives the last statement.
\end{proof}

\begin{lemma}[Special edges on trees of the composite graphs]
\label{lem:c11-trees}
Let $O$ be a local functional. The composite graphs of the first type with the special edge on a tree, at an output of the cycle or at a spectator letter of $O$, together with the special-edge parts of the counterterm trees of the arity-one counterterm of $O$, are $e^{-\varepsilon|p|^2}-1$ at a generic momentum $p$ times renormalised composite cycles, and tend to $0$. The composite graphs of the second type with the special edge on a tree tend to $0$, and so do the special-edge parts of the counterterm trees of $I^{CT}_{\mathfrak{psl}(4|4)}[\varepsilon]^{(1)}(x)$.
\end{lemma}

\begin{proof}
A special edge on a tree edge of momentum $p$ factorises: the graph is $(e^{-\varepsilon|p|^2}-1)$ times the composite cycle through the pair with its trees; the counterterm vertex joined by the same special edge gives $(e^{-\varepsilon|p|^2}-1)$ times minus its singular part, so the sum is $(e^{-\varepsilon|p|^2}-1)$ times the renormalised composite cycle, which has a limit at generic momenta (Proposition \ref{prop:c9-structure} (i)). At a spectator letter of $O$ the same holds with the special-edge parts of $\{I,O\}_\varepsilon$ and of the bracket of $I$ with the arity-one counterterm of $O$. The graphs of the second type are $(e^{-\varepsilon|p|^2}-1)$ times $\Xi_x(\varepsilon)$, which is $\log\varepsilon\,L^{(1)}(x)$ plus a convergent functional (Lemma \ref{lem:c11-Cid}), and $(e^{-\varepsilon|p|^2}-1)\log\varepsilon\to0$; the counterterm trees carry $(e^{-\varepsilon|p|^2}-1)\log\varepsilon$.
\end{proof}

\begin{proposition}[The image of the composite graphs of each type]
\label{prop:c11-typeI}
For every local functional $O$ and every $\varepsilon>0$, with (A) the special-edge composite cycles of the pairs of letters of $O$ with trees of $I$ on their outputs (Definition \ref{def:c11-sp}), (B) the composite graphs of the first type with the special edge on a tree, (D) those of the second type with the special edge on a tree, and $\epsilon_\Xi\in\{\pm1\}$ the sign with which the bracket $\{\Xi_x,O\}$ enters $Q\,\mathrm{tr}(\Xi_x,O)$,
\[
\begin{aligned}
Q\big(\text{composite graphs of the first type of }O\big)&=(\mathrm A)+(\mathrm B)-\epsilon_\Xi\{\Xi_x(\varepsilon),O\}-\big(\text{composite graphs of }sO\big)\mp\mathrm{tr}(\Xi_x,sO),\\
Q\,\mathrm{tr}(\Xi_x,O)&=\mathrm{tr}(s\,\Xi_x,O)+\epsilon_\Xi\{\Xi_x,O\}\pm\mathrm{tr}(\Xi_x,sO)+(\mathrm D),
\end{aligned}
\]
the terms with $sO$ being present only for $sO\neq0$.
\end{proposition}

\begin{proof}
By Proposition \ref{prop:c11-master} the image under $Q$ of all composite graphs of $O$ is $-\tfrac\hbar2\,\partial$ of the composite kernel applied to the trees with one vertex $sO$ or one special edge; the trees with one vertex $sO$ closed by the composite kernel are the composite graphs of $sO$, of both types. The special edge lies on the cycle through two letters (A), on a tree of a graph of the first type (B), on the cycle of a composite loop of $I$, which gives $\mathrm{tr}(\Xi^\sigma_x,O)$, or on a tree of a graph of the second type (D). The image of the second type is computed with $\Xi_x$ as a vertex, by Proposition \ref{prop:c4-defectflow} applied to $I+\delta_1\Xi_x+\delta_2O$, the part linear in $\delta_1$ and $\delta_2$: the vertex $s\,\Xi_x$ with $O$, the vertex $sO$ with $\Xi_x$, the direct contraction $\{\Xi_x,O\}$, and the special edges on trees, which are (D). Subtract: $\mathrm{tr}(s\,\Xi_x,O)=\mathrm{tr}(\Xi^\sigma_x,O)$ by Lemma \ref{lem:c11-Cid}, so the loops with the special edge on the cycle cancel, (D) cancels, and $\pm\mathrm{tr}(\Xi_x,sO)$ moves to the first line with the opposite sign.
\end{proof}

\begin{theorem}[The composite term is two-body and local]
\label{thm:c11-local}
Let $O=\int\omega\,X_1\cdots X_n$ be a local functional with $sO=0$, $x\in\{\Delta,P_b\}$, $T=\infty$, generic momenta, and $\hbar F_1^{(0)}$ as in (i) of Definition \ref{def:c10-hyp}.
\begin{itemize}
\item (i) The special-edge composite cycles of $(X_a,X_b)$ along $x$ are $\sum_{m\ge1}\varepsilon^{-m}$ times local coefficients, plus the composite pair kernel, plus a remainder tending to $0$; they contain no $\log\varepsilon$.
\item (ii) $s$ annihilates the coefficient of $\log\varepsilon$ in $I^{CT}_{\mathfrak{psl}(4|4)}[\varepsilon]^{(1)}(x)$, and $I^{CT}_{\mathfrak{psl}(4|4)}[\varepsilon]^{(1)}(x)$ is $\log\varepsilon$ times that coefficient in every sector.
\item (iii) (Arity-one BRST covariance.) $s$ of the coefficient of $\log\varepsilon$ in the arity-one counterterm of $O$ along $x$ is minus the bracket of $O$ with the coefficient of $\log\varepsilon$ in $I^{CT}_{\mathfrak{psl}(4|4)}[\varepsilon]^{(1)}(x)$, and the sum over the pairs, with the Koszul signs of \eqref{eq:c10-decomp}, of the coefficients of $\varepsilon^{-m}$ of the special-edge composite cycles is $s$ of the coefficient of $\varepsilon^{-m}$ in the arity-one counterterm of $O$.
\item (iv) The composite term of $O$ along $\Delta$ (Definition \ref{def:c10-hyp}) and along $P_b$ (Proposition \ref{prop:c10-P1}) is
\begin{equation}
\big(\text{composite term of }O\text{ along }x\big)=\sum_{a<b}(-1)^{|X_b|(|X_{a+1}|+\cdots+|X_{b-1}|)}\int\omega\,X_1\cdots X_{a-1}\,\big(\text{composite pair kernel of }(X_a,X_b)\text{ along }x\big)\,X_{a+1}\cdots\widehat{X_b}\cdots X_n,
\label{eq:c11-local}
\end{equation}
a local two-body functional with no one-body term.
\item (v) Hence the composite terms of the pairs and of the legs of Theorem \ref{thm:c10-D1}, summed, equal the sum over the pairs of the composite pair kernels along $\Delta$: the nonlocal parts of the brackets of the letters $A|_V$, $\lambda$, $c$ with the renormalised composite two-point functional, $\log|p|^2$ of the leg momentum and $(p,a)(p,b)|p|^{-2}$ on $A|_V$ (Remark \ref{rem:c9-Pi}), cancel against those of the images under $Q$ of the renormalised composite cycles through the pairs containing these letters.
\item (vi) For $sO\ne0$ the right side of \eqref{eq:c11-local} acquires the renormalised composite graphs of $sO$, minus $\tfrac\hbar2\,\partial$ of the composite kernel applied to the trees with one vertex $sO$; these contain the nonlocal bracket with the renormalised composite two-point functional on the letters of $sO$.
\end{itemize}
\end{theorem}

\begin{proof}
(i) Theorem \ref{thm:c11-simplex} applied to each chain of Definition \ref{def:c11-sp}. (ii) Lemma \ref{lem:c11-Cid}.

(iii), (iv) The composite term is the limit of $Q$ of the composite graphs minus their counterterm trees, minus the tree insertion of the composite part. Write the counterterm trees as $-\mathrm{tr}(\mathrm{ct}^{(1)}_x(O)[\varepsilon])$, with $\mathrm{ct}^{(1)}_x(O)[\varepsilon]=\sum_{m\ge1}\varepsilon^{-m}\mathrm{ct}^{(1)}_{x,m}(O)+\log\varepsilon\,\mathrm{lc}^{(1)}_x(O)$ the arity-one counterterm of $O$ along $x$, and $-\mathrm{tr}(I^{CT,(1)},O)$, with $I^{CT,(1)}$ short for $I^{CT}_{\mathfrak{psl}(4|4)}[\varepsilon]^{(1)}(x)$, an odd vertex. By Proposition \ref{prop:c4-defectflow} and Lemma \ref{lem:c6-bvexact}, with $sO=0$,
\[
\begin{aligned}
Q\big(-\mathrm{tr}(\mathrm{ct}^{(1)}_x(O)[\varepsilon])\big)&=-\mathrm{tr}\big(s\,\mathrm{ct}^{(1)}_x(O)[\varepsilon]\big)+(\text{special edges}),\\
Q\big(-\mathrm{tr}(I^{CT,(1)},O)\big)&=-\mathrm{tr}\big(s\,I^{CT,(1)},O\big)-\epsilon_\Xi\{I^{CT,(1)},O\}+(\text{special edges}),
\end{aligned}
\]
the sign of the bracket being that for $\Xi_x$, both vertices having degree $-1$ and being joined to $O$ by a tree. By Proposition \ref{prop:c11-typeI} the sum of the images is
\[
(\mathrm A)+(\mathrm B)+(\mathrm D)+(\text{special edges})+\mathrm{tr}\big(s(\Xi_x-I^{CT,(1)}),O\big)-\epsilon_\Xi\{I^{CT,(1)},O\}-\mathrm{tr}\big(s\,\mathrm{ct}^{(1)}_x(O)[\varepsilon]\big),
\]
the brackets $\pm\epsilon_\Xi\{\Xi_x,O\}$ of the two types cancelling. By Lemma \ref{lem:c11-trees}, (B), (D) and the special edges tend to $0$. By Lemma \ref{lem:c11-Cid}, $\mathrm{tr}(s(\Xi_x-I^{CT,(1)}),O)$ tends to $\mathrm{tr}(s\,\Xi^{\rm ren}_x,O)$, the tree insertion of the composite part, which is subtracted in the definition. What remains is (A) minus $\epsilon_\Xi\log\varepsilon\{L^{(1)}(x),O\}$ minus $\mathrm{tr}(s\,\mathrm{ct}^{(1)}_x(O)[\varepsilon])$. The left side has a limit, the renormalised composite graphs converging and $Q$ multiplying plane-wave values by polynomials in the momenta, and (A) has no $\log\varepsilon$ by (i): the coefficient of $\log\varepsilon$ gives $\mathrm{tr}\big(s\,\mathrm{lc}^{(1)}_x(O)+\epsilon_\Xi\{L^{(1)}(x),O\}\big)=0$, and the coefficient of $\varepsilon^{-m}$, with $a_m(X_a,X_b)$ the coefficient of $\varepsilon^{-m}$ of the special-edge composite cycles of the pair, gives $\mathrm{tr}\big(\sum_{a<b}\pm a_m(X_a,X_b)-s\,\mathrm{ct}^{(1)}_{x,m}(O)\big)=0$; the trees with these local vertices are the same at every $\varepsilon$, and the injectivity of the tree-level flow gives (iii). At order $\varepsilon^0$ the trees on the outputs of (A) converge to trees with edges $P(0,T)$, and the limit is the tree-level flow of the sum over the pairs of the composite pair kernels; Theorem \ref{thm:c8-L} (i) gives \eqref{eq:c11-local}.

(v) By Proposition \ref{prop:c11-typeI}, the image of the first type contains $-\epsilon_\Xi\{\Xi_x,O\}$ and that of the second type $+\epsilon_\Xi\{\Xi_x,O\}$; the counterterm trees replace $\Xi_x$ by $\Xi^{\rm ren}_x$ in the second and add $-\epsilon_\Xi\log\varepsilon\{L^{(1)}(x),O\}$, which (iii) removes from the first. The renormalised composite loops on a letter $X$ give the bracket of $\Xi^{\rm ren}_x$ with $X$, whose quadratic part along $\Delta$ is the functional of Remark \ref{rem:c9-Pi}, local on $A|_{W_{\mathbb R}^\perp}$ and containing $\log|p|^2$ on $A|_V$, $\lambda$, $c$. The form of Theorem \ref{thm:c10-D1} splits the same sum into the images of the renormalised composite cycles through the pairs, the composite terms of the pairs, and the brackets with $\Xi^{\rm ren}_x$ on the letters together with the tree insertion of $s\,\Xi^{\rm ren}_x$, the composite terms of the legs; the brackets cancel between the two, and the sum is local.

(vi) For $sO\neq0$, Proposition \ref{prop:c11-master} keeps the trees with one vertex $sO$, whose closures by the composite kernel are the composite graphs of $sO$ of both types (Proposition \ref{prop:c11-typeI}); with their counterterm trees they converge to the renormalised composite graphs of $sO$, nonlocal through the brackets with $\Xi^{\rm ren}_x$ on the letters of $sO$.
\end{proof}

\begin{remark}[What the locality uses and what it fixes]
\label{rem:c11-local}
The argument uses $sO=0$ exactly, the classical master equation, the anticommutation of $Q$ with $\partial$ of the composite kernel, which rests on $[v_x.,Q]=0$ and $QP\star+P\star Q=e^{-\varepsilon H}$ (Lemma \ref{lem:c11-kernels} (v)), the absence of $\log\varepsilon$ in cycles with one special edge, a property of the cutoff on every edge (Theorem \ref{thm:c11-simplex}), the absence of a term $\varepsilon^{-1}$ in the composite loops of $I$ (Theorem \ref{thm:c11-local} (ii)), and generic momenta at $T=\infty$ for the convergence of the cycles and of the composite loops, which are of order $|p|^{-3}$ (Proposition \ref{prop:c9-composite} (iv)), or any finite $T$ (Theorem \ref{thm:c8-L} (ii)); (i) of Definition \ref{def:c10-hyp} enters through the terms with $QW$ of the proof of Theorem \ref{thm:c10-D1}, which are outside the composite term. For a functional of the fields $A|_V$, $A|_{W_{\mathbb R}^\perp}$, $\lambda$ alone, $sO=0$ is the gauge invariance of its density. Statement (iv) fixes the one-body part of the composite term: a term $t$ units times $v.X_a$ on every letter, added to the right side of \eqref{eq:c11-local}, would replace each sector coefficient $\theta$ by $\theta+t$ in Sections 11.3--11.5 and leave every spectrum unchanged; the operator-level value $\theta=-\tfrac13$ of Proposition \ref{prop:c11-theta-sector} is the value in the normalisation of \eqref{eq:c11-local}, and it agrees with the action-level value of Theorem \ref{thm:c9-theta2}. The steps of the proof are Lemma \ref{lem:c11-Cid} for the composite loops of the action, Lemma \ref{lem:c11-trees} for the special edges on trees and Proposition \ref{prop:c11-typeI} for the cancellation of the brackets $\pm\epsilon_\Xi\{\Xi_x,O\}$ between the two types, forced by Proposition \ref{prop:c11-master}, whose right side contains no bracket; the sign $\epsilon_\Xi$ never needs to be evaluated. The composite term enters the second evaluation of $D_1$ and $P_1$; the components themselves are given by the master formula of Theorem \ref{thm:c11-classes}, which holds along every generator. Statement (iv) is the input of Corollary \ref{cor:c11-D1local}, Propositions \ref{prop:c11-routes} and \ref{prop:c11-scalar} and Sections 11.3--11.5; its consequences are evaluated on the bilinears, where the single chain of Example \ref{ex:c11-bilinear} reproduces the value of the route of the special edges, on the planar scalar sector, and through the spectra of Table \ref{tab:c11-results}, which a common one-body term leaves unchanged.
\end{remark}

\begin{corollary}[Local two-body form of $D_1$ and $P_1$]
\label{cor:c11-D1local}
Under the hypotheses of Definition \ref{def:c10-hyp}, and of the same definition with $P_b$ in place of $\Delta$, for $O=\int\omega\,X_1\cdots X_n$, with the Koszul signs of \eqref{eq:c10-decomp},
\[
\begin{aligned}
D_1(O)=\ &\sum_{a<b}\pm\int\omega\cdots\Big[4\big(\text{coefficient of }\log\varepsilon\text{ in }\partial_\delta I^{CT}[\varepsilon](X_aX_b)\big)+\big(\text{composite pair kernel of }(X_a,X_b)\text{ along }\Delta\big)\Big]\cdots\\
&+\sum_a\int\omega\cdots\Big[-4\,\text{units}\cdot(\text{weight of }X_a)\,X_a-\big\{\hbar F_1^{(1)}(\Delta)-4\,\text{units}\cdot(\text{field-counting functional}),X_a\big\}\Big]\cdots,\\
P_1(O)=\ &\sum_{a<b}\pm\int\omega\cdots\big(\text{composite pair kernel of }(X_a,X_b)\text{ along }P_b\big)\cdots-\sum_a\int\omega\cdots\big\{\hbar F_1^{(1)}(P_b),X_a\big\}\cdots,
\end{aligned}
\]
every term local; $D_1$ minus four times the map of Theorem \ref{thm:c10-brst} is the sum over the pairs of the composite pair kernels along $\Delta$ minus the bracket of $O$ with $\hbar F_1^{(1)}(\Delta)-4$ units times the field-counting functional.
\end{corollary}

\begin{proof}
Insert \eqref{eq:c11-local} into Theorem \ref{thm:c10-D1} and Proposition \ref{prop:c10-P1}; four times the coefficient of $\log\varepsilon$ is two-body by Theorem \ref{thm:c10-decomp} (iii), and the brackets with the field-counting functional, with $\hbar F_1^{(1)}(\Delta)$ and with $\hbar F_1^{(1)}(P_b)$ are derivations.
\end{proof}

\begin{proposition}[The route of the special edges]
\label{prop:c11-routes}
For $O$ as in Corollary \ref{cor:c11-D1local}, the limit as $\varepsilon\to0$ of the cycles through the pairs of letters of $O$ with one special edge and one vertex $S^{(1)}(\Delta)$ anywhere on the cycle, as the composite kernel or with one leg on the special edge, minus the trees of $s$ of the arity-one counterterm of $O$, the graphs of Lemma \ref{lem:c8-transfer} used in Theorem \ref{thm:c9-D1}, is the sum over the pairs of four times the coefficient of $\log\varepsilon$ of the pair kernel plus the composite pair kernel along $\Delta$.
\end{proposition}

\begin{proof}
Both routes compute the local representative $D_1(O)$, which is unique (Theorem \ref{thm:c8-L} (i)). In the route of the special edges the one-body terms are the trees with $-\{\hbar F_1^{(1)}(\Delta),O\}$, the arity-one defect of the action joined to $O$ vanishing by \eqref{eq:c9-arity1} (proof of Theorem \ref{thm:c9-D1}); in Corollary \ref{cor:c11-D1local} the one-body terms are $-4$ units times the weights minus the bracket with $\hbar F_1^{(1)}(\Delta)-4$ units times the field-counting functional, which is $-\{\hbar F_1^{(1)}(\Delta),-\}$. The two-body sums therefore agree.
\end{proof}

\begin{example}[The scalar bilinears]
\label{ex:c11-bilinear}
For $O=\int\omega\,(\iota_{w_1}A,\iota_{w_2}A)$, on the part with two legs $A|_{W_{\mathbb R}^\perp}$, Theorem \ref{thm:c9-D1} and Remark \ref{rem:c9-twoeval} give the value $4\,(\text{coefficient of }\log\varepsilon\text{ in }\partial_\delta I^{CT}[\varepsilon](O))-\tfrac43\,O-\tfrac13\{S^{(1)}(\Delta),O\}$ units of the route of the special edges, so by Proposition \ref{prop:c11-routes}
\[
\big(\text{composite pair kernel of }(A|_{W_{\mathbb R}^\perp},A|_{W_{\mathbb R}^\perp})\text{ along }\Delta\big)\Big|_{\text{outputs }A|_{W_{\mathbb R}^\perp}A|_{W_{\mathbb R}^\perp}}=\Big[-\tfrac13\,v.O-\tfrac43\,O\Big]\ \text{units},
\]
for the traceless and the trace parts alike. It is the single chain of Table \ref{tab:c11-fieldpairs},
\[
A|_{W_{\mathbb R}^\perp}\xrightarrow{\ K_\varepsilon-K_0\ }\big([A|_{W_{\mathbb R}^\perp},c],A^*|_{W_{\mathbb R}^\perp}\big)\xrightarrow{\ \text{composite kernel}\ }A|_V(A|_{W_{\mathbb R}^\perp})^2\xrightarrow{\ P\ }A|_{W_{\mathbb R}^\perp}.
\]
Its position part along $b$ is $-\tfrac23$ units times $\partial_b$, so the composite pair kernel along $P_b$ is $-\tfrac13$ units times $\{S^{(1)}(P_b),O\}$, in agreement with Theorem \ref{thm:c9-P1}, the composite triangle and the composite two-point functional along $P_b$ giving $-2+\tfrac53$; its translation-invariant part is $-\tfrac23\cdot2\,O-\tfrac43\,O=-\tfrac83\,O$ units.
\end{example}

\section{The free directions and the arity-two equation}

\begin{definition}[Finite counterterms along $\Delta$ and $P_b$]
\label{def:c11-free}
Write the solutions of the arity-one equations along $\Delta$ and along $P_b$ as
\[
\begin{aligned}
\hbar F_1^{(1)}(\Delta)=\ &4\,\text{units}\cdot(\text{field-counting functional})+\Big[-\tfrac23\int(A,A^*)+\tfrac13\int(\lambda^*,\lambda)+\tfrac43\int(c^*,c)\Big]\,\text{units}\\
&+\theta\,\text{units}\int\big(2E(A|_{W_{\mathbb R}^\perp})+2A|_{W_{\mathbb R}^\perp},A^*|_{W_{\mathbb R}^\perp}\big)-\theta_\lambda\,\text{units}\int\big(2E(\lambda)+3\lambda,\lambda^*\big)\\
&+\theta_V\,\text{units}\Big[\int\big(2E(A|_V)+2A|_V,A^*|_V\big)-\int\big(2E(c),c^*\big)\Big]+\text{units}\cdot Q\Big(\kappa_V\int(\iota_xA|_V,c^*)+\kappa_\lambda\int(\rho(x)\lambda^*,\lambda^*)\Big)+(\text{cubic and higher}),\\
\hbar F_1^{(1)}(P_b)=\ &\theta_P\,\text{units}\int\big(\partial_b(A|_{W_{\mathbb R}^\perp}),A^*|_{W_{\mathbb R}^\perp}\big)-\theta_{P,\lambda}\,\text{units}\int(\partial_b\lambda,\lambda^*)+\theta_{P,V}\,\text{units}\Big[\int\big(\partial_b(A|_V),A^*|_V\big)-\int(\partial_bc,c^*)\Big]\\
&+Q\big(\text{a combination of }\textstyle\int(\iota_bA|_V,c^*)\text{ and }\int(\rho(b)\lambda^*,\lambda^*)\big)+(\text{cubic and higher}),
\end{aligned}
\]
with numbers $\theta,\theta_\lambda,\theta_V,\kappa_V,\kappa_\lambda,\theta_P,\theta_{P,\lambda},\theta_{P,V}$; $\theta$ is the coefficient of \eqref{eq:c9-Fperp} and $\theta_P$ that of Theorem \ref{thm:c9-P1}. The three terms with $\theta$, $\theta_\lambda$, $\theta_V$ are the sector parts of $S^{(1)}(\Delta)$ in the sectors of $A|_{W_{\mathbb R}^\perp}$, of $\lambda$ and of $(A|_V,c)$ times these numbers, and those with $\theta_P$, $\theta_{P,\lambda}$, $\theta_{P,V}$ the sector parts of $S^{(1)}(P_b)$ (Definition \ref{def:c11-kernels}). The scheme is translation covariant if the position part along $b$ of $\hbar F_1^{(1)}(\Delta)$ is $2\hbar F_1^{(1)}(P_b)$; then $\theta_P=\theta$, $\theta_{P,\lambda}=\theta_\lambda$, $\theta_{P,V}=\theta_V$. The arity-two equations along $(\Delta,P_b)$ and $(\Delta,\xi)$ fix $\theta=\theta_\lambda=\theta_V=-\tfrac13$ (Theorem \ref{thm:c9-theta2}), and those along $(\xi,P_b)$ give $\theta_P=\theta_{P,\lambda}=\theta_{P,V}$ (Theorem \ref{thm:c11-oddar2}).
\end{definition}

\begin{proposition}[Completeness of the quadratic directions]
\label{prop:c11-tc0}
Let $\mathcal F_{-1}$ and $\mathcal F_{-2}$ be the quadratic local functionals of degree $-1$ and $-2$, of Weyl weight $4$ with the position at weight $-1$, at most linear in the position and invariant under $\mathrm{Spin}(4)\times\mathrm{Spin}(6)$. On the parity-even subspaces of Theorem \ref{thm:c9-quadratic} (iii), invariant under the reflection of one direction of $V$ together with one direction of $W_{\mathbb R}^\perp$,
\[
\frac{\ker\big(Q|_{\mathcal F_{-1}}\big)}{Q\,\mathcal F_{-2}}=\operatorname{span}\big\{\text{the three sector parts of }S^{(1)}(\Delta)\big\},\qquad Q\,\mathcal F_{-2}=\operatorname{span}\Big\{Q\int(\iota_xA|_V,c^*),\ Q\int(\rho(x)\lambda^*,\lambda^*)\Big\}.
\]
Without the reflection the quotient has one further class, $\int(\lambda^*,\rho(e_1e_2e_3e_4)\lambda)$, with $e_1,\dots,e_4$ an oriented orthonormal basis of $V$.
\end{proposition}

\begin{proof}
By Schur's lemma, $\mathcal F_{-1}$ is spanned by $\int(A^*|_V,D\,A|_V)$ with the structures $D$ built from $1$, $x\cdot\partial$, $x_\mu\partial_\nu$, $x_\nu\partial_\mu$ on the vector index and $\epsilon_{\mu\nu\kappa\tau}x^\kappa\partial^\tau$; by $\int(A^*|_{W_{\mathbb R}^\perp},D\,A|_{W_{\mathbb R}^\perp})$ with $D\in\{1,x\cdot\partial\}$; by $\int(\lambda^*,D\lambda)$ with $D\in\{1,\ x\cdot\partial,\ \rho(x)\rho(\partial)\}$ times $\{1,\rho(e_1e_2e_3e_4)\}$; and by $\int(c^*,Dc)$ with $D\in\{1,x\cdot\partial\}$. $\mathcal F_{-2}$ is spanned by $\int x^\mu(A_\mu,c^*)$, $\int(\lambda^*,\rho(x)\lambda^*)$ and $\int(\lambda^*,\rho(e_1e_2e_3e_4)\rho(x)\lambda^*)$. The structures with $\epsilon$ or with one factor $\rho(e_1e_2e_3e_4)$ change sign under the reflection. With $Q$ of Proposition \ref{prop:c1-Q}, the ranks on the plane-wave normal forms, computed exactly, are, in the sectors of $A|_{W_{\mathbb R}^\perp}$, of $\lambda$ and of $(A|_V,c)$: $\dim\mathcal F_{-1}=2,3,6$, $\dim\ker Q=1,2,2$, $\dim Q\mathcal F_{-2}=0,1,1$ for the reflection-invariant subspaces. The kernel is spanned by the sector part of $A|_{W_{\mathbb R}^\perp}$; by the sector part of $\lambda$ and $Q\int(\rho(x)\lambda^*,\lambda^*)$; by the sector part of $(A|_V,c)$ and $Q\int(\iota_xA|_V,c^*)$. Without the reflection, $\dim\mathcal F_{-1}=15$, the $\epsilon$-structure is not $Q$-closed, $\int(\lambda^*,\rho(e_1e_2e_3e_4)\lambda)$ is closed and not exact, and $\int(\lambda^*,\rho(e_1e_2e_3e_4)\rho(x)\lambda^*)=0$.
\end{proof}

\begin{proposition}[What the arity-one equations fix]
\label{prop:c11-arity1}
\begin{itemize}
\item (i) Along $\Delta$ the quadratic part of \eqref{eq:c9-arity1} fixes the translation-invariant part $\big[-\tfrac23\int(A,A^*)+\tfrac13\int(\lambda^*,\lambda)+\tfrac43\int(c^*,c)\big]$ units (Theorem \ref{thm:c9-quadratic} (iii)) and leaves $\theta$, $\theta_\lambda$, $\theta_V$, $\kappa_V$, $\kappa_\lambda$ free. Along $P_b$ the quadratic part of the arity-one obstruction vanishes (Proposition \ref{prop:c11-position} (i)), so the quadratic part of $\hbar F_1^{(1)}(P_b)$ is $Q$-closed and every coefficient of Definition \ref{def:c11-free} along $P_b$ is free in arity one, the arity-two equations relating them to the coefficients along $\Sigma$ (Theorem \ref{thm:c11-oddar2}); the translation-covariant choice, $\hbar F_1^{(1)}(P_b)$ equal to half the position part along $b$ of $\hbar F_1^{(1)}(\Delta)$, solves the equation along $P_b$ at every order in the fields.
\item (ii) Two solutions along $\Delta$ differ by an $s$-closed local functional of degree $-1$ and Weyl weight $4$. The common shift of $\theta$, $\theta_\lambda$, $\theta_V$ by $t$ adds $t$ units times $S^{(1)}(\Delta)$, which is $s$-closed (Lemma \ref{lem:c9-SDelta} (iii)); the directions $Q\int(\iota_xA|_V,c^*)$ and $Q\int(\rho(x)\lambda^*,\lambda^*)$ complete to $s\int(\iota_xA|_V,c^*)$ and $s\int(\rho(x)\lambda^*,\lambda^*)$, and an $s$-exact direction $sG$ contributes $s\{G,O\}$ to $D_1(O)$ on cocycles, zero on cohomology.
\item (iii) Since $\{I,S^{(1)}(\Delta)\}=0$ while the brackets of $I$ with the sector parts are nonzero, the cubic part of \eqref{eq:c9-arity1} depends on the sector coefficients through $(\theta-\theta_V)$ times the bracket of the cubic part of $I$ with the sector part of $A|_{W_{\mathbb R}^\perp}$ plus $(\theta_\lambda-\theta_V)$ times its bracket with the sector part of $\lambda$. It is solvable exactly when this sum is congruent to the cubic part of the composite part minus the bracket of the cubic part of $I$ with the translation-invariant part, modulo the images under $Q$ of the cubic local functionals of degree $-1$. If the classes of the two brackets modulo the images under $Q$ of the cubic local functionals of degree $-1$, Weyl weight $4$ and degree at most one in the position are linearly independent, the cubic equation fixes $\theta-\theta_V$ and $\theta_\lambda-\theta_V$.
\end{itemize}
\end{proposition}

\begin{proof}
(i) Theorem \ref{thm:c9-quadratic} (iii) along $\Delta$, Proposition \ref{prop:c11-position} (i) along $P_b$. For the covariant choice: $s$ maps letters to letters and carries no position, so it commutes with taking the position part, and the position part of the arity-one obstruction along $\Delta$ is twice the obstruction along $P_b$ by Proposition \ref{prop:c11-position} (i), the term $4$ units times $s$ of the field-counting functional having no position part. (ii) Theorem \ref{thm:c8-obstruction} (iii) restricted to arity one along $\Delta$, and Lemma \ref{lem:c10-deriv}: $\{sG,O\}=-s\{G,O\}$ for $sO=0$ and $G$ of degree $-2$. (iii) $s$ of the sector terms is the sum of $\theta$, $\theta_\lambda$, $\theta_V$ units times the brackets of $I$ with the sector parts, whose sum over the three sectors vanishes; at cubic order $s$ of the solution is $Q$ of its cubic part plus the bracket of the cubic part of $I$ with its quadratic part.
\end{proof}

\begin{theorem}[The position part of $D_1$ is exact]
\label{thm:c11-arity2}
Assume a one-loop quantisation through arity two along $\Delta$ and $P_b$, that is, \eqref{eq:c8-qme} modulo $\hbar^2$ in its components of arity at most two along these generators, and local representatives through arity two along $(\Delta,P_b)$ on a space $\mathfrak F$ of local functionals stable under $\Lambda$, with $P_1$ translation covariant, $P_1(\int\omega\,D)=\int\omega\,P_1(D)$ for a map $P_1$ of densities commuting with the derivatives; this holds when $\hbar F_1^{(1)}(P_b)$ is translation invariant. Then for every $O=\int\omega\,D\in\mathfrak F$ with $sO=0$
\[
\big(\text{position part of }D_1\text{ along }b\big)(O)=\big[D_1,\{S^{(1)}(P_b),-\}\big](O)=\pm\,s\Big(\hbar L_1(O)^{(2)}(\Delta,P_b)\Big):
\]
the position part of $D_1$ is $s$-exact on cocycles, with the arity-two component of $O$ as primitive. With the composite pair kernels this gives an operator-level determination of the sector coefficients along $\Delta$ (Proposition \ref{prop:c11-theta-sector}), which reproduces the values of the action-level arity-two equations (Theorem \ref{thm:c9-theta2}).
\end{theorem}

\begin{proof}
By Corollary \ref{cor:c8-homotopy} (ii), $[\Lambda^{(1)}(\Delta),\Lambda^{(1)}(P_b)]+\Lambda^{(1)}([\Delta,P_b])$ is a graded commutator of $\Lambda^{(0)}$ with an operator built from $\Lambda^{(2)}(\Delta,P_b)$. With $[\Delta,P_b]=-2P_b$ (Remark \ref{rem:c2-dilatation}) and $S^{(2)}(\Delta,P_b)=0$ (Definition \ref{def:c2-S2}), the order-$\hbar^0$ part of that operator vanishes and the order-$\hbar$ part of the identity reads
\[
\big[D_1,\{S^{(1)}(P_b),-\}\big]+\big[v.,P_1\big]-2P_1=\big[s,H_1\big],
\]
with $H_1$ built from $\hbar L_1^{(2)}(\Delta,P_b)$. For translation-covariant $P_1$, $v.(P_1O)=\int\omega\,2(x\cdot\partial+w_D+1)P_1(D)$ and $P_1(v.O)=\int\omega\,\big(2x^\mu P_1(\partial_\mu D)+2w_DP_1(D)\big)$, so $[v.,P_1]=2P_1$. Lemma \ref{lem:c11-canonical} identifies $[D_1,\{S^{(1)}(P_b),-\}]$ with the position part, and on a cocycle $[s,H_1]O=s(H_1O)$.
\end{proof}

\begin{lemma}[Position part of a bracket]
\label{lem:c11-bracketpos}
For a local functional $G=\int(x^\mu g_\mu+g_0)$, the position part along $b$ of $\{G,-\}$ is $\{\int b^\mu g_\mu,-\}$.
\end{lemma}

\begin{proof}
The bracket contracts a slot of $g_\mu$ at $y$ with a letter $\partial^\alpha\varphi$ of $D$ at $x$ through $K_0$, which gives $\partial^\alpha_x(x^\mu\cdots)$; the term with no derivative on $x^\mu$ is $x^\mu$ times the bracket with $\int g_\mu$, and the others carry no position.
\end{proof}

\begin{corollary}[The composite term along $P_b$ and the translation component]
\label{cor:c11-theta}
Under the hypotheses of Theorem \ref{thm:c11-arity2} and of Definition \ref{def:c10-hyp}, on every cocycle $O\in\mathfrak F$:
\begin{itemize}
\item (i) twice the composite term of $O$ along $P_b$ minus the bracket of $O$ with the position part along $b$ of $\hbar F_1^{(1)}(\Delta)$ is $s$-exact; this position part is
\[
2\theta\,\text{units}\int\big(\partial_b(A|_{W_{\mathbb R}^\perp}),A^*|_{W_{\mathbb R}^\perp}\big)-2\theta_\lambda\,\text{units}\int(\partial_b\lambda,\lambda^*)+2\theta_V\,\text{units}\Big[\int\big(\partial_b(A|_V),A^*|_V\big)-\int(\partial_bc,c^*)\Big]+\text{units}\cdot Q\Big(\kappa_V\int(\iota_bA|_V,c^*)+\kappa_\lambda\int(\rho(b)\lambda^*,\lambda^*)\Big)
\]
plus the position parts of the higher terms;
\item (ii) if $\theta=\theta_\lambda=\theta_V$, the composite term along $P_b$ induces $\theta$ units times $\{S^{(1)}(P_b),-\}$ on cohomology on the parts with the fewest fields of the cocycles of $\mathfrak F$, on which $\hbar F_1^{(1)}(\Delta)$ acts through its quadratic part, and on every family in which $\{S^{(1)}(P_b),O\}$ is not exact for some cocycle the composite term along $P_b$ determines $\theta$;
\item (iii) in the translation-covariant scheme the position part of $D_1$ is $2P_1$, and the arity-two equation along $(\Delta,P_b)$ is equivalent to $P_1\cong0$ on cohomology: $P_1(O)=\pm\tfrac12\,s\big(\hbar L_1(O)^{(2)}(\Delta,P_b)\big)$;
\item (iv) if $\theta=\theta_\lambda=\theta_V$ and $\theta_P=\theta_{P,\lambda}=\theta_{P,V}$, then $P_1$ induces $(\theta-\theta_P)$ units times $\{S^{(1)}(P_b),-\}$ on cohomology, and the finite part of the operators $\hbar\psi'=\tfrac a2\{S^{(1)}(\Delta),-\}$ (Definition \ref{def:c8-renorm}; the letter $\psi$ of Section 11.5 is a different object), $a=(\theta-\theta_P)$ units, of Proposition \ref{prop:c9-scheme} gives $P_1\cong0$ on every family at once, leaves $D_1$, $M_1$, $R_1$ unchanged and changes the components along $K_b$, $\Sigma$ and $\Sigma^*$ by $a$, $-\tfrac a2$ and $\tfrac a2$ times the classical transformation. If the differences of the coefficients along $P_b$ and along $\Delta$ depend on the sector, a multiple of $\{S^{(1)}(\Delta),-\}$ removes $P_1$ on the families with letters of one sector, and the translation-covariant scheme removes it on every family.
\end{itemize}
\end{corollary}

\begin{proof}
(i) By Corollary \ref{cor:c11-D1local} and Proposition \ref{prop:c11-position} (i), the position part of $D_1$ is twice the sum over the pairs of the composite pair kernels along $P_b$ minus the bracket with the position part of $\hbar F_1^{(1)}(\Delta)$, by Lemma \ref{lem:c11-bracketpos}; the coefficients of $\log\varepsilon$ and the bracket with the field-counting functional commute with the translations. Theorem \ref{thm:c11-arity2} gives exactness, and the position part of each sector part of $S^{(1)}(\Delta)$ is twice the sector part of $S^{(1)}(P_b)$ (Definition \ref{def:c11-kernels}). (ii) For equal coefficients the $\theta$-terms are $2\theta$ units times $S^{(1)}(P_b)$, which is $s$-closed, and the other directions are $s$-exact after completion, with $\{sG,O\}=-s\{G,O\}$. (iii) In the covariant scheme the position part of $\hbar F_1^{(1)}(\Delta)$ is $2\hbar F_1^{(1)}(P_b)$, and $P_1$ is the sum of the composite pair kernels along $P_b$ minus the bracket with $\hbar F_1^{(1)}(P_b)$ (Corollary \ref{cor:c11-D1local}). (iv) By (ii), $P_1\cong\theta$ units times $\{S^{(1)}(P_b),-\}$ minus $\theta_P$ units times the same; Proposition \ref{prop:c9-scheme} with $\lambda_{P_b}=-2$. For sector-dependent differences $P_1$ induces the sum over the sectors of the differences times the brackets with the sector parts of $S^{(1)}(P_b)$, plus the completions, which is a multiple of $\{S^{(1)}(P_b),-\}$ on the cocycles with letters of one sector.
\end{proof}

\section{The dilatation identities in kernel form}

Corollary \ref{cor:c10-P7} requires that on cohomology the composite term minus the bracket with the translation-invariant finite counterterm beyond the field-counting functional induce $-\tfrac13$ units times $\{S^{(1)}(\Delta),-\}$, the operator cancelled by the sector coefficients $-\tfrac13$ of Theorem \ref{thm:c9-theta2}. By Corollary \ref{cor:c11-D1local} this is a statement about the composite pair kernels, which splits into a statement along $P_b$ and a statement about the translation-invariant part.

\begin{lemma}[One-body terms]
\label{lem:c11-onebody}
For $\theta=\theta_\lambda=\theta_V$, $\kappa_V=\kappa_\lambda=0$ and the quadratic part of $\hbar F_1^{(1)}(\Delta)$ of Definition \ref{def:c11-free}, the translation-invariant part of $-\big\{\hbar F_1^{(1)}(\Delta)-4\,\text{units}\cdot(\text{field-counting functional}),X\big\}$ on a field letter $X=\partial^\alpha\varphi$ is $(c_\varphi-2\theta\,w_X)$ units times $X$, with $c_\varphi=\tfrac23$ for $A|_V$ and $A|_{W_{\mathbb R}^\perp}$, $\tfrac13$ for $\lambda$ and $\tfrac43$ for $c$, and its position part along $b$ is $-2\theta$ units times $\partial_bX$. At $\theta=-\tfrac13$ the translation-invariant part is $\tfrac23(2+|\alpha|)$ units times $X$ for every field slot and for $c$.
\end{lemma}

\begin{proof}
$\{\int(A,A^*),X\}=X$, $\{\int(\lambda^*,\lambda),\lambda\}=-\lambda$ and $\{\int(c^*,c),c\}=-c$ (Lemma \ref{lem:c4-count}), and the bracket of the sector part with a letter of its sector is $v.X$, with translation-invariant part $2w_XX$ and position part $2\partial_bX$ (Lemma \ref{lem:c11-canonical}). At $\theta=-\tfrac13$: $\tfrac23+\tfrac23(1+|\alpha|)$, $\tfrac13+\tfrac23(\tfrac32+|\alpha|)$ and $\tfrac43+\tfrac23|\alpha|$ are all $\tfrac23(2+|\alpha|)$.
\end{proof}

\begin{definition}[The identities in kernel form]
\label{def:c11-P7}
Let $\mathfrak F$ be a family of gauge-invariant cocycles of field letters, preserved by $D_1$ and contained in a space of local functionals stable under $\Lambda$, satisfying (ii) and (iii) of Definition \ref{def:c10-hyp}. The letter count is the derivation that multiplies a monomial $\int\omega\,X_1\cdots X_n$ of letters $X_a=\partial^{\alpha_a}\varphi_a$, $\varphi_a\in\{A|_V,A|_{W_{\mathbb R}^\perp},\lambda\}$, by $\sum_a(2+|\alpha_a|)$; it is defined on canonical forms, and its commutator with $\{S^{(1)}(P_b),-\}$ is $\{S^{(1)}(P_b),-\}$. The three identities on $\mathfrak F$ are:
\begin{itemize}
\item (a) the cubic part of \eqref{eq:c9-arity1} forces $\theta=\theta_\lambda=\theta_V$;
\item (b) the sum over the pairs, with the Koszul signs of \eqref{eq:c10-decomp}, of the composite pair kernels along $P_b$ induces $-\tfrac13$ units times $\{S^{(1)}(P_b),-\}$ on the cohomology of $\mathfrak F$;
\item (c) the sum over the pairs of the translation-invariant parts of the composite pair kernels along $\Delta$ induces $-\tfrac23$ units times the letter count on the cohomology of $\mathfrak F$.
\end{itemize}
\end{definition}

\begin{remark}[Dependences among the identities]
\label{rem:c11-P7dep}
(1) By Proposition \ref{prop:c11-position} (ii), pair by pair the commutator of the translation-invariant part of the composite pair kernel along $\Delta$ with $\{S^{(1)}(P_b),-\}$ is twice the composite pair kernel along $P_b$, and the commutator of the letter count with $\{S^{(1)}(P_b),-\}$ is $\{S^{(1)}(P_b),-\}$; hence (c) on $\mathfrak F$ and on $\{S^{(1)}(P_b),\mathfrak F\}$ gives (b) along $b$, and the content of (b) beyond (c) lies in the directions $b$ with $\{S^{(1)}(P_b),\mathfrak F\}$ outside $\mathfrak F$. (2) Identity (a) follows from (b) on three families with letters of one sector each, under the hypotheses of Theorem \ref{thm:c11-arity2} (Proposition \ref{prop:c11-theta-sector}); the cubic arity-one equation is then a consistency condition. (3) Identities (b) and (c) are identities between graph sums at the regulator of Definition \ref{def:c8-PT}; the numbers $-\tfrac13$ and $-\tfrac23$ are finite parts of special-edge cycles, which depend on the cutoff shape (Remark \ref{rem:c10-scheme1} (d), Proposition \ref{prop:c10-shape}). The statement independent of the scheme is the spectrum of $D_1$ at the coefficients fixed by the arity-two equation (Remark \ref{rem:c10-scheme1}).
\end{remark}

\begin{theorem}[Consequences of the identities]
\label{thm:c11-P7}
Assume (a), (b) and (c) of Definition \ref{def:c11-P7} on $\mathfrak F$.
\begin{itemize}
\item (i) On the parts with the fewest fields of the cocycles of $\mathfrak F$, on which $\hbar F_1^{(1)}(\Delta)$ acts through its quadratic part, $D_1$ minus four times the map of Theorem \ref{thm:c10-brst} induces $-(\tfrac13+\theta)$ units times $\{S^{(1)}(\Delta),-\}$ on cohomology: the assumption of Corollary \ref{cor:c10-P7}, for the solution with sector coefficient $\theta$.
\item (ii) Under the hypotheses of Theorem \ref{thm:c11-arity2}, and if $\{S^{(1)}(P_b),O\}$ is not exact for some cocycle $O\in\mathfrak F$, $\theta=-\tfrac13$, and $D_1$ induces four times the map of Theorem \ref{thm:c10-brst}; its spectrum on $\mathfrak F$ in a fixed classical Weyl weight is twice that of the one-loop dilatation coefficients, $-8$ units times the eigenvalues of $\sum_lH_{l,l+1}$ of [B03c, (1.1)--(1.3)] on the planar families of Corollary \ref{cor:c10-beisert} and of Table \ref{tab:c11-results}.
\item (iii) Under the hypotheses of (ii), in the translation-covariant scheme $P_1\cong0$ on $\mathfrak F$; otherwise Corollary \ref{cor:c11-theta} (iv) applies.
\end{itemize}
Conversely, if $D_1$ minus four times the map induces $-(\tfrac13+\theta)$ units times $\{S^{(1)}(\Delta),-\}$ for equal sector coefficients, then (b) and (c) hold on $\mathfrak F$.
\end{theorem}

\begin{proof}
(i) By Corollary \ref{cor:c11-D1local}, $D_1$ minus four times the map is the sum of the composite pair kernels along $\Delta$ minus the bracket with $\hbar F_1^{(1)}(\Delta)-4$ units times the field-counting functional. Its position part along $b$ is, by (b) and Lemma \ref{lem:c11-onebody}, $2\cdot(-\tfrac13)\partial_b-2\theta\,\partial_b=-2(\tfrac13+\theta)\partial_b$ units. Its translation-invariant part is, by (c) and Lemma \ref{lem:c11-onebody}, $-\tfrac23$ units times the letter count plus $\sum_a(c_{\varphi_a}-2\theta w_{X_a})$ units, which is $-2(\tfrac13+\theta)w_O$ units, by $\sum_a(c_{\varphi_a}+\tfrac23w_{X_a})=\tfrac23\sum_a(2+|\alpha_a|)$. Together they give $-(\tfrac13+\theta)(2x\cdot\partial+2w_O)$ units, which is $-(\tfrac13+\theta)$ units times $v.$ (Lemma \ref{lem:c11-canonical}). (ii) By Theorem \ref{thm:c11-arity2} the position part of $D_1$, which is $-2(\tfrac13+\theta)$ units times $\partial_b$ on cohomology, is exact, and $\partial_bO$ is not exact for that cocycle, so $\theta=-\tfrac13$; then Theorem \ref{thm:c10-brst} and Corollary \ref{cor:c10-beisert}. (iii) Corollary \ref{cor:c11-theta} (iii). The converse reverses the two computations of (i).
\end{proof}

\begin{proposition}[The arity-two equation fixes each sector coefficient]
\label{prop:c11-theta-sector}
Assume the hypotheses of Theorem \ref{thm:c11-arity2} and of Definition \ref{def:c10-hyp}. Let $\mathfrak F_\perp$, $\mathfrak F_\lambda$, $\mathfrak F_V$ be families of gauge-invariant cocycles whose two-field parts have letters of the sector of $A|_{W_{\mathbb R}^\perp}$, of $\lambda$ and of $A|_V$ alone, each containing a cocycle $O$ for which $\{S^{(1)}(P_b),O\}$ is not exact. If (b) holds along $b$ on $\mathfrak F_\perp$, then $\theta=-\tfrac13$; if it holds along $b$ on $\mathfrak F_\lambda$, then $\theta_\lambda=-\tfrac13$; if it holds along $b$ on $\mathfrak F_V$, then $\theta_V=-\tfrac13$. In particular (b) on the three families gives $\theta=\theta_\lambda=\theta_V=-\tfrac13$, the equality asserted in (a), and the cubic part of \eqref{eq:c9-arity1} admits the solution $\theta-\theta_V=\theta_\lambda-\theta_V=0$. These are the values computed from the action-level arity-two equations along $(\Delta,P_b)$ and $(\Delta,\xi)$ (Theorem \ref{thm:c9-theta2}); on each family the operator-level determination is a cross-check of that computation.
\end{proposition}

\begin{proof}
By Corollary \ref{cor:c11-D1local}, Proposition \ref{prop:c11-position} (i) and Lemma \ref{lem:c11-bracketpos}, the position part of $D_1$ is twice the sum of the composite pair kernels along $P_b$ minus the bracket with the position part of $\hbar F_1^{(1)}(\Delta)$. On the two-field part $\hbar F_1^{(1)}(\Delta)$ acts through its quadratic part; the bracket of the sector part of the sector of $X$ has the position part $2\partial_bX$ and the other sector parts have bracket $0$ with $X$, and the $\kappa$-directions are $s$-exact (Proposition \ref{prop:c11-arity1} (ii)). With (b), on the cohomology of the family of the sector with coefficient $\theta_\star\in\{\theta,\theta_\lambda,\theta_V\}$,
\[
\big(\text{position part of }D_1\text{ along }b\big)\cong-\tfrac23\,\text{units}\cdot\partial_b-2\theta_\star\,\text{units}\cdot\partial_b=-2\big(\tfrac13+\theta_\star\big)\,\text{units}\cdot\partial_b .
\]
The left side is $s$-exact on cocycles (Theorem \ref{thm:c11-arity2}) and $\partial_bO$ is not exact, so $\theta_\star=-\tfrac13$. The last statement is Proposition \ref{prop:c11-arity1} (iii); the existence of the quantisation is a hypothesis of Theorem \ref{thm:c11-arity2}.
\end{proof}

\begin{corollary}[The dilatation component without the cubic equation]
\label{cor:c11-noA}
Under the hypotheses of Proposition \ref{prop:c11-theta-sector}, if (b) holds on $\mathfrak F_\perp$, $\mathfrak F_\lambda$ and $\mathfrak F_V$, and (b) and (c) hold on $\mathfrak F$, then on the cohomology of $\mathfrak F$, $D_1$ induces four times the map of Theorem \ref{thm:c10-brst}, with spectrum $-8$ units times the eigenvalues of $\sum_lH_{l,l+1}$ on the planar families of Corollary \ref{cor:c10-beisert} and of Table \ref{tab:c11-results}, and $P_1\cong0$ in the translation-covariant scheme.
\end{corollary}

\begin{proof}
Proposition \ref{prop:c11-theta-sector} gives $\theta=\theta_\lambda=\theta_V=-\tfrac13$, and Theorem \ref{thm:c11-P7} (i)--(iii) applies.
\end{proof}

\begin{proposition}[The planar scalar sector]
\label{prop:c11-scalar}
On the part with the outputs $A|_{W_{\mathbb R}^\perp}$ of the planar scalar sector (Proposition \ref{prop:c10-scalarsector}), (b) and (c) hold as identities modulo terms whose sum over the cyclically adjacent pairs of every trace vanishes: per adjacent pair the composite pair kernel along $\Delta$ is $\big[-\tfrac16v.(X_aX_b)-\tfrac23X_aX_b\big]$ units and along $P_b$ it is $-\tfrac16$ units times $\partial_b(X_aX_b)$, so that on $O=\operatorname{tr}(\iota_{w_1}A\cdots\iota_{w_L}A)$ the composite term along $\Delta$ is $-\big(\tfrac13v.+\tfrac23L\big)O$ units, the composite term along $P_b$ is $-\tfrac13$ units times $\partial_bO$, and the sum of the translation-invariant parts is $-\tfrac43L$ units, $-\tfrac23$ units times the letter count. On the bilinears these are exact identities (Example \ref{ex:c11-bilinear}).
\end{proposition}

\begin{proof}
By Proposition \ref{prop:c10-scalarsector} (i), (e), the composite term of an adjacent pair is $-v.(X_aX_b)$ units, and the bracket with the renormalised composite two-point functional on a letter is $(\tfrac53v.X-\tfrac23X)$ units (Remark \ref{rem:c9-Pi}); by Theorem \ref{thm:c11-local} (v) their sum is the sum of the composite pair kernels, and the one-body term distributed over the two adjacent pairs of each letter gives the density per pair. For $L$ letters without derivatives the letter count is $2L$ and the translation-invariant part of $v.$ is $2L$: $-\tfrac13\cdot2L-\tfrac23L=-\tfrac43L$. Along $P_b$: Theorem \ref{thm:c9-P1}, $-2+\tfrac53=-\tfrac13$.
\end{proof}

\section{The composite kernels on the families}

Throughout this section $\mathfrak g=\mathfrak{su}(N)$ at leading order in $N$ and $O$ is a single trace, with the conventions of Section 10.7. The projection to cohomology on the two-field part is the evaluation on the free equations of motion, $|p|^2=0$ on the scalar legs, the Weyl equation on the gaugino legs and transverse on-shell gauge legs, which removes the $s$-exact terms $s(\text{antifield}\times\text{field})$; total derivatives are kept, as descendants through $\omega$. The letters are written in the components of Definition \ref{def:c8-twocomp}.

\begin{definition}[Test families]
\label{def:c11-families}
Let $n_m:=\sigma^m_{1\dot1}=(0,0,1,i)$, a null vector, and $D_n:=n_m(\partial_m+[A_m,-])$, the covariant derivative $\iota_nd_{A|_V}$ along $n$; let $Z=\iota_wA$ with $w\in W_{\mathbb R}^\perp\otimes\mathbb C$, $(w,w)=0$; let $\psi:=\lambda_1{}^i$ for a fixed $i$, the component $\alpha=1$ of $\lambda|_{S_+\otimes W}$; and let $F_{11}$ be the component $\alpha=\beta=1$ of $(\sigma^m\bar\sigma^n)_\alpha{}^\gamma\varepsilon_{\gamma\beta}F_{mn}$, $F_{mn}=\partial_mA_n-\partial_nA_m+[A_m,A_n]$, whose two-field part is a component of $dA|_V$ and whose three-field part is a component of $[A|_V\wedge A|_V]$. The test families are the spans of the following single traces and their gauge completions:
\begin{itemize}
\item the scalars: $\operatorname{tr}(\iota_{w_1}A\cdots\iota_{w_L}A)$, the bilinears, and the $\mathfrak{su}(2)$ sector of Proposition \ref{prop:c10-su2} for $L\le8$;
\item twist two: $\operatorname{tr}\big(D_n^kZ\,D_n^{S'-k}Z\big)$, $0\le k\le S'$ (Proposition \ref{prop:c10-twist});
\item the $\mathfrak{su}(1|1)$ sector: single traces of the letters $Z$ and $\psi$ of Proposition \ref{prop:c10-su11}, $L\le6$;
\item the fermion $\mathfrak{sl}(2)$ sector: $\operatorname{tr}\big(D_n^k\psi\,D_n^{S'-k}\psi\big)$;
\item two gaugino components: $\operatorname{tr}\big(D_n^k\psi^{(3)}\,D_n^{S'-k}\psi^{(4)}\big)$ with $\psi^{(3)}=\lambda_1{}^3$, $\psi^{(4)}=\lambda_1{}^4$, together with $\operatorname{tr}\big(D_n^kZ'\,D_n^{S'-k}F_{11}\big)$, $Z'$ the scalar letter with the R-symmetry weights of $\psi^{(3)}\psi^{(4)}$, with which the family closes;
\item the gluon $\mathfrak{sl}(2)$ sector: $\operatorname{tr}\big(D_n^kF_{11}\,D_n^{S'-k}F_{11}\big)$;
\item the mixed gluon--gaugino family: $\operatorname{tr}\big(D_n^kF_{11}\,D_n^{S'-k}\psi\big)$.
\end{itemize}
Twist two, the fermion and the gluon $\mathfrak{sl}(2)$ sectors are families of Proposition \ref{prop:c11-theta-sector} for the sectors of $A|_{W_{\mathbb R}^\perp}$, of $\lambda$ and of $A|_V$.
\end{definition}

The two-field part of a gauge-invariant operator with covariant derivatives receives the pairs of its two leading letters with two outputs and the pairs of a letter $A|_V$ from a covariant derivative, or from the three-field part of a field strength, with another letter, with one output (Proposition \ref{prop:c10-twist}). For the composite pair kernels these are the leg sets $A|_{W_{\mathbb R}^\perp}$ of $(A|_V,A|_{W_{\mathbb R}^\perp})$, $\lambda$ of $(A|_V,\lambda)$ and $A|_V$ of $(A|_V,A|_V)$ of Table \ref{tab:c11-fieldpairs}, all with $k=1$, with derivatives.

\begin{proposition}[Spectra from the harmonic action]
\label{prop:c11-targets}
On the two-site graded cyclic states, single traces of two letters with $S'$ derivatives $D_n$ in total, the eigenvalues of $\sum_lH_{l,l+1}$ of [B03c, (1.1)--(1.3)], computed from the harmonic action [B03c, (4.27)--(4.29)] on the oscillator letters [B03c, (4.9)], are:
\begin{itemize}
\item (i) scalars $D_n^kZ$: $4h(S)$, $S\le S'$ even [B03c, (2.31)];
\item (ii) gluons $D_n^kF_{11}$: $4h(S+2)$, $S\le S'$ even, so that $\operatorname{tr}(F_{11}F_{11})$ has the Konishi value $6$;
\item (iii) one gaugino component $D_n^k\psi$: $4h(S+1)$, $S\le S'$ odd;
\item (iv) two gaugino components $\psi^{(3)}$, $\psi^{(4)}$ with $Z'F_{11}$: $0,6$ at $S'=0$; $0,6,6,6$ at $S'=1$; $0,6,6,6,6,\tfrac{25}3$ at $S'=2$; $0$, four times $6$ and three times $\tfrac{25}3$ at $S'=3$; $0$, four times $6$, four times $\tfrac{25}3$ and $\tfrac{49}5$ at $S'=4$;
\item (v) $(F_{11},\psi)$: $6$; $6,6$; $6,6,\tfrac{25}3$; $6,6,\tfrac{25}3,\tfrac{25}3$; $6,6,\tfrac{25}3,\tfrac{25}3,\tfrac{49}5$ at $S'=0,\dots,4$.
\end{itemize}
On the $\mathfrak{su}(1|1)$ sector the two-site Hamiltonian is one minus the graded exchange (Proposition \ref{prop:c10-su11}), with the cyclic spectra of Table \ref{tab:c11-su11}.
\end{proposition}

\begin{proof}
Exact evaluation of [B03c, (4.27)--(4.29)], with the harmonic numbers of the coefficients $c_{n,0,0}=h(\tfrac n2)$, on bases of the families closed under $H_{12}$, with the graded exchange for $H_{21}$ and the projection to graded cyclic states, and exact diagonalisation. The same evaluation gives [B03c, (2.31)] on twist two and the values $0$, twenty times, and $6$ on the scalar bilinears, which fixes the central charge per site, the signs of the fermionic oscillators and the graded cyclicity.
\end{proof}

\begin{tabular}{l|l}
$(L,M)$ & eigenvalues on the graded cyclic states of $\operatorname{tr}(Z^{L-M}\psi^M)$\\
\hline
$M\le1$ & $0$\\
$(3,2)$, $(3,3)$ & $6$\\
$(4,2)$, $(4,3)$ & $4$\\
$(5,2)$, $(5,3)$ & $5\pm\sqrt5$\\
$(5,4)$, $(5,5)$ & $10$\\
$(6,2)$ & $2$, $6$\\
$(6,3)$ & $2$, $6$, $8$, $8$\\
$(6,4)$ & $8$, $8$, $8$\\
$(6,5)$ & $8$\\
\end{tabular}
\tabcaption{The eigenvalues of $\sum_lH_{l,l+1}$ on the $\mathfrak{su}(1|1)$ sector, $L\le6$; the sectors $(2,2)$, $(4,4)$, $(6,6)$ have no graded cyclic state.\label{tab:c11-su11}}

\begin{theorem}[The composite pair kernels on the families]
\label{thm:c11-families}
\begin{itemize}
\item (i) (The two-site density.) For every ordered adjacent pair $(X_a,X_b)$ of letters among $A|_{W_{\mathbb R}^\perp}$, $\lambda|_{S_+\otimes W}$, $\lambda|_{S_-\otimes W^*}$ without derivatives, on the free shell, and for two letters $A|_{W_{\mathbb R}^\perp}$ also off the shell and for all their components, the planar composite pair kernels have, on the same word $X_aX_b$, the translation-invariant part $-\tfrac43$ units times $X_aX_b$ along $\Delta$ and the value $-\tfrac16$ units times $\partial_b(X_aX_b)$ along $P_b$; on the exchanged word and on the trace they vanish.
\item (ii) On the parts with the fewest fields of the families of Definition \ref{def:c11-families}, the two-field parts for two letters, on the free shell, the sum over the pairs of the translation-invariant parts of the composite pair kernels along $\Delta$ is the diagonal operator of Table \ref{tab:c11-results}, which is $-\tfrac23$ units times the letter count, and the sum of the composite pair kernels along $n$ is $-\tfrac13$ units times $\partial_n$; on twist two for $S'\le4$ the sum along every $b\in V$ is $-\tfrac13$ units times $\partial_b$.
\item (iii) Hence (b) along $n$ and (c) of Definition \ref{def:c11-P7} hold on each family in the ranges of Table \ref{tab:c11-results}, and (b) along every $b$ holds on twist two for $S'\le4$.
\end{itemize}
\end{theorem}

\begin{proof}
Exact evaluation of the special-edge composite cycles by Theorem \ref{thm:c11-simplex}, Propositions \ref{prop:c11-genfun} and \ref{prop:c11-k1}, with the colour words reduced to the planar adjacent form of Lemma \ref{lem:c10-colour} (iv), on the leg sets of Table \ref{tab:c11-fieldpairs} with the derivatives of the families, the outputs evaluated on the free shell (Remark \ref{rem:c11-method}). On the families with derivatives the outputs depend on the output momenta through $(n,p)$, every derivative being contracted with $n$ and $(n,n)=0$; the sums are diagonal in the bases of the families. (iii) is (i), (ii) and the values of the letter count: $2L$ on $L$ letters without derivatives, $4+S'$ on twist two and on two gaugino letters, $5+S'$ on $Z'F_{11}$ and on $(F_{11},\psi)$, $6+S'$ on two gluon letters, a letter $F_{11}$ counting $3$ through its two-field part $dA|_V$.
\end{proof}

\begin{tabular}{l|l|l|l}
family & range & sum of the translation-invariant parts along $\Delta$ & spectrum of $D_1$ at $\theta=\theta_\lambda=\theta_V=-\tfrac13$, in units\\
\hline
scalars & bilinears; $\mathfrak{su}(2)$, $L\le8$ & $-\tfrac43L$ & $0$, $-48$; $-8$ times the Heisenberg spectra\\
twist two & $S'\le6$ & $-\tfrac23(4+S')$ & $0$, $-48$, $-\tfrac{200}3$, $-\tfrac{392}5$: $-32h(S)$, $S\le S'$ even\\
$\mathfrak{su}(1|1)$ & $L\le6$ & $-\tfrac43L$ & $-8$ times the $\mathfrak{su}(1|1)$ spectra\\
fermion $\mathfrak{sl}(2)$ & $S'\le5$ & $-\tfrac23(4+S')$ & $-48$, $-\tfrac{200}3$, $-\tfrac{392}5$: $-32h(S+1)$, $S\le S'$ odd\\
$\psi^{(3)}\psi^{(4)}$ with $Z'F_{11}$ & $S'\le2$ & $-\tfrac23(4+S')$ on $\psi^{(3)}\psi^{(4)}$, $-\tfrac23(5+S')$ on $Z'F_{11}$ & $0,-48$; $0,-48,-48,-48$; $0$, four times $-48$, $-\tfrac{200}3$\\
gluon $\mathfrak{sl}(2)$ & $S'\le2$ & $-\tfrac23(6+S')$ & $-48$, $-\tfrac{200}3$: $-32h(S+2)$, $S\le S'$ even\\
$(F_{11},\psi)$ & $S'\le4$ & $-\tfrac23(5+S')$ & at $S'=4$: $-48$, $-48$, $-\tfrac{200}3$, $-\tfrac{200}3$, $-\tfrac{392}5$\\
\end{tabular}
\tabcaption{The composite pair kernels and the spectrum of $D_1$ on the families: the third column is identity (c), the fourth is $-8$ units times the eigenvalues of Proposition \ref{prop:c11-targets}, of Proposition \ref{prop:c10-su2} on the $\mathfrak{su}(2)$ sector and of Table \ref{tab:c11-su11} on the $\mathfrak{su}(1|1)$ sector.\label{tab:c11-results}}

\begin{corollary}[$D_1$ and $P_1$ on the families]
\label{cor:c11-D1fam}
Under the hypotheses of Definition \ref{def:c10-hyp}, on the cohomology of the part with the fewest fields of every family of Definition \ref{def:c11-families}, with every free direction of Definition \ref{def:c11-free},
\[
D_1\cong4\,(\text{the map})-\sum_a\big(\theta_{X_a}+\tfrac13\big)\,\text{units}\cdot v.X_a,\qquad P_1\cong-\sum_a\big(\theta_{P,X_a}+\tfrac13\big)\,\text{units}\cdot\partial_bX_a,
\]
with the map of Theorem \ref{thm:c10-brst}, the position part of $D_1$ and the component $P_1$ taken along $b=n$, and along every $b\in V$ on twist two for $S'\le4$ and on the letters without derivatives (Theorem \ref{thm:c11-families}), where $\theta_{X_a}$ is $\theta$, $\theta_\lambda$ or $\theta_V$ and $\theta_{P,X_a}$ is $\theta_P$, $\theta_{P,\lambda}$ or $\theta_{P,V}$ according to the sector of the letter $X_a$, $v.X_a$ and $\partial_bX_a$ act on the letter $X_a$ alone, and the directions $Q\int(\iota_xA|_V,c^*)$, $Q\int(\rho(x)\lambda^*,\lambda^*)$ contribute $s$-exact terms. At the values $\theta=\theta_\lambda=\theta_V=-\tfrac13$ of Theorem \ref{thm:c9-theta2}, which twist two, the fermion and the gluon $\mathfrak{sl}(2)$ sectors reproduce under the hypotheses of Theorem \ref{thm:c11-arity2} (Proposition \ref{prop:c11-theta-sector}), $D_1$ induces four times the map of Theorem \ref{thm:c10-brst} on every family, with the spectra of Table \ref{tab:c11-results}, $-8$ units times the eigenvalues of $\sum_lH_{l,l+1}$, which is twice the one-loop anomalous dimensions; and $P_1\cong0$ in the translation-covariant scheme, while for equal coefficients $\theta_P$ along $P_b$, $P_1\cong(\theta-\theta_P)$ units times $\partial_b$ with $\theta=-\tfrac13$.
\end{corollary}

\begin{proof}
Corollary \ref{cor:c11-D1local} and Theorem \ref{thm:c11-families}: on the two-field part $\hbar F_1^{(1)}(\Delta)$ acts through its quadratic part, the bracket of the sector part of the sector of $X_a$ with $X_a$ being $v.X_a$ and the others vanishing. As in the proof of Theorem \ref{thm:c11-P7} (i), letter by letter, the position part of the composite pair kernels and of the bracket is $-2(\tfrac13+\theta_{X_a})\partial_b$ units and the translation-invariant part is $-2(\tfrac13+\theta_{X_a})w_{X_a}$ units; along $P_b$ the composite pair kernels give $-\tfrac13$ units times $\partial_b$ and the bracket with $\hbar F_1^{(1)}(P_b)$ gives $-\theta_{P,X_a}$ units times $\partial_bX_a$. The $\kappa$-directions are $s$-exact (Proposition \ref{prop:c11-arity1} (ii)). The rest is Proposition \ref{prop:c11-theta-sector}, Theorem \ref{thm:c11-P7} and Corollary \ref{cor:c11-theta} (iii), (iv).
\end{proof}

\begin{remark}[Method of computation]
\label{rem:c11-method}
Each special-edge composite cycle is evaluated exactly: the Gaussian loop integral and the integrals of \eqref{eq:c11-simplex} over the ordered subsimplices of the region where $\alpha_\sigma$ is the smallest coordinate, the letter derivatives through Proposition \ref{prop:c11-genfun}, the position through the injected momentum $r$ (Proposition \ref{prop:c11-position}), the chains with $k=1$ through Proposition \ref{prop:c11-k1}, the fibre algebra in explicit realisations of $V_{10}$ and $\Sigma$ over $\mathbb Q(i)$, and the colour words by Lemma \ref{lem:c10-colour}, with its planar adjacent rule also checked in realisations of $\mathfrak{su}(2)$, $\mathfrak{su}(3)$ and $\mathfrak{gl}(3)$. In this evaluation every chain carries the sign of Definition \ref{def:c11-sp}, the factor $(-1)^{[c<\sigma]}$ included, and the families are evaluated at generic null leg momenta. A second evaluation works in momentum space with the kernel form of Lemma \ref{lem:c11-kernels} (ii), $v.$ acting on the Fourier components as an Euler operator plus constants, Feynman parameters for the edges $P$ and the special Schwinger time integrated last, applies the sign of Definition \ref{def:c11-sp} as written, and evaluates the families on on-shell legs with momenta proportional to one null vector $\bar n$, $(n,\bar n)\neq0$, which determine the translation-invariant parts and the parts along $n$ (proof of Theorem \ref{thm:c11-families}). The two evaluations give the same off-shell entries with at most two outputs and no letter derivatives, for the pairs of bosonic letters and for the pairs with a gaugino letter or output, whose values lie in $\mathbb Q+\mathbb Q\log2+\mathbb Q\log3$ for $k\le2$, and the same family matrices on the scalars, on twist two for $S'\le4$, on the $\mathfrak{su}(1|1)$ sector, on the fermion $\mathfrak{sl}(2)$ sector for $S'\le5$ and on the gluon $\mathfrak{sl}(2)$ sector for $S'\le2$; the first evaluation extends twist two to $S'\le6$, and the second evaluation gives the two families with $F_{11}$ and a gaugino letter. A third evaluation, a Taylor expansion in momentum space at generic null leg momenta, gives the sums on twist two for $S'\le4$ along a generic $b$. Every integral is a four-dimensional Schwinger integral with the cutoff $\varepsilon$ on every edge.
\end{remark}

\section{The two-body kernels of $Q_1$, $S_1$ and $K_1$}

For $x$ in a homogeneous basis of $\mathfrak{psl}(4|4)$, $\lambda_x$ is its $\mathrm{ad}_\Delta$ eigenvalue, $-1$ on $\Sigma$, $+1$ on $\Sigma^*$ and $+2$ for $K_b$, and the degree of $x$ in the position is $0$, $1$ and $2$ (Section 10.9). The two-body kernel of a pair of letters along $x$ is the kernel of Proposition \ref{prop:c10-arity1}; this section computes it along the supercharges, the special supercharges and the special conformal fields. Along $\Sigma$ the vertex $S^{(1)}(\xi)$ meets the gauge-fixing operator in a new way: the commutator $[\xi.,Q^{GF}]$ joins the gaugino to the ghost, and $S^{(1)}(\xi)$ sits on a cycle as a line from $\lambda$ to $c$ whose image under $Q$ is nonzero, in contrast to the composite lines along $\Delta$ and $P_b$. The section shows that the component along $x$ of every local functional is nevertheless a sum over its pairs of letters of local two-body kernels, the finite parts of the cycles with one vertex $S^{(1)}(x)$ and one special edge, plus a one-body term; writes these kernels as simplex integrals, independent of the endpoint along $K_b$; computes the action-level obstructions and solutions along the odd generators, the rotations and $K_b$ on two legs with all their free directions, and the arity-two equations that relate the free coefficients; assembles $Q_1$, $S_1$, $D_1$, $P_1$, $M_1$, $R_1$ and $K_1$ on the scalar bilinears, where the two-body part of $Q_1$ on the traceless part is $-2\log\tfrac32$ units times the classical transformation, as representatives with their scheme-independent content; and evaluates the pair entries of the twisted differential along a rank-one supercharge. Edges carry $P(\varepsilon,\infty)$ along $\Sigma$ and $\Sigma^*$ (Proposition \ref{prop:c8-Tsafe}) and $P(\varepsilon,T)$ at a finite $T$ along $K_b$ (Definition \ref{def:c8-K}). In this section $\psi$ denotes the finite part of the operators of Definition \ref{def:c8-renorm}, and $L:=\log\tfrac32$.

\begin{definition}[The terms of $S^{(1)}(x)$]
\label{def:c11-terms}
For $\xi\in S_+\otimes W$ the functional $S^{(1)}(\xi)$ of \eqref{eq:c2-S1Q} is the sum of the seven terms of Table \ref{tab:c11-terms}: the quadratic terms $T_1,\dots,T_4$ and the three cubic terms $T_5$, with $F_A=dA+\tfrac12[A\wedge A]$ reduced to $V$ and $dA|_V$, $dA|_{W_{\mathbb R}^\perp}$ the parts of $dA$ in $\wedge^2V$ and in $V\wedge W_{\mathbb R}^\perp$; $T_3$ and $T_4$ carry one derivative on the leg $A$. For $\xi\in S_-\otimes W^*$ the table holds with $\lambda|_{S_+\otimes W}$ and $\lambda|_{S_-\otimes W^*}$ exchanged and $\lambda^*|_{S_+\otimes W^*}$ and $\lambda^*|_{S_-\otimes W}$ exchanged. The quadratic part of $S^{(1)}(\xi)$ is $S^V_2(\xi)+S^\perp_2(\xi)$, with $S^V_2(\xi):=T_1+T_3$, the part with the legs $(\lambda,A^*|_V)$ and $(A|_V,\lambda^*)$, and $S^\perp_2(\xi):=T_2+T_4$, the part with the legs $(\lambda,A^*|_{W_{\mathbb R}^\perp})$ and $(A|_{W_{\mathbb R}^\perp},\lambda^*)$. For $\eta\in\Sigma^*$ and an orthonormal basis $e_\mu$ of $V$, \eqref{eq:c2-S1S} reads
\begin{equation}
S^{(1)}(\eta)=\sum_\mu\int x^\mu s_\mu+\int s_0,\qquad \int s_\mu:=S^{(1)}\big(\rho(e_\mu)\eta\big),\qquad s_0:=-\big(\rho(A|_{W_{\mathbb R}^\perp})\eta,\lambda^*\big);
\label{eq:c11-Seta}
\end{equation}
for $\eta\in S_+\otimes W^*$ the spinors $\rho(e_\mu)\eta$ lie in $S_-\otimes W^*$ and $s_0$ has the legs $A|_{W_{\mathbb R}^\perp}$, $\lambda^*|_{S_+\otimes W^*}$, and for $\eta\in S_-\otimes W$ they lie in $S_+\otimes W$ and $s_0$ has the legs $A|_{W_{\mathbb R}^\perp}$, $\lambda^*|_{S_-\otimes W}$. For an even generator with vector field $v$, $S^{(1)}(x)=\int\big((v.A,A^*)-(v.\lambda,\lambda^*)-(v.c,c^*)\big)$ (Definition \ref{def:c7-vdot}).
\end{definition}

\begin{tabular}{l|l|l}
term & functional, $\int\mathrm{dvol}$ understood & legs\\
\hline
$T_1$ & $-\big(\Gamma(\xi,\lambda)|_V,A^*|_V\big)$ & $\lambda|_{S_-\otimes W^*}$, $A^*|_V$\\
$T_2$ & $-\big(\Gamma(\xi,\lambda)|_{W_{\mathbb R}^\perp},A^*|_{W_{\mathbb R}^\perp}\big)$ & $\lambda|_{S_+\otimes W}$, $A^*|_{W_{\mathbb R}^\perp}$\\
$T_3$ & $\tfrac12\big(\rho(dA|_V)\xi,\lambda^*\big)$ & $A|_V$, $\lambda^*|_{S_+\otimes W^*}$\\
$T_4$ & $\tfrac12\big(\rho(dA|_{W_{\mathbb R}^\perp})\xi,\lambda^*\big)$ & $A|_{W_{\mathbb R}^\perp}$, $\lambda^*|_{S_-\otimes W}$\\
$T_5$ & $\tfrac12\big(\rho(\tfrac12[A\wedge A])\xi,\lambda^*\big)$ & $(A|_V,A|_V,\lambda^*|_{S_+\otimes W^*})$, $(A|_V,A|_{W_{\mathbb R}^\perp},\lambda^*|_{S_-\otimes W})$, $(A|_{W_{\mathbb R}^\perp},A|_{W_{\mathbb R}^\perp},\lambda^*|_{S_+\otimes W^*})$\\
\end{tabular}
\tabcaption{The terms of $S^{(1)}(\xi)$ for $\xi\in S_+\otimes W$ and their legs; their two-component forms are in Remark \ref{rem:c11-twocomp}.\label{tab:c11-terms}}

\begin{lemma}[The generators on the regulator and their lines]
\label{lem:c11-lines}
For an odd generator $x$, $x.$ is the odd operator on $\mathscr E$ that induces the bracket with the quadratic part of $S^{(1)}(x)$ (Proposition \ref{prop:c7-susy}); for an even one it is $v.$ of Definition \ref{def:c7-vdot}. Let $g(H):=\int_\varepsilon^Te^{-tH}dt$, so that $P\star=Q^{GF}g(H)$ and $QP\star+P\star Q=e^{-\varepsilon H}-e^{-TH}=:E(H)$ (Proposition \ref{prop:c9-composite} (iii)).
\begin{itemize}
\item (i) For $\xi\in\Sigma$, $\xi.$ maps the slot of $\lambda$ to the slot of $A$ by $\lambda\mapsto+\Gamma(\xi,\lambda)$, the slot of $A$ to the slot of $\lambda$ by $A\mapsto+\tfrac12\rho(dA)\xi$, the slots of $A^*$ and of $\lambda^*$ to each other by the transposed arrows, and vanishes on the slots of $c$ and $c^*$. In particular
\[
\Big\{S^{(1)}(\xi),\int\omega\,(\iota_{w_1}A,\iota_{w_2}A)\Big\}=\int\omega\,\Big[\big(\iota_{w_1}A,(\rho(w_2)\xi,\lambda)\big)+\big(\iota_{w_2}A,(\rho(w_1)\xi,\lambda)\big)\Big].
\]
\item (ii) $[\xi.,Q]=0$, $[\xi.,H]=0$, and $C_\xi:=[\xi.,Q^{GF}]$ is an even operator of degree $-1$ with constant coefficients and of order one, with four components, each nonzero for $\xi\neq0$: $\lambda\mapsto-d^*\big(\Gamma(\xi,\lambda)|_V\big)$ in the slot of $c$, on a plane wave $e^{i(p,x)}$ the map $\lambda\mapsto+i\big(p,\Gamma(\xi,\lambda)|_V\big)$; a component from $A^*|_V$ to $\lambda$, a nonzero multiple of $A^*|_V\mapsto(p,A^*|_V)\,\xi$; and their transposes from $\lambda^*$ to $A|_V$ and from $c^*$ to $\lambda^*$. For $\xi\in S_+\otimes W$ these are $\lambda|_{S_-\otimes W^*}\to c$, $A^*|_V\to\lambda|_{S_+\otimes W}$, $\lambda^*|_{S_+\otimes W^*}\to A|_V$ and $c^*\to\lambda^*|_{S_-\otimes W}$. The variation kernel $(\xi.\otimes1+1\otimes\xi.)P(\varepsilon,T)=(C_\xi\,g(H)\otimes1)K_0$ has its components in $\lambda\otimes A|_V$, $c\otimes\lambda^*$ and their transposes.
\item (iii) The term $T_1$ with both legs on edges $P$, its leg $\lambda$ on the component $\lambda\otimes\lambda$ and its leg $A^*|_V$ on the component $A^*|_V\otimes c$, is the composite line along $\xi$, with the kernel $2(P\star\,\xi.\,P\star\otimes1)K_0$, whose components lie in $c\otimes\lambda$ and $\lambda\otimes c$. For $\xi\in S_+\otimes W$ and $T=\infty$ it is $-2\,\xi\,e^{-2\varepsilon|p|^2}|p|^{-2}(-,-)^{-1}$ on $c\otimes\lambda|_{S_+\otimes W}$ and $+2\,\xi\,e^{-2\varepsilon|p|^2}|p|^{-2}(-,-)^{-1}$ on $\lambda|_{S_+\otimes W}\otimes c$, and
\[
P\star\,\xi.\,P\star=C_\xi\,Q^{GF}g(H)^2-\xi.\,(Q^{GF})^2g(H)^2,\qquad \big[Q,P\star\,\xi.\,P\star\big]=E(H)\,C_\xi\,g(H)\neq0 .
\]
The terms of $S^{(1)}(\xi)$ with a leg $A^*|_{W_{\mathbb R}^\perp}$ or $\lambda^*$ enter a cycle with that leg on the special edge, and a term $T_5$ also with its leg $\lambda^*$ as an output.
\item (iv) For $\eta\in\Sigma^*$, with $\xi_\mu:=\rho(e_\mu)\eta$,
\[
\eta.=\sum_\mu x^\mu\,\xi_\mu.+\eta^0.+Y_\eta,\qquad [\eta.,Q]=0,\qquad [\eta.,H]=2\sum_\mu\xi_\mu.\,\partial_\mu,
\]
where $\eta^0.$ is the constant operator of order zero induced by $s_0$, with the arrows $A|_{W_{\mathbb R}^\perp}\to\lambda$ and $\lambda^*\to A^*|_{W_{\mathbb R}^\perp}$, and $Y_\eta$ is the constant operator of order zero from $\lambda^*$ to $A^*$ produced when the derivative of $T_3$, $T_4$ in $\sum_\mu\int x^\mu\big(S^V_2(\xi_\mu)+S^\perp_2(\xi_\mu)\big)$ is moved off $A$, $\partial_\nu(x^\mu\,\cdot\,)=x^\mu\partial_\nu+\delta^\mu_\nu$. On a cycle $S^{(1)}(\eta)$ forms the line of (iii) at the spinor $\rho(x)\eta$, linear in the position, and $s_0$ enters with its leg $\lambda^*$ on the special edge.
\item (v) For $K_b$, with $v(x)=(x,x)b-2(b,x)x$ and $\operatorname{div}v=-8(b,x)$, on a slot of Weyl weight $w$ with the rotation generators $\Sigma^{\rho\nu}$, and with $\sigma^\rho(D):=[D,x^\rho]$ for a constant-coefficient operator $D$,
\[
\begin{aligned}
&[v.,Q^{GF}]=4(b,x)\,Q^{GF}+C_K,\qquad C_K:=2\,\sigma^\rho(Q^{GF})\big(w\,b_\rho-b_\nu\Sigma^{\rho\nu}\big),\\
&[v.,H]=4(b,x)\,H-4(w-1)(b,\partial)+4b_\nu\Sigma^{\rho\nu}\partial_\rho,
\end{aligned}
\]
$C_K$ of order zero with constant coefficients, zero on $A^*|_{W_{\mathbb R}^\perp}$, with $C_KA|_V=-4(b,A|_V)$ in the slot of $c$. $S^{(1)}(K_b)$ forms on a cycle the composite line from $c$ to $A|_V$ with the kernel $P\star\,v.\,P\star=v.(P\star)^2+[P\star,v.]P\star$, quadratic in the position, at a finite $T$ (Definition \ref{def:c8-K}).
\end{itemize}
In each case every letter, $A^*|_{W_{\mathbb R}^\perp}$, $\lambda^*$ and $c^*$ included, is active (Lemma \ref{lem:c10-active1}).
\end{lemma}

\begin{proof}
(i) The terms $-(\Gamma(\xi,\lambda),A^*)$ and $\tfrac12(\rho(dA)\xi,\lambda^*)$ of \eqref{eq:c2-S1Q} pair the slots $(\lambda,A^*)$ and $(A,\lambda^*)$; the bracket of a quadratic functional with a coordinate functional is the coordinate functional composed with the corresponding arrow (proof of Lemma \ref{lem:c9-SDelta} (i)), and each pairing term gives one arrow between field slots and the transposed arrow between antifield slots (proof of Proposition \ref{prop:c7-susy}). In the reading of Proposition \ref{prop:c1-Q} the term $(dc,A^*)$ of \eqref{eq:c1-free} gives $c\mapsto-dc$; the term $-(\Gamma(\xi,\lambda),A^*)$, whose letter $\lambda$ is odd like $c$, gives $\lambda\mapsto+\Gamma(\xi,\lambda)$, and $[\xi.,Q]=0$ on the slot of $A$, the part of polynomial degree two of $\{S,S^{(1)}(\xi)\}=0$, fixes $A\mapsto+\tfrac12\rho(dA)\xi$. With $\iota_w\Gamma(\xi,\lambda)=(\rho(w)\xi,\lambda)$ [ESW, (28)] the classical component on the bilinear is the first functional of Corollary \ref{cor:c7-Q}.

(ii) $[\xi.,Q]=0$ and $[\xi.,H]=0$ are Proposition \ref{prop:c7-susy}. $Q^{GF}$ maps $A^*\mapsto-2A^*$ into the slot of $A$, $A\mapsto-d^*(A|_V)$ into the slot of $c$, $c^*\mapsto dc^*$ into the slot of $A^*|_V$ and $\lambda^*\mapsto-{\not d}\lambda^*$ into the slot of $\lambda$, and vanishes on $c$ and $\lambda$ (Chapter 3). Hence $\xi.Q^{GF}$ has the arrows $A^*\to A\to\lambda$, $c^*\to A^*|_V\to\lambda^*$, $\lambda^*\to\lambda\to A$, and $Q^{GF}\xi.$ the arrows $\lambda\to A\to c$, $A^*\to\lambda^*\to\lambda$, $\lambda^*\to A^*\to A$; each composite has order one. On $\lambda$, $C_\xi\lambda=Q^{GF}\Gamma(\xi,\lambda)=-d^*\big(\Gamma(\xi,\lambda)|_V\big)$, and $d^*=-\iota_\partial$ is $-i\iota_p$ on plane waves. On $A^*=a\,e^{i(p,x)}$ the two composites to $\lambda$ are opposite multiples of $\rho(p)\rho(a)\xi$ plus, through $\rho(q(p\wedge a))=\rho(p)\rho(a)-(p,a)$ [ESW, (29)] in the first, a multiple of $(p,a)\xi$; for $a\in W_{\mathbb R}^\perp$ the two cancel. The components from $\lambda^*$ and $c^*$ are the transposes, since $\xi.$ preserves the pairing and $Q^{GF}$ has the adjointness of Proposition \ref{prop:c3-adjoint}, so $(C_\xi\otimes1)K_0$ is graded symmetric. $\Gamma(\xi,-)|_V$ is nonzero on one chirality of $\lambda$ for $\xi\neq0$, and the component from $A^*$ to $\lambda$ is nonzero by Proposition \ref{prop:c7-susy}. The variation kernel is the third item of Proposition \ref{prop:c7-susy} with $T$ in place of $\infty$.

(iii) As in Proposition \ref{prop:c9-composite} (i), with the signs of Lemma \ref{lem:c9-ops} for the odd $\xi.$, both orders of the two legs give $(P\star\,\xi.\,P\star\otimes1)K_0$. Since $[\xi.,g(H)]=0$, $[\xi.,P\star]=C_\xi g(H)$, which gives the first identity; with $[\xi.,Q]=0$,
\[
Q\,(P\star\xi.P\star)=E\,\xi.P\star-P\star Q\,\xi.P\star=E\,\xi.P\star+P\star\xi.E-P\star\xi.P\star Q,
\]
so $[Q,P\star\xi.P\star]=E\,(\xi.P\star+P\star\xi.)=E\,C_\xi g$. The components follow the arrows $\lambda^*\to\lambda\to A\to c$ and $c^*\to A^*|_V\to\lambda^*\to\lambda$. On a plane wave the composite on $c^*$ is a multiple of $\rho(p)\rho(p)\xi=|p|^2\xi$ times $g(H)^2=e^{-2\varepsilon|p|^2}|p|^{-4}$ at $T=\infty$, a scalar line of momentum degree $-2$; the multiple, $-2$ on $c\otimes\lambda$, is the value of both evaluations of Remark \ref{rem:c11-oddmethod}. A leg $A^*|_{W_{\mathbb R}^\perp}$ or $\lambda^*$ has no component of $P$ (Proposition \ref{prop:c3-P}).

(iv) By \eqref{eq:c11-Seta} the field arrows of $\int x^\mu S_2(\xi_\mu)$, $S_2$ the quadratic part, are $x^\mu$ times those of $\xi_\mu.$; its antifield arrows, the transposes, are obtained by integrating by parts the derivative of $T_3$, $T_4$ against $x^\mu\lambda^*$, which gives $x^\mu$ times the antifield arrows of $\xi_\mu.$ plus $Y_\eta$. For the odd $Q$, $[x^\mu\xi_\mu.,Q]=x^\mu[\xi_\mu.,Q]+[Q,x^\mu]\xi_\mu.$, and the vanishing of $[\eta.,Q]$ is the part of polynomial degree two of $\{S,S^{(1)}(\eta)\}=0$ (Lemma \ref{lem:c2-components}); it requires $Y_\eta$. $H$ is $-\partial^2$ on every component, $[x^\mu,-\partial^2]=2\partial_\mu$, and $\eta^0.$, $Y_\eta$ and $\xi_\mu.$ have constant coefficients.

(v) For a conformal Killing field, $\partial_\mu v_\nu=\tfrac14(\operatorname{div}v)\delta_{\mu\nu}+\omega_{\mu\nu}$ with $\omega$ antisymmetric and $v.=v^\mu\partial_\mu+\tfrac12\omega_{\mu\nu}\Sigma^{\mu\nu}+\tfrac w4\operatorname{div}v$; the transport, spin and weight terms and the rotation covariance of a constant-coefficient operator $D$ of weight $-2$ give
\[
\begin{aligned}
&[v.,D]=-\tfrac12(\operatorname{div}v)D-\tfrac w4\sigma^\rho(D)\,\partial_\rho(\operatorname{div}v)-\tfrac12\sigma^\rho(D)(\partial_\rho\omega_{\mu\nu})\Sigma^{\mu\nu},\\
&[v.,H]=-\tfrac12(\operatorname{div}v)H+\tfrac{w-1}2(\partial_\rho\operatorname{div}v)\partial_\rho+(\partial_\rho\omega_{\mu\nu})\Sigma^{\mu\nu}\partial_\rho,
\end{aligned}
\]
with $\partial^2v_\mu=-\tfrac12\partial_\mu\operatorname{div}v$. For $K_b$, $\omega_{\mu\nu}=2(x_\mu b_\nu-x_\nu b_\mu)$ and $\partial_\rho\omega_{\mu\nu}=2(\delta_{\rho\mu}b_\nu-\delta_{\rho\nu}b_\mu)$. On $A^*|_{W_{\mathbb R}^\perp}$, $Q^{GF}=-2$ has order zero; on $A|_V$, $\sigma^\rho(Q^{GF})Y=Y_\rho$, $w=1$ and $\sigma^\rho(Q^{GF})\Sigma^{\rho\nu}A=3A_\nu$, so $C_KA|_V=2(b,A|_V)-6(b,A|_V)$. The line is Lemma \ref{lem:c11-kernels} (iii), (iv) with this $v$.
\end{proof}

\begin{definition}[Defect cycles and two-body kernels]
\label{def:c11-defect}
Let $x$ be a generator, $X_a$, $X_b$ two letters of any slots and $0<\varepsilon<T\le\infty$.
\begin{itemize}
\item (i) A defect cycle of $(X_a,X_b)$ along $x$ is a chain $X_a-e_0-V_1-e_1-\cdots-V_n-e_n-X_b$ of Definition \ref{def:c10-chain} whose vertices are vertices of $I$ and exactly one vertex $V_s$ given by a term of $S^{(1)}(x)$, entered and left through two of its legs, the third leg of a cubic term being an output, and whose edges are components of $P(\varepsilon,T)$ and exactly one special edge $e_\sigma=K_\varepsilon-K_0$, with the components of $K_0$. The vertex $V_s$ is placed on a line, with both legs on edges $P$; at the special edge, with one leg on $e_\sigma$; or cubic, a term $T_5$ with two legs on the cycle.
\item (ii) $\mathcal D_x(X_a,X_b;\varepsilon,T)$ is the sum of the defect cycles of the pair, amputated, with the outputs left open as in Definition \ref{def:c10-pair}: the $\delta$-linear graphs of arity one along $x$ of the flow of the defect vertex of Lemma \ref{lem:c8-transfer}, applied to $e^{(J+\delta M(\varepsilon))/\hbar}$ with $M=\int\omega\,X_aX_b$, whose cycle is a defect cycle of the pair. A chain read from $X_a$ carries, beyond the signs of Lemma \ref{lem:c10-sign}, the factor
\[
-(-1)^{n(<\sigma)+|S^{(1)}(x)|\,(n(<s)+[s<\sigma])},
\]
with $n(<j)$ the number of odd outputs at the vertices before the position $j$, $[s<\sigma]=1$ if $V_s$ precedes $e_\sigma$ and $0$ otherwise, and $|S^{(1)}(x)|$ the total parity of $S^{(1)}(x)$, $0$ for odd $x$ and $1$ for even $x$; the leading sign is that of the special edge (Proposition \ref{prop:c11-master}), and for even $x$ the factor is that of Definition \ref{def:c11-sp} times this sign.
\item (iii) Let $\mathrm{ct}_m(M)$ and $\mathrm{lc}(M)$ be the coefficients of $\varepsilon^{-m}$ and of $\log\varepsilon$ in $\partial_\delta I^{CT}[\varepsilon](M)$, and $\mathrm{ct}^{(1)}_{x,m}(M)$ and $\mathrm{lc}^{(1)}_x(M)$ those of the arity-one counterterm of $M$ along $x$, the $\hbar\,\mathrm{Sing}_\varepsilon$ of the cycles through two letters of $M$ with one vertex $S^{(1)}(x)$ on the cycle and no special edge (Definition \ref{def:c11-graphs} for $x\in\{\Delta,P_b\}$). The finite special vertex of the pair is
\[
\mathcal F_x(X_a,X_b):=-\sum_{m\ge1}\frac{(-1)^m}{m!}\Big(\big\{I,\mathrm{ct}^{(1)}_{x,m}(X_aX_b)\big\}_{H^m}+\big\{S^{(1)}(x),\mathrm{ct}_m(X_aX_b)\big\}_{H^m}\Big),
\]
with $\{-,-\}_{H^m}$ the bracket of \eqref{eq:c10-finitesp}: the coefficient of $\varepsilon^0$ of the special-edge parts of the counterterm brackets $-\{I,\mathrm{ct}^{(1)}_x(X_aX_b)[\varepsilon]\}_\varepsilon$ and $-\{S^{(1)}(x),\partial_\delta I^{CT}[\varepsilon](X_aX_b)\}_\varepsilon$. It is cancelled in the components by the graphs with the special edge on a bridge (Theorem \ref{thm:c11-classes} (ii)); the two-body kernel is the coefficient of $\varepsilon^0$ of the defect cycles alone.
\item (iv) The two-body kernel of $(X_a,X_b)$ along $x$ is
\[
\Upsilon_x(X_a,X_b):=[\varepsilon^0]\,\mathcal D_x(X_a,X_b;\varepsilon,T),
\]
with $[\varepsilon^0]$ the coefficient of $\varepsilon^0$ in the expansion of Theorem \ref{thm:c11-oddsimplex}: a local density of the outputs at the point of the pair, with values in the tensor product of the fibres and colour spaces of the two letters.
\item (v) In arity zero the special-edge cycles of the pair are the chains of Definition \ref{def:c10-chain} with vertices of $I$, edges $P(\varepsilon,T)$ and exactly one special edge, and the arity-zero two-body kernel $\Upsilon_0(X_a,X_b)$ is the coefficient of $\varepsilon^0$ of their sum.
\end{itemize}
For $x\in\{\Delta,P_b\}$ the defect cycles with $V_s$ on a line are the special-edge composite cycles of Definition \ref{def:c11-sp}.
\end{definition}

\begin{theorem}[Two-body form of the components along every generator]
\label{thm:c11-classes}
Let $x$ be a generator and $O=\int\omega\,X_1\cdots X_n$ a monomial local functional, gauge-variant ones included. Let $J=I_{\mathfrak{psl}(4|4)}-I^{CT}_{\mathfrak{psl}(4|4)}[\varepsilon]-\hbar F_1$ with $\hbar F_1$ a one-loop quantisation through arity one along $x$, that is, \eqref{eq:c8-qme} modulo $\hbar^2$ in arity zero and in arity one along $x$, equivalently the first line of \eqref{eq:c9-arity2}, and let $\psi=0$. The endpoint is $T=\infty$ along $\Sigma$, $\Sigma^*$, $\Delta$, $P_b$ and the rotations, and any $T\in(0,\infty)$ along $K_b$.
\begin{itemize}
\item (i) The $\delta$-linear graphs of arity one along $x$ and of order $\hbar$ of the flow of the defect vertex of Lemma \ref{lem:c8-transfer}, applied to $e^{(J+\delta O(\varepsilon))/\hbar}$, form the classes: (F) the one-loop graphs of the vertex $\delta\{S^{(1)}(x),O\}$ and of the vertex $\delta(Q+\{I,-\})O$ with a separate vertex $S^{(1)}(x)$; (C2) the defect cycles of the pairs of letters of $O$, with trees of $I$ on the outputs; (C0) the cycles through two letters of $O$ with the special edge or a contraction through $K_\varepsilon$, with $S^{(1)}(x)$ on a tree; (Br) the graphs with the special edge on a tree; (Ac) the graphs whose cycle passes through vertices of the action alone; (Ct) the trees with one counterterm vertex of order $\hbar$.
\item (ii) (F) is the one-loop part of the arity-one flow of $sO$, and (C0) with the arity-zero counterterm vertices of (Ct) tends to the trees of $\hbar L_1(O)^{(0)}$ with one vertex $S^{(1)}(x)$, both on the left side of \eqref{eq:c8-rep}. A graph of (Br) separates at its special edge, of generic momentum $p$, a one-loop subgraph $G$; when $G$ is divergent the bare graph has the finite limit $\sum_m\frac{(-1)^m}{m!}|p|^{2m}a_m(G)$, $a_m(G)$ the coefficients of $\varepsilon^{-m}$ of $G$, and its counterterm partner, the special-edge part of the counterterm bracket across the same edge, has the opposite limit: the sum is $(e^{-\varepsilon|p|^2}-1)(G-\mathrm{Sing}\,G)$ and tends to $0$. The partners are the special-edge parts of $-\{I,\mathrm{ct}^{(1)}_x(O)[\varepsilon]\}_\varepsilon$, $-\{S^{(1)}(x),\partial_\delta I^{CT}[\varepsilon](O)\}_\varepsilon$ and $-\{I^{CT}[\varepsilon],O\}_\varepsilon$, whose coefficients of $\varepsilon^0$ are the finite special vertices of Definition \ref{def:c11-defect} (iii) and the vacuum-pole terms of \eqref{eq:c10-finitesp}; they are used up here. (Ac) with the defect vertices of the action tends to $0$ by the master equation in arity zero and in arity one along $x$.
\item (iii) (Master formula.) With the Koszul signs of \eqref{eq:c10-decomp},
\begin{equation}
\begin{aligned}
\hbar L_1(O)^{(1)}(x)=\ &\sum_{1\le a<b\le n}(-1)^{|X_b|(|X_{a+1}|+\cdots+|X_{b-1}|)}\int\omega\,X_1\cdots X_{a-1}\,\Upsilon_x(X_a,X_b)\,X_{a+1}\cdots\widehat{X_b}\cdots X_n\\
&-\big\{\hbar F_1^{(1)}(x),O\big\},
\end{aligned}
\label{eq:c11-decompx}
\end{equation}
the two-body kernels $\Upsilon_x=[\varepsilon^0]\mathcal D_x$ of Definition \ref{def:c11-defect} summed over the unordered pairs of letters, and a one-body term, the bracket with the arity-one finite counterterm along $x$. In arity zero, with the same signs,
\begin{equation}
\hbar L_1(O)^{(0)}=\sum_{a<b}\pm\int\omega\cdots\Upsilon_0(X_a,X_b)\cdots-\big\{\hbar F_1^{(0)},O\big\}.
\label{eq:c11-decomp0}
\end{equation}
Every number of a component is a coefficient of $\varepsilon^0$ of a defect cycle or a coefficient of the action-side solution $\hbar F_1$; these formulas are the definition of the components in the construction.
\item (iv) The coefficients of $\log\varepsilon$ and of $\varepsilon^{-m}$ cancel. The defect cycles carry no $\log\varepsilon$ (Theorem \ref{thm:c11-oddsimplex}), and the first cancellation is the logarithmic covariance
\begin{equation}
\mathrm{lc}\big(\{S^{(1)}(x),O\}\big)-\big\{S^{(1)}(x),\mathrm{lc}(O)\big\}=s\,\mathrm{lc}^{(1)}_x(O)-\mathrm{lc}^{(1)}_x(sO)+\big\{L^{(1)}(x),O\big\},
\label{eq:c11-logcov}
\end{equation}
with $L^{(1)}(x)$ the coefficient of $\log\varepsilon$ in $I^{CT}_{\mathfrak{psl}(4|4)}[\varepsilon]^{(1)}(x)$, which has no term $\varepsilon^{-m}$ along $\Sigma$, $\Sigma^*$, $\Delta$, $P_b$ and $K_b$. The second is the pole identity, for $m\ge1$,
\begin{equation}
\sum_{a<b}\pm[\varepsilon^{-m}]\,\mathcal D_x(X_a,X_b)=s\,\mathrm{ct}^{(1)}_{x,m}(O)+\big\{S^{(1)}(x),\mathrm{ct}_m(O)\big\}-\mathrm{ct}_m\big(\{S^{(1)}(x),O\}\big)-\mathrm{ct}^{(1)}_{x,m}(sO),
\label{eq:c11-pole}
\end{equation}
the special-edge parts of the counterterm brackets, with all their powers of $\varepsilon$, having been grouped with the bridges in (ii); for a single pair $O=\int\omega\,X_aX_b$ the right side contains the poles of the pairs of the cubic monomials of $\{S^{(1)}(x),O\}$ and of $sO$ times their third letter.
\item (v) \eqref{eq:c11-decompx} holds along $K_b$ at every $T\in(0,\infty)$ with the same right side. Given the master equation through arity one along $x$, the local representative of $L(O)$ exists through arity one along $x$ for every local functional $O$: (iii) of Definition \ref{def:c10-hyp} holds along every generator.
\end{itemize}
\end{theorem}

\begin{proof}
(i) At order $\hbar$ a graph of the flow has one cycle and vertices of order $\hbar^0$, or no cycle and one vertex of order $\hbar$ (Proposition \ref{prop:c3-graphs}); the part of arity one along $x$ and linear in $\delta$ places $x$ on one vertex and $\delta$ on one vertex. Expand the defect vertex in $\delta$, in the arity and in $\hbar$, with $I_{\mathfrak{psl}(4|4)}=I+S^{(1)}+S^{(2)}$ and $O(\varepsilon)=O-\partial_\delta I^{CT}_{\mathfrak{psl}(4|4)}[\varepsilon](O)$. The classical master equation \eqref{eq:c2-mc} leaves at order $\hbar^0$ the vertices $\delta\{S^{(1)}(x),O\}$ and $\delta(Q+\{I,-\})O$ and the special-edge parts of $\{S^{(1)}(x),O\}_\varepsilon$, $\{I,O\}_\varepsilon$, $\{I,S^{(1)}(x)\}_\varepsilon$, $\tfrac12\{I,I\}_\varepsilon$, together with $\delta\hbar\Delta_\varepsilon O$; $\hbar\Delta_\varepsilon S^{(1)}(x)$ vanishes, its quadratic terms pairing a field with the antifield of another field and its cubic terms carrying $\lambda^*$ and no $\lambda$. A cycle through a vertex of the first two kinds is in (F). A cycle with a special edge passes through $O$ along two letters (Theorem \ref{thm:c10-decomp} (i)) or avoids $O$, and the positions of $S^{(1)}(x)$ and of the special edge sort it into (C2), (C0), (Br) or (Ac).

(ii) By Theorem \ref{thm:c8-L} (iv) the left side of \eqref{eq:c8-rep} in arity one along $x$ is $\big(L(O)^{(1)}(x)\big)[T]+\big(L(O)^{(0)}\big)[T]^{(1)}(x)$ with $L(O)=sO+\hbar L_1(O)$; the one-loop part of $(sO)[T]^{(1)}(x)$ consists of the graphs (F) and their counterterm trees. For (C0), the arity-zero analysis of Theorem \ref{thm:c10-decomp} and of the present argument at arity zero applies with a tree carrying $S^{(1)}(x)$ attached at generic momenta. (Br): the special edge on a tree edge of momentum $p$ multiplies the subgraph $G$ it separates by $e^{-\varepsilon|p|^2}-1=\sum_{j\ge1}(-\varepsilon|p|^2)^j/j!$. For $G=\sum_{m\ge1}\varepsilon^{-m}a_m+a_0+o(1)$, a cycle through two letters of $O$ with $S^{(1)}(x)$ on it or a cycle of the action such as the self-energy with the pole of $I^{CT}$, the terms $j=m$ give the finite limit stated. The counterterm vertex of $G$ joined by the same special edge is $-(e^{-\varepsilon|p|^2}-1)\sum_m\varepsilon^{-m}a_m$; its coefficient of $\varepsilon^0$ is $-\sum_m\frac{(-1)^m}{m!}|p|^{2m}a_m$, which in the bracket notation of \eqref{eq:c10-finitesp} is $-\frac{(-1)^m}{m!}\{f,g_m\}_{H^m}$ with $(f,g_m)=(I,\mathrm{ct}^{(1)}_{x,m}(O))$, $(S^{(1)}(x),\mathrm{ct}_m(O))$ or $(\text{coefficient of }\varepsilon^{-m}\text{ in }I^{CT}[\varepsilon],O)$. The sum of a bridge and its partner is $(e^{-\varepsilon|p|^2}-1)$ times the renormalised subgraph, which has a limit (Definition \ref{def:c8-quant}), and tends to $0$; the partners with $\log\varepsilon$ carry $(e^{-\varepsilon|p|^2}-1)\log\varepsilon\to0$ (Theorem \ref{thm:c8-obstruction} (ii), Proposition \ref{prop:c4-defect1}). Along $K_b$ at finite $T$ the subgraphs carry at fixed $\varepsilon$ the term $\tfrac3{8\pi^2}\hbar\log T$ times a tree with the same factor (remark after Theorem \ref{thm:c8-K}), which also tends to $0$. (Ac): the cycle and its counterterm vertices form the flow of the defect of the action in arity zero or in arity one along $x$, joined to $O$ by a tree at generic momenta, whose limit is the tree-level flow of $\hbar\,\mathrm{Ob}-s(\hbar F_1)$ in these arities (Definition \ref{def:c8-ob}), zero by the master equation.

(iii) What remains is (C2) with the arity-one vertices of (Ct) without special edge, against the tree of $\hbar L_1(O)^{(1)}(x)$ and the counterterm trees of $-\partial_\delta I^{CT}[\varepsilon](\{S^{(1)}(x),O\})$ and $-\mathrm{ct}^{(1)}_x(sO)$ on the left. The vertex $-\{\hbar F_1^{(1)}(x),O\}_\varepsilon$ gives the one-body term, its special-edge part tending to $0$. Each defect cycle is a sum of poles, $[\varepsilon^0]\mathcal D_x$ and a vanishing remainder (Theorem \ref{thm:c11-oddsimplex}); the trees on its outputs converge, and Theorem \ref{thm:c8-L} (i) identifies the limit with the tree-level flow of \eqref{eq:c11-decompx}. The sum over the pairs with the Koszul signs is that of Theorem \ref{thm:c10-decomp} (iii), each defect cycle meeting $O$ in exactly two letters. In arity zero the same argument with the special-edge cycles of the pairs, the vertex $-\{\hbar F_1^{(0)},O\}_\varepsilon$ and the bridges through the self-energy cycles grouped with the special-edge parts of $-\{I^{CT}[\varepsilon],O\}_\varepsilon$ gives \eqref{eq:c11-decomp0}.

(iv) The left side has a limit; collecting the coefficients of $\log\varepsilon$ and of $\varepsilon^{-m}$ of (C2), of the brackets without special edge and of the counterterm trees moved from the left side gives \eqref{eq:c11-logcov} and \eqref{eq:c11-pole}, the special-edge parts of the counterterm brackets cancelling against the bridges at every power of $\varepsilon$ by (ii). A coefficient of $\varepsilon^{-m}$ in $I^{CT}_{\mathfrak{psl}(4|4)}[\varepsilon]^{(1)}(x)$ is a colour-singlet local functional of degree $-1$ and total Weyl weight $4-\tfrac12\lambda_x-2m$, the position at weight $-1$ and of degree $p$ at most that of $x$ (Proposition \ref{prop:c10-bound1}); a colour singlet of degree $-1$ contains an antifield of weight at least $\tfrac52$ and a second adjoint field of weight at least $1$, or $c^*$ and $c$, so its density has weight at least $\tfrac72$ and the total weight is at least $\tfrac72-p$. For $m\ge1$ the total weight is at most $\tfrac52$, $\tfrac32$, $2$, $3$, $1$ along $\Sigma$, $\Sigma^*$, $\Delta$, $P_b$, $K_b$, below $\tfrac72$, $\tfrac52$, $\tfrac52$, $\tfrac72$, $\tfrac32$. Read for one pair, the cubic terms of $\{S^{(1)}(x),O\}$ and of $sO$ have pairs with the empty leg set of Theorem \ref{thm:c10-decomp} (iv), whose poles stand on the right with the third letter.

(v) The argument uses $sO$ through (F), which cancels between the two sides, and through the counterterm trees. Along $K_b$ at a fixed finite $T$ the edges $P(\varepsilon,T)$ are bounded uniformly in $T$ and decay as $e^{-|y|^2/4T}$ (Definition \ref{def:c8-PT}); $\Upsilon_{K_b}$ is independent of $T$ by Theorem \ref{thm:c11-oddsimplex} (iii), and Theorem \ref{thm:c8-K} gives the independence of the local representative. The classes (C2) and (Ct) have local limits, so Theorem \ref{thm:c8-L} (i) gives the local representative.
\end{proof}

\begin{theorem}[Defect cycles as simplex integrals]
\label{thm:c11-oddsimplex}
For a defect cycle with $n$ vertices and the edges $e_0,\dots,e_n$, let $q_j$, the centre, the momentum form and the $m$-th simplex integrand $F_m$ be those of Definition \ref{def:c10-cycle} with $n$ in place of $k$, the numerator containing the function of the term of $S^{(1)}(x)$ at $V_s$ (Table \ref{tab:c11-terms}, Definition \ref{def:c7-vdot}) and the special-edge factor $-|l+q_\sigma|^2$ times the fibre factor of $K_0$; along $\eta$ and $K_b$ a momentum $r$ is injected at $V_s$ and carried by $e_s,\dots,e_n$ to $X_b$ and $\omega$. Let $D$ be the superficial degree of Proposition \ref{prop:c11-oddbound}.
\begin{itemize}
\item (i) At generic output momenta and $T=\infty$, for $x\in\Sigma\cup\Sigma^*\cup\{\Delta,P_b\}\cup\mathfrak{so}(V)$, the value of the cycle is $\sum_{m\ge1}\varepsilon^{-m}a_m+a_0+o(1)$ with
\begin{equation}
\begin{aligned}
16\pi^2a_m&=\frac1m\int_{\alpha_\sigma\text{ smallest}}\Big[\big(\min_{f\ne\sigma}\alpha_f\big)^m-\alpha_\sigma^m\Big]F_m\,d\alpha\qquad(m\ge1),\\
16\pi^2a_0&=\int_{\alpha_\sigma\text{ smallest}}F_0\,\log\frac{\min_{f\ne\sigma}\alpha_f}{\alpha_\sigma}\,d\alpha,
\end{aligned}
\label{eq:c11-oddsimplex}
\end{equation}
and no term in $\log\varepsilon$; $a_m$ is a polynomial of degree $D-2m$ in the output momenta and in $r$, and $a_0=0$ for $D<0$.
\item (ii) Along $\eta\in\Sigma^*$, at the point $x_O$ of the pair, $\Upsilon_\eta=\sum_\mu x_O^\mu\,\Upsilon_{\rho(e_\mu)\eta}+\Upsilon^0_\eta$: the position part along $b$ is the kernel along $\rho(b)\eta\in\Sigma$, and the translation-invariant part $\Upsilon^0_\eta$ is $-i\sum_\mu\partial_{r_\mu}$ at $r=0$ of \eqref{eq:c11-oddsimplex} with $S^{(1)}(\rho(e_\mu)\eta)$ at $V_s$, plus \eqref{eq:c11-oddsimplex} with $s_0$ at $V_s$.
\item (iii) Along $K_b$, with $z=x-x_O$,
\[
v(x)=v(x_O)+\big(-2(b,x_O)z+\omega_{x_O}z\big)+v(z),\qquad \omega_{x_O}z:=2\big((x_O,z)b-(b,z)x_O\big),
\]
and correspondingly
\begin{equation}
\Upsilon_{K_b}=v(x_O)^\mu\,\Upsilon_{P_{e_\mu}}-(b,x_O)\,\Upsilon^0_\Delta+\Upsilon^0_{M_{\omega_{x_O}}}+\Upsilon^0_{K_b},
\label{eq:c11-Kdec}
\end{equation}
with $\Upsilon^0_y$ the translation-invariant part along the generator $y$ centred at $x_O$, the rotation $M_{\omega_{x_O}}$ having the vector field $\omega_{x_O}z$; the monomials $z^\mu$ and $z^\mu z^\nu$ of the Killing data act as $-i\partial_{r_\mu}$ and $-\partial_{r_\mu}\partial_{r_\nu}$ at $r=0$. At every finite $T$ the coefficients of the expansion in $\varepsilon$ are given by \eqref{eq:c11-oddsimplex}, and $T\partial_T$ of every defect cycle is $O(\varepsilon)$: $\Upsilon_{K_b}$ is independent of $T$.
\item (iv) For a cycle with $n_t=n+1$ edges and a constant integrand $F_0$ ($D=0$), $16\pi^2a_0=F_0\log\frac{n_t}{n_t-1}\big/(n_t-1)!$, so that $16\pi^2a_0=\tfrac12F_0\log\tfrac32$ for $n_t=3$ and $\tfrac16F_0\log\tfrac43$ for $n_t=4$; for a constant $F_1$ ($D=2$), $16\pi^2a_1=F_1\big/\big(n_t^2(n_t-1)(n_t-1)!\big)$. For $n=1$, the vertex $V_s$ at the special edge, with $\alpha:=\alpha_\sigma\in[0,\tfrac12]$,
\[
16\pi^2a_0=\int_0^{1/2}F_0(\alpha)\log\frac{1-\alpha}\alpha\,d\alpha,\qquad 16\pi^2a_m=\frac1m\int_0^{1/2}\big((1-\alpha)^m-\alpha^m\big)F_m(\alpha)\,d\alpha,
\]
in $\mathbb Q+\mathbb Q\log2$ for polynomial numerators; for $n=2$ with $V_s$ on a line between a letter or a vertex of $I$ and the special edge, Proposition \ref{prop:c11-k1} applies, with values in $\mathbb Q+\mathbb Q\log2+\mathbb Q\log3$.
\end{itemize}
\end{theorem}

\begin{proof}
(i) The domain $\{t_\sigma\le\varepsilon,\ t_f\ge\varepsilon\ (f\ne\sigma)\}$ is the difference of the domains with the cutoff shapes $1/\min_{f\ne\sigma}\alpha_f$ and $1/\min_f\alpha_f$, which differ on the region where $\alpha_\sigma$ is the smallest coordinate; Proposition \ref{prop:c10-shape} gives for each shape the poles, the same coefficient of $\log\varepsilon$ and the finite part, and the difference is \eqref{eq:c11-oddsimplex}, as in Lemma \ref{lem:c10-special}. The region $\tau\ge1$ converges at $T=\infty$: the edges carry distinct generic momenta (proof of Theorem \ref{thm:c10-decomp} (ii)), the composite line along $\xi$ is of order $|p|^{-2}$ at small momentum (Lemma \ref{lem:c11-lines} (iii)), the composite kernels along $\Delta$, $P_b$ and $\mathfrak{so}(V)$ of order $|p|^{-3}$ (Proposition \ref{prop:c9-composite} (iv)), and along $\eta$ the line carries the spinor $\rho(b)\eta$ in the position part and one derivative in $r$ in the translation-invariant part; the measure of $\{\alpha_\sigma\le\varepsilon/\tau\}$ is of order $\varepsilon$. The coefficient of $\tau^{-1-m}$ after the Gaussian integral is homogeneous of degree $D-2m$. (ii) By \eqref{eq:c11-Seta}, insert $x^\mu=x_O^\mu+(x-x_O)^\mu$ and argue as in Proposition \ref{prop:c11-position} (i): on plane waves the factor $(x-x_O)^\mu$ is $-i\partial/\partial r_\mu$ of the momentum injected at $V_s$ and absorbed by $\omega$, which commutes with the Gaussian integral and the simplex integration; the same holds for the special edges, the counterterm vertices and the finite special vertices. (iii) Expand $(x_O+z,x_O+z)b-2(b,x_O+z)(x_O+z)$; $S^{(1)}$ of Definition \ref{def:c7-vdot} is linear in $v$ with its weight and spin terms, and $\operatorname{div}v(x)=-8(b,x_O)-8(b,z)$, $\operatorname{div}(2z)=8$, $\operatorname{div}(\omega_{x_O}z)=0$. At a finite $T$ the region $\tau\ge\tau_0$ is compact in $\tau$, so the expansion in $\varepsilon$ is that of the region of small $\tau$. $T\partial_T$ replaces one edge $e$ by $TL_T$, of scalar factor $Te^{-T|l+q_e|^2}$, which confines the loop momentum, while the special edge carries $e^{-\varepsilon|l+q_\sigma|^2}-1=O\big(\varepsilon(1+|l|)^2\big)$; the polynomial dependence on the position enters through derivatives in $r$, which preserve these bounds, and the finite special vertices are built from the coefficients $\mathrm{ct}_m$, $\mathrm{ct}^{(1)}_{x,m}$, independent of $T$ (Proposition \ref{prop:c8-Tindep}). (iv) For a constant integrand the integrals of \eqref{eq:c11-oddsimplex} over the region where $\alpha_\sigma$ is the smallest coordinate are elementary; for $n=1$ the simplex is $[0,1]$, the region is $\alpha_\sigma\le\tfrac12$ and $\min_{f\ne\sigma}\alpha_f=1-\alpha_\sigma$, and $\int_0^{1/2}\alpha^a(1-\alpha)^b\log\frac{1-\alpha}\alpha\,d\alpha\in\mathbb Q+\mathbb Q\log2$. Proposition \ref{prop:c11-k1} uses that the two segments adjacent to $V_s$ carry the same momentum at $r=0$ and that the special edge is the third edge.
\end{proof}

\begin{example}[The cubic term at the special edge]
\label{ex:c11-cubic}
The defect cycles with $n=1$ and an output have a cubic term of $S^{(1)}(\xi)$ at the special edge. For $\xi\in S_+\otimes W$ and the pair $(A|_{W_{\mathbb R}^\perp},\lambda|_{S_+\otimes W})$, the cycle
\[
A|_{W_{\mathbb R}^\perp}\xrightarrow{\ P\ }T_5\big(A|_{W_{\mathbb R}^\perp},A|_{W_{\mathbb R}^\perp},\lambda^*|_{S_+\otimes W^*}\big)\xrightarrow{\ K_\varepsilon-K_0\ }\lambda|_{S_+\otimes W},
\]
with the output $A|_{W_{\mathbb R}^\perp}$ of momentum $q$, has by Theorem \ref{thm:c11-oddsimplex} (iv) and Lemma \ref{lem:c10-special} the value
\[
\frac{\hbar}{16\pi^2}\Big(-\frac1{2\varepsilon}-\frac38|q|^2\Big)
\]
times the value of $T_5$ on the fibre elements and colour factors of the letters and of the output, with the fibre factor $2(-,-)^{-1}$ of the edge $A\otimes A$ and the inverse fibre pairing of $(\lambda,\lambda^*)$ on the special edge. Its contribution to $\Upsilon_\xi\big(A|_{W_{\mathbb R}^\perp},\lambda|_{S_+\otimes W}\big)$ on the leg set $A|_{W_{\mathbb R}^\perp}$ is $\tfrac38\frac{\hbar}{16\pi^2}$ times $\partial^2$ of the output times that factor, and its pole enters \eqref{eq:c11-pole} against $\mathrm{ct}_1$ of the pair $(\lambda|_{S_+\otimes W},\lambda|_{S_+\otimes W})$ of the classical transformation $\{S^{(1)}(\xi),A|_{W_{\mathbb R}^\perp}\}=\Gamma(\xi,\lambda)|_{W_{\mathbb R}^\perp}$, whose coefficient of $\varepsilon^{-1}$ has the leg set $A|_{W_{\mathbb R}^\perp}$ (Table \ref{tab:c10-legsets}).
\end{example}

\begin{proposition}[Output bounds along $\Sigma$, $\Sigma^*$ and $K_b$]
\label{prop:c11-oddbound}
A defect cycle of $(X_a,X_b)$ along $x$ whose outputs have total Weyl weight $w$, fields and derivatives, has the superficial degree
\[
D=w_{X_a}+w_{X_b}+\delta-w,
\]
with $\delta=\tfrac12$ along $\xi\in\Sigma$ and for the position part along $\eta$, $\delta=-\tfrac12$ for $\Upsilon^0_\eta$, and $\delta=1,0,-1$ for the parts of $\Upsilon_{K_b}$ multiplying $v(x_O)$, the first derivatives of $v$ at $x_O$ and its second derivatives. Hence:
\begin{itemize}
\item (i) along $\xi\in\Sigma$: $\Upsilon_\xi(X_a,X_b)$ is translation invariant, of Weyl weight $w_{X_a}+w_{X_b}+\tfrac12$, of the pair degree and of parity $|X_a|+|X_b|+1$, with at most $w_{X_a}+w_{X_b}+\tfrac12$ plus the pair degree output fields, as for the second-order operator on pairs of letters of [BK, (3.7)];
\item (ii) along $\eta\in\Sigma^*$: $\Upsilon_\eta=\sum_\mu x^\mu\,\Upsilon_{\rho(e_\mu)\eta}+\Upsilon^0_\eta$, of parity $|X_a|+|X_b|+1$, with $\Upsilon^0_\eta$ translation invariant of Weyl weight $w_{X_a}+w_{X_b}-\tfrac12$ and at most $w_{X_a}+w_{X_b}-\tfrac12$ plus the pair degree output fields;
\item (iii) along $K_b$: $\Upsilon_{K_b}$ is a polynomial of degree at most two in the position, of Weyl weight $w_{X_a}+w_{X_b}-1$ with the position at weight $-1$, of the pair degree and of parity $|X_a|+|X_b|$, with at most $w_{X_a}+w_{X_b}+1$ plus the pair degree output fields; $\Upsilon^0_{K_b}$ has Weyl weight $w_{X_a}+w_{X_b}-1$.
\end{itemize}
The same bounds hold for the counterterm cycles, whose outputs carry the pair degree minus one, and for the finite special vertices, which are supported on the leg sets with $D\ge2$ and cancel against the bridges in the components (Theorem \ref{thm:c11-classes} (ii)). These are the bounds of Proposition \ref{prop:c10-bound1} with $\lambda_x=-1,+1,+2$, valid without (iii) of Definition \ref{def:c10-hyp} by Theorem \ref{thm:c11-classes} (v).
\end{proposition}

\begin{proof}
Lemma \ref{lem:c10-degree} with one vertex of weight $4+\delta$: every vertex of $I$ has weight $4$, the term of $S^{(1)}(x)$ has weight $4+\delta$ on its fields and derivatives once the Killing data carry the weights of Proposition \ref{prop:c7-select} ($-\tfrac12$ for a constant spinor, $-1$, $0$, $+1$ for the value and the first and second derivatives of a vector field, the position at weight $-1$), every edge $P$ has numerator degree $w_1+w_2-2$, and the special edge, $-|l+q|^2$ over a Schwinger interval $[0,\varepsilon]$, has degree $0=w+(4-w)-4$. The degree of the coefficient of $\varepsilon^0$ is Theorem \ref{thm:c11-oddsimplex} (i). $S^{(1)}(x)$ has degree $-1$ and total parity $0$ for odd $x$ and $1$ for even $x$, the special edge degree $+1$ and parity $1$. Weyl weight plus degree is at least $1$ for every field and antifield, with equality for $c$ and $A$ (Proposition \ref{prop:c10-bound}), which bounds the number of outputs. The position parts are Theorem \ref{thm:c11-oddsimplex} (ii), (iii). On a leg set of degree $D$ the term of $\mathcal F_x$ with $\mathrm{ct}_m$ or $\mathrm{ct}^{(1)}_{x,m}$ contains a coefficient of $\varepsilon^{-m}$, of degree $D-2m\ge0$ in the momenta, and $m\ge1$.
\end{proof}

\begin{proposition}[Census of the defect cycles]
\label{prop:c11-oddcensus}
\begin{itemize}
\item (i) Every defect cycle conserves the degree and the centre $\mathbb Z/4$ of $SL(4,\mathbb C)$, $W\mapsto iW$, once $x$ carries its charge, $+1$ on $S_+\otimes W$ and on $S_-\otimes W$, $-1$ on $S_-\otimes W^*$ and on $S_+\otimes W^*$, $0$ on the even part; it shifts the total parity by that of $S^{(1)}(x)$ plus one.
\item (ii) Along every generator, defect cycles exist for $45$ of the $55$ unordered pairs of the ten slots; the other ten are the pairs of two letters among $A^*|_{W_{\mathbb R}^\perp}$, $\lambda^*|_{S_+\otimes W^*}$, $\lambda^*|_{S_-\otimes W}$, $c^*$, each of which needs the special edge.
\item (iii) A defect cycle without vertices of $I$ has $V_s$ at the special edge, letter--$P$--$V_s$--special edge--letter; a cycle without outputs is homogeneous in $\varepsilon$ of degree $-\tfrac D2$ (Lemma \ref{lem:c10-self}), so its coefficient of $\varepsilon^0$ is nonzero at $D=0$ alone.
\item (iv) Along $\xi\in S_+\otimes W$ the leg sets through the pairs of field and ghost letters without derivatives are those of Table \ref{tab:c11-defectpairs}; the numbers of leg sets with $n\le3$ derivatives on the two letters, along $\xi\in S_+\otimes W$ and for the translation-invariant parts along $\eta\in S_+\otimes W^*$ and along $K_b$, are those of Table \ref{tab:c11-defectcounts}. Along $S_-\otimes W^*$ and $S_-\otimes W$ the tables hold with the chiralities exchanged, and for the position parts with $\delta$ of Proposition \ref{prop:c11-oddbound}.
\item (v) Along $\xi\in S_+\otimes W$ the pair $(c,c)$ has the empty set of leg sets at the pair weight without derivatives, and the pairs of $c$ with $A|_V$, $A|_{W_{\mathbb R}^\perp}$, $\lambda|_{S_+\otimes W}$, $\lambda|_{S_-\otimes W^*}$ have leg sets with one ghost output.
\end{itemize}
\end{proposition}

\begin{proof}
(i) Each vertex of $I$, each term of $S^{(1)}(x)$ and each component of $P$ and of $K_0$ carries these charges (Table \ref{tab:c11-terms}, Definition \ref{def:c1-vertex}, Proposition \ref{prop:c3-P}). (ii)--(v) Exhaustive enumeration of the chains with up to six vertices on the cycle, exactly one term of $S^{(1)}(x)$ in each of its placements and exactly one special edge, with output weights up to $w_{X_a}+w_{X_b}+3+\delta$; every chain is checked for degree, parity and the centre, and its superficial degree against the count of Lemma \ref{lem:c10-degree}. A letter $A^*|_{W_{\mathbb R}^\perp}$, $\lambda^*$ or $c^*$ has no component of $P$ and lies on the special edge, of which a defect cycle has one.
\end{proof}

\begin{tabular}{l|c|l}
pair & output weight at most & leg sets along $\xi\in S_+\otimes W$: fewest vertices of $I$; number of chains\\
\hline
$(A|_V,A|_V)$ & $5/2$ & $\lambda|_{S_-\otimes W^*}$ $(1;2)$, $A|_V\,\lambda|_{S_-\otimes W^*}$ $(1;8)$, $A|_{W_{\mathbb R}^\perp}\,\lambda|_{S_+\otimes W}$ $(1;7)$, $c\,\lambda^*|_{S_+\otimes W^*}$ $(1;4)$\\
$(A|_V,A|_{W_{\mathbb R}^\perp})$ & $5/2$ & $\lambda|_{S_+\otimes W}$ $(1;4)$, $A|_V\,\lambda|_{S_+\otimes W}$ $(1;15)$, $A|_{W_{\mathbb R}^\perp}\,\lambda|_{S_-\otimes W^*}$ $(1;15)$, $c\,\lambda^*|_{S_-\otimes W}$ $(1;6)$\\
$(A|_V,\lambda|_{S_+\otimes W})$ & $3$ & $\emptyset$ $(0;2)$, $A|_V$ $(0;6)$, $A|_V^2$ $(1;13)$, $(A|_{W_{\mathbb R}^\perp})^2$ $(1;11)$, $A|_V^3$ $(1;25)$, $A|_V\,(A|_{W_{\mathbb R}^\perp})^2$ $(1;55)$, $c\,A^*|_V$ $(2;9)$, $\lambda|_{S_+\otimes W}\,\lambda|_{S_-\otimes W^*}$ $(2;15)$\\
$(A|_V,\lambda|_{S_-\otimes W^*})$ & $3$ & $A|_{W_{\mathbb R}^\perp}$ $(0;5)$, $A|_V\,A|_{W_{\mathbb R}^\perp}$ $(1;19)$, $A|_V^2\,A|_{W_{\mathbb R}^\perp}$ $(1;51)$, $(A|_{W_{\mathbb R}^\perp})^3$ $(1;19)$, $c\,A^*|_{W_{\mathbb R}^\perp}$ $(2;2)$, $(\lambda|_{S_-\otimes W^*})^2$ $(2;7)$, $(\lambda|_{S_+\otimes W})^2$ $(2;6)$\\
$(A|_V,c)$ & $3/2$ & $\lambda|_{S_-\otimes W^*}\,c$ $(2;10)$\\
$(A|_{W_{\mathbb R}^\perp},A|_{W_{\mathbb R}^\perp})$ & $5/2$ & $\lambda|_{S_-\otimes W^*}$ $(1;2)$, $A|_V\,\lambda|_{S_-\otimes W^*}$ $(1;7)$, $A|_{W_{\mathbb R}^\perp}\,\lambda|_{S_+\otimes W}$ $(1;8)$, $c\,\lambda^*|_{S_+\otimes W^*}$ $(1;2)$\\
$(A|_{W_{\mathbb R}^\perp},\lambda|_{S_+\otimes W})$ & $3$ & $A|_{W_{\mathbb R}^\perp}$ $(0;6)$, $A|_V\,A|_{W_{\mathbb R}^\perp}$ $(1;22)$, $A|_V^2\,A|_{W_{\mathbb R}^\perp}$ $(1;59)$, $(A|_{W_{\mathbb R}^\perp})^3$ $(1;21)$, $c\,A^*|_{W_{\mathbb R}^\perp}$ $(2;7)$, $(\lambda|_{S_-\otimes W^*})^2$ $(2;7)$, $(\lambda|_{S_+\otimes W})^2$ $(2;8)$\\
$(A|_{W_{\mathbb R}^\perp},\lambda|_{S_-\otimes W^*})$ & $3$ & $\emptyset$ $(0;2)$, $A|_V$ $(0;5)$, $A|_V^2$ $(1;10)$, $(A|_{W_{\mathbb R}^\perp})^2$ $(1;11)$, $A|_V^3$ $(1;18)$, $A|_V\,(A|_{W_{\mathbb R}^\perp})^2$ $(1;52)$, $\lambda|_{S_+\otimes W}\,\lambda|_{S_-\otimes W^*}$ $(2;13)$\\
$(A|_{W_{\mathbb R}^\perp},c)$ & $3/2$ & $\lambda|_{S_+\otimes W}\,c$ $(2;9)$\\
$(\lambda|_{S_+\otimes W},\lambda|_{S_+\otimes W})$ & $7/2$ & $\lambda|_{S_+\otimes W}$ $(1;3)$, $A|_V\,\lambda|_{S_+\otimes W}$ $(1;11)$, $A|_{W_{\mathbb R}^\perp}\,\lambda|_{S_-\otimes W^*}$ $(1;9)$, $c\,\lambda^*|_{S_-\otimes W}$ $(2;7)$, $A|_V^2\,\lambda|_{S_+\otimes W}$ $(2;27)$, $A|_V\,A|_{W_{\mathbb R}^\perp}\,\lambda|_{S_-\otimes W^*}$ $(2;42)$, $A|_V\,c\,\lambda^*|_{S_-\otimes W}$ $(2;30)$, $(A|_{W_{\mathbb R}^\perp})^2\,\lambda|_{S_+\otimes W}$ $(2;21)$, $A|_{W_{\mathbb R}^\perp}\,c\,\lambda^*|_{S_+\otimes W^*}$ $(3;9)$\\
$(\lambda|_{S_+\otimes W},\lambda|_{S_-\otimes W^*})$ & $7/2$ & $\lambda|_{S_-\otimes W^*}$ $(1;5)$, $A|_V\,\lambda|_{S_-\otimes W^*}$ $(1;18)$, $A|_{W_{\mathbb R}^\perp}\,\lambda|_{S_+\otimes W}$ $(1;17)$, $c\,\lambda^*|_{S_+\otimes W^*}$ $(2;7)$, $A|_V^2\,\lambda|_{S_-\otimes W^*}$ $(2;44)$, $A|_V\,A|_{W_{\mathbb R}^\perp}\,\lambda|_{S_+\otimes W}$ $(2;79)$, $A|_V\,c\,\lambda^*|_{S_+\otimes W^*}$ $(2;30)$, $(A|_{W_{\mathbb R}^\perp})^2\,\lambda|_{S_-\otimes W^*}$ $(2;40)$, $A|_{W_{\mathbb R}^\perp}\,c\,\lambda^*|_{S_-\otimes W}$ $(2;21)$\\
$(\lambda|_{S_+\otimes W},c)$ & $2$ & $c$ $(1;5)$, $A|_V\,c$ $(1;17)$, $A|_V^2\,c$ $(2;39)$, $(A|_{W_{\mathbb R}^\perp})^2\,c$ $(2;18)$\\
$(\lambda|_{S_-\otimes W^*},\lambda|_{S_-\otimes W^*})$ & $7/2$ & $\lambda|_{S_+\otimes W}$ $(1;2)$, $A|_V\,\lambda|_{S_+\otimes W}$ $(1;7)$, $A|_{W_{\mathbb R}^\perp}\,\lambda|_{S_-\otimes W^*}$ $(1;8)$, $A|_V^2\,\lambda|_{S_+\otimes W}$ $(2;17)$, $A|_V\,A|_{W_{\mathbb R}^\perp}\,\lambda|_{S_-\otimes W^*}$ $(2;37)$, $(A|_{W_{\mathbb R}^\perp})^2\,\lambda|_{S_+\otimes W}$ $(2;19)$, $A|_{W_{\mathbb R}^\perp}\,c\,\lambda^*|_{S_+\otimes W^*}$ $(2;12)$\\
$(\lambda|_{S_-\otimes W^*},c)$ & $2$ & $A|_{W_{\mathbb R}^\perp}\,c$ $(1;10)$, $A|_V\,A|_{W_{\mathbb R}^\perp}\,c$ $(2;44)$\\
$(c,c)$ & $1/2$ & $-$\\
\end{tabular}
\tabcaption{The leg sets of the defect cycles along $\xi\in S_+\otimes W$ through the pairs of field and ghost letters without derivatives, with the fewest vertices of $I$ and the number of chains; along $\xi\in S_-\otimes W^*$ the chiralities of $\lambda$ and $\lambda^*$ are exchanged. A leg set enters $\Upsilon_\xi$ when its output weight is at most the bound of the second column, and the empty leg set $\emptyset$ enters the poles.\label{tab:c11-defectpairs}}

\begin{tabular}{l|cccc|cccc|cccc}
pair & $\xi$, $n=0$ & $1$ & $2$ & $3$ & $\Upsilon^0_\eta$, $n=0$ & $1$ & $2$ & $3$ & $\Upsilon^0_{K_b}$, $n=0$ & $1$ & $2$ & $3$\\
\hline
$(A|_V,A|_V)$ & 4 & 9 & 20 & 43 & 1 & 4 & 9 & 20 & 2 & 4 & 8 & 19\\
$(A|_V,A|_{W_{\mathbb R}^\perp})$ & 4 & 9 & 20 & 43 & 1 & 4 & 9 & 20 & 1 & 2 & 7 & 16\\
$(A|_V,\lambda|_{S_+\otimes W})$ & 8 & 19 & 35 & 68 & 2 & 7 & 16 & 33 & 1 & 4 & 9 & 20\\
$(A|_V,\lambda|_{S_-\otimes W^*})$ & 7 & 16 & 33 & 61 & 4 & 8 & 19 & 35 & 1 & 4 & 9 & 20\\
$(A|_V,c)$ & 1 & 4 & 9 & 20 & 0 & 1 & 4 & 9 & 1 & 2 & 4 & 8\\
$(A|_V,A^*|_V)$ & 8 & 22 & 49 & 97 & 3 & 8 & 22 & 49 & 1 & 6 & 15 & 35\\
$(A|_V,A^*|_{W_{\mathbb R}^\perp})$ & 8 & 22 & 49 & 97 & 3 & 8 & 22 & 49 & 1 & 5 & 13 & 29\\
$(A|_V,\lambda^*|_{S_+\otimes W^*})$ & 3 & 11 & 26 & 62 & 1 & 6 & 15 & 35 & 1 & 3 & 8 & 21\\
$(A|_V,\lambda^*|_{S_-\otimes W})$ & 6 & 15 & 35 & 73 & 0 & 3 & 11 & 26 & 1 & 3 & 8 & 21\\
$(A|_V,c^*)$ & 3 & 13 & 34 & 80 & 0 & 3 & 13 & 34 & 1 & 3 & 10 & 25\\
$(A|_{W_{\mathbb R}^\perp},A|_{W_{\mathbb R}^\perp})$ & 4 & 9 & 19 & 39 & 1 & 4 & 9 & 19 & 2 & 4 & 7 & 14\\
$(A|_{W_{\mathbb R}^\perp},\lambda|_{S_+\otimes W})$ & 7 & 16 & 33 & 61 & 4 & 7 & 16 & 30 & 1 & 3 & 7 & 16\\
$(A|_{W_{\mathbb R}^\perp},\lambda|_{S_-\otimes W^*})$ & 7 & 16 & 30 & 60 & 2 & 7 & 16 & 33 & 1 & 3 & 7 & 16\\
$(A|_{W_{\mathbb R}^\perp},c)$ & 1 & 4 & 9 & 20 & 0 & 1 & 4 & 9 & 0 & 1 & 2 & 7\\
$(A|_{W_{\mathbb R}^\perp},A^*|_V)$ & 8 & 22 & 49 & 97 & 3 & 8 & 22 & 49 & 1 & 5 & 13 & 29\\
$(A|_{W_{\mathbb R}^\perp},A^*|_{W_{\mathbb R}^\perp})$ & 7 & 18 & 40 & 81 & 3 & 7 & 18 & 40 & 0 & 1 & 5 & 16\\
$(A|_{W_{\mathbb R}^\perp},\lambda^*|_{S_+\otimes W^*})$ & 3 & 10 & 26 & 59 & 1 & 5 & 14 & 31 & 0 & 1 & 4 & 14\\
$(A|_{W_{\mathbb R}^\perp},\lambda^*|_{S_-\otimes W})$ & 5 & 14 & 31 & 70 & 0 & 3 & 10 & 26 & 0 & 1 & 4 & 14\\
$(A|_{W_{\mathbb R}^\perp},c^*)$ & 3 & 13 & 34 & 80 & 0 & 3 & 13 & 34 & 0 & 1 & 4 & 19\\
$(\lambda|_{S_+\otimes W},\lambda|_{S_+\otimes W})$ & 9 & 20 & 42 & 79 & 3 & 7 & 16 & 35 & 2 & 6 & 12 & 24\\
$(\lambda|_{S_+\otimes W},\lambda|_{S_-\otimes W^*})$ & 9 & 20 & 43 & 82 & 4 & 9 & 20 & 43 & 4 & 7 & 15 & 27\\
$(\lambda|_{S_+\otimes W},c)$ & 4 & 8 & 19 & 35 & 1 & 2 & 7 & 16 & 0 & 1 & 4 & 9\\
$(\lambda|_{S_+\otimes W},A^*|_V)$ & 15 & 35 & 73 & 143 & 5 & 14 & 31 & 70 & 3 & 8 & 21 & 46\\
$(\lambda|_{S_+\otimes W},A^*|_{W_{\mathbb R}^\perp})$ & 14 & 31 & 70 & 135 & 3 & 10 & 26 & 59 & 1 & 4 & 14 & 35\\
$(\lambda|_{S_+\otimes W},\lambda^*|_{S_+\otimes W^*})$ & 8 & 21 & 47 & 88 & 3 & 8 & 21 & 47 & 0 & 2 & 7 & 19\\
$(\lambda|_{S_+\otimes W},\lambda^*|_{S_-\otimes W})$ & 8 & 21 & 46 & 86 & 1 & 3 & 11 & 29 & 0 & 1 & 5 & 14\\
$(\lambda|_{S_+\otimes W},c^*)$ & 10 & 25 & 62 & 116 & 2 & 5 & 19 & 49 & 0 & 2 & 10 & 29\\
$(\lambda|_{S_-\otimes W^*},\lambda|_{S_-\otimes W^*})$ & 7 & 16 & 35 & 68 & 4 & 9 & 20 & 42 & 2 & 6 & 12 & 24\\
$(\lambda|_{S_-\otimes W^*},c)$ & 2 & 7 & 16 & 33 & 2 & 4 & 8 & 19 & 0 & 1 & 4 & 9\\
$(\lambda|_{S_-\otimes W^*},A^*|_V)$ & 14 & 31 & 70 & 135 & 6 & 15 & 35 & 73 & 3 & 8 & 21 & 46\\
$(\lambda|_{S_-\otimes W^*},A^*|_{W_{\mathbb R}^\perp})$ & 10 & 26 & 59 & 123 & 5 & 14 & 31 & 70 & 1 & 4 & 14 & 35\\
$(\lambda|_{S_-\otimes W^*},\lambda^*|_{S_+\otimes W^*})$ & 3 & 11 & 29 & 60 & 3 & 8 & 21 & 46 & 0 & 1 & 5 & 14\\
$(\lambda|_{S_-\otimes W^*},\lambda^*|_{S_-\otimes W})$ & 8 & 21 & 47 & 88 & 3 & 8 & 21 & 47 & 0 & 2 & 7 & 19\\
$(\lambda|_{S_-\otimes W^*},c^*)$ & 5 & 19 & 49 & 101 & 3 & 10 & 25 & 62 & 0 & 2 & 10 & 29\\
$(c,c)$ & 0 & 1 & 4 & 9 & 0 & 0 & 1 & 4 & 0 & 1 & 2 & 4\\
$(c,A^*|_V)$ & 9 & 20 & 43 & 82 & 4 & 9 & 20 & 43 & 4 & 8 & 19 & 35\\
$(c,A^*|_{W_{\mathbb R}^\perp})$ & 9 & 20 & 43 & 82 & 4 & 9 & 20 & 43 & 2 & 7 & 16 & 32\\
$(c,\lambda^*|_{S_+\otimes W^*})$ & 6 & 14 & 30 & 53 & 4 & 8 & 19 & 35 & 1 & 4 & 9 & 20\\
$(c,\lambda^*|_{S_-\otimes W})$ & 8 & 19 & 35 & 64 & 2 & 6 & 14 & 30 & 1 & 4 & 9 & 20\\
$(c,c^*)$ & 8 & 22 & 49 & 91 & 3 & 8 & 22 & 49 & 1 & 6 & 15 & 35\\
$(A^*|_V,A^*|_V)$ & 13 & 34 & 80 & 169 & 3 & 13 & 34 & 80 & 3 & 10 & 25 & 61\\
$(A^*|_V,A^*|_{W_{\mathbb R}^\perp})$ & 13 & 34 & 80 & 169 & 3 & 13 & 34 & 80 & 1 & 4 & 19 & 49\\
$(A^*|_V,\lambda^*|_{S_+\otimes W^*})$ & 5 & 19 & 49 & 101 & 3 & 10 & 25 & 62 & 0 & 2 & 10 & 29\\
$(A^*|_V,\lambda^*|_{S_-\otimes W})$ & 10 & 25 & 62 & 116 & 2 & 5 & 19 & 49 & 0 & 2 & 10 & 29\\
$(A^*|_V,c^*)$ & 4 & 15 & 46 & 101 & 1 & 4 & 15 & 46 & 0 & 1 & 7 & 24\\
\end{tabular}
\tabcaption{The number of leg sets of the defect cycles entering $\Upsilon_\xi$, $\xi\in S_+\otimes W$, the translation-invariant part $\Upsilon^0_\eta$, $\eta\in S_+\otimes W^*$, and the translation-invariant part $\Upsilon^0_{K_b}$, for $n$ derivatives on the two letters and up to six vertices on the cycle, on the $45$ pairs of slots with defect cycles; the other ten pairs, of two letters among $A^*|_{W_{\mathbb R}^\perp}$, $\lambda^*|_{S_+\otimes W^*}$, $\lambda^*|_{S_-\otimes W}$, $c^*$, are those of Proposition \ref{prop:c11-oddcensus} (ii).\label{tab:c11-defectcounts}}

\begin{proposition}[The quadratic free directions along $\Sigma$ and $\Sigma^*$]
\label{prop:c11-freexi}
\begin{itemize}
\item (i) The quadratic local functionals of degree $-1$, of Weyl weight $\tfrac92$ in the fields and derivatives, linear in $\xi\in S_+\otimes W$, translation invariant and covariant under $\mathrm{Spin}(4)\times\mathrm{Spin}(6)\times\mathfrak g$, are spanned by $T_1$, $T_2$, $T_3$, $T_4$ and $\int\big((d^*A|_V)\xi,\lambda^*\big)$; their $Q$-closed combinations are spanned by $S^V_2(\xi)$ and $S^\perp_2(\xi)$, and the $Q$-exact ones are $0$.
\item (ii) The arity-one equation along $\xi$, the arity-one part along $\xi$ of $s(\hbar F_1)=\hbar\,\mathrm{Ob}$ (Theorem \ref{thm:c8-obstruction} (ii)), is
\begin{equation}
(Q+\{I,-\})\,\hbar F_1^{(1)}(\xi)+\big\{S^{(1)}(\xi),\hbar F_1^{(0)}\big\}=\hbar\,\mathrm{Ob}^{(1)}(\xi).
\label{eq:c11-arity1xi}
\end{equation}
At fixed $\hbar F_1^{(0)}$ two solutions differ, in the quadratic sector, by $\theta^\xi_V$ units times $S^V_2(\xi)$ plus $\theta^\xi_\perp$ units times $S^\perp_2(\xi)$, and the cubic part of \eqref{eq:c11-arity1xi} then contains $\theta^\xi_V\{I_3,S^V_2(\xi)\}+\theta^\xi_\perp\{I_3,S^\perp_2(\xi)\}$ units, $I_3$ the cubic part of $I$. For $\theta^\xi_V=\theta^\xi_\perp=:\theta^\xi$ the direction extends to $\theta^\xi$ units times $S^{(1)}(\xi)$, the arity-one value of $G_B$ (Definition \ref{def:c8-B}) and of the field redefinition by the R-rotation $S^{(1)}(R_\xi)$, $[R_\xi,\xi]=\xi$ (Theorem \ref{thm:c8-twist}). The cubic part fixes $\theta^\xi_V-\theta^\xi_\perp$ when the class of $\{I_3,S^V_2(\xi)\}$ modulo the images under $Q$ of the cubic local functionals of degree $-1$ is nonzero, as for the sector coefficients along $\Delta$ (Proposition \ref{prop:c11-arity1} (iii)).
\item (iii) The one-body kernel of a letter $X$ along $\xi$ is $-\{\hbar F_1^{(1)}(\xi),X\}$; the free directions add $-\theta^\xi_V$ units times the arrows of $S^V_2(\xi)$ and $-\theta^\xi_\perp$ units times those of $S^\perp_2(\xi)$, which are $\{S^V_2(\xi),A|_V\}=\Gamma(\xi,\lambda)|_V$ and $\{S^\perp_2(\xi),A|_{W_{\mathbb R}^\perp}\}=\Gamma(\xi,\lambda)|_{W_{\mathbb R}^\perp}$ on the gauge field and the scalars and $\tfrac12\rho(dA|_V)\xi$ and $\tfrac12\rho(dA|_{W_{\mathbb R}^\perp})\xi$ on $\lambda$ (Lemma \ref{lem:c11-lines} (i)).
\item (iv) Along $\Sigma^*$, with the position included at weight $-1$ and $\eta$ at weight $\tfrac12$, the $Q$-closed quadratic directions of degree $-1$ and Weyl weight $4$, linear in $\eta$, are spanned by $\sum_\mu\int x^\mu S^V_2(\rho(e_\mu)\eta)$ and $\sum_\mu\int x^\mu S^\perp_2(\rho(e_\mu)\eta)+\int s_0$, whose sum is the quadratic part of $S^{(1)}(\eta)$.
\end{itemize}
The same holds along $S_-\otimes W^*$ with the chiralities exchanged.
\end{proposition}

\begin{proof}
(i) The field contents of degree $-1$ and field weight $\tfrac92$ linear in a constant spinor are $(\lambda,A^*)$ without derivatives and $(A,\lambda^*)$ with one derivative, since $c$ and $c^*$ would need a further weight $\tfrac12$; the covariant pairings are $\Gamma(\xi,\lambda)$ with $A^*$ in its parts along $V$ and $W_{\mathbb R}^\perp$, and $\rho(\partial A)\xi$ with $\lambda^*$, $\partial A$ being $dA|_V$, $d^*A|_V$ or the mixed part $dA|_{W_{\mathbb R}^\perp}$, the volume element acting by a sign on each chirality. On functionals $Q$ substitutes $A\mapsto-dc$, $A^*\mapsto-\tfrac12d^*dA$, $\lambda^*\mapsto{\not d}\lambda$, $c^*\mapsto d^*(A^*|_V)$ and annihilates $c$ and $\lambda$ (Proposition \ref{prop:c1-Q}). $Q\int((d^*A|_V)\xi,\lambda^*)$ contains $-\int((d^*dc)\xi,\lambda^*)$, produced by no other structure, so its coefficient vanishes in a closed combination; $Q$ preserves the split into the parts along $V$ and along $W_{\mathbb R}^\perp$, and $S^V_2(\xi)$, $S^\perp_2(\xi)$ are closed, being the two parts of the quadratic part of $\{S,S^{(1)}(\xi)\}=0$. A $Q$-exact functional is $Q$ of a functional of degree $-2$ and field weight $\tfrac92$, and two antifields have weight at least $5$, while $c^*$ needs a further field of weight $\tfrac12$. Both evaluations of Remark \ref{rem:c11-oddmethod} give the same ranks on plane-wave normal forms. (ii) Theorem \ref{thm:c8-obstruction} (iii) restricted to arity one along $\xi$ at fixed arity-zero part. The cubic part of $s\big(\theta^\xi_VS^V_2+\theta^\xi_\perp S^\perp_2\big)$ is the bracket with $I_3$, and $\{I_3,S^V_2(\xi)+S^\perp_2(\xi)\}=-QS_3(\xi)$ with $S_3(\xi)$ the cubic part of $S^{(1)}(\xi)$, the cubic part of $\{S,S^{(1)}(\xi)\}=0$; $G_B^{(1)}(\xi)=S^{(1)}(B\xi)=S^{(1)}(\xi)$ for $\xi\in S_+\otimes W$. (iii) Theorem \ref{thm:c11-classes} (iii) and Lemma \ref{lem:c11-lines} (i). (iv) The same enumeration with $\int x^\mu$ times the structures of (i) at $\rho(e_\mu)\eta$, $\int s_0$ and $\int(\rho(A|_V)\eta,\lambda^*)$, and the exact ranks of $Q$ on plane-wave normal forms; $s_0$ enters the closed direction together with the part along $W_{\mathbb R}^\perp$.
\end{proof}

\begin{proposition}[Two-leg obstructions and solutions along the odd generators, the rotations and $K_b$]
\label{prop:c11-oddsol}
Let $\hbar F_1^{(0)}$ contain the finite counterterm $\tfrac{27}{16}$ units times $\int(d^*A|_V,d^*A|_V)$ of Theorem \ref{thm:c4-finite}, write $L=\log\tfrac32$, and, for $\xi\in S_+\otimes W$, $\Gamma_V:=\Gamma(\xi,\lambda)|_V$, $\Gamma_\perp:=\Gamma(\xi,\lambda)|_{W_{\mathbb R}^\perp}$,
\[
g_1=\int(\Gamma_V,d^*d\,A|_V),\qquad g_2=\int(\Gamma_V,dd^*A|_V),\qquad g_3=\int(\Gamma_\perp,d^*d\,A|_{W_{\mathbb R}^\perp}),\qquad g_4=\int\big(d^*dc,(\xi,\lambda^*)\big),\qquad T_0=\int\big(d^*A|_V,(\xi,\lambda^*)\big).
\]
The obstructions are the sums of the coefficients of $\varepsilon^0$ of the defect cycles with two legs along $x$, as in Proposition \ref{prop:c9-structure} (ii), with the placements of Lemma \ref{lem:c11-lines}.
\begin{itemize}
\item (i) Along $\xi$, for both chiralities: $QT_1=-QT_3=-\tfrac12g_1$, $QT_2=-QT_4=-\tfrac12g_3$, $QT_0=-(g_2+g_4)$; the cohomology of $Q$ on the quadratic functionals of degree $0$, Weyl weight $\tfrac92$, linear in $\xi$ and covariant is zero, and
\[
\hbar\,\mathrm{Ob}^{(1)}(\xi)=\text{units}\cdot\Big[\big(\tfrac{79}{24}-\tfrac12L\big)(g_1+g_3)+\tfrac{27}8\,g_2+0\cdot g_4\Big].
\]
The arity-one equation \eqref{eq:c11-arity1xi} is solvable on two legs, with the solutions
\[
\hbar F_1^{(1)}(\xi)=\big(L-\tfrac{79}{12}\big)\,\text{units}\cdot(T_1+T_2)+\theta^\xi_V\,\text{units}\cdot S^V_2(\xi)+\theta^\xi_\perp\,\text{units}\cdot S^\perp_2(\xi),
\]
the part $\tfrac{27}8g_2$ being the bracket $\{\hbar F_1^{(0)},S^{(1)}(\xi)\}$; the field-counting directions of $\hbar F_1^{(0)}$ (Theorem \ref{thm:c9-F3}) shift this by brackets with $S^{(1)}(\xi)$ (Proposition \ref{prop:c11-freedom} (i)). On two legs the coefficient of $\log\varepsilon$ in $I^{CT}_{\mathfrak{psl}(4|4)}[\varepsilon]^{(1)}(\xi)$ is $\{\mathcal N,S^{(1)}(\xi)\}$ units, $\mathcal N=2\int(A,A^*)-4\int(\lambda^*,\lambda)-\int(c^*,c)$ the field-counting functional: $2(T_1+T_2-T_3-T_4)$ units.
\item (ii) Along $\eta\in\Sigma^*$, for both chiralities, with position degree at most one, writing $u_0:=\int(\rho(A|_V)\eta,\lambda^*)$ and $s_0$ of \eqref{eq:c11-Seta},
\[
\hbar\,\mathrm{Ob}^{(1)}(\eta)=\sum_\mu x^\mu\,\hbar\,\mathrm{Ob}^{(1)}(\rho(e_\mu)\eta)+\text{units}\cdot\Big[-\tfrac52\int(\eta,\lambda)\,d^*A|_V-\tfrac52\int(\eta,\rho(dA|_V)\lambda)+\big(3L-\tfrac32\big)\int(\eta,\rho(dA|_{W_{\mathbb R}^\perp})\lambda)+\tfrac52\int(\rho(dc)\eta,\lambda^*)\Big]+Q(\text{local}),
\]
the cohomology of $Q$ at this weight is one-dimensional and the residual has no component along it, and the solutions are
\[
\hbar F_1^{(1)}(\eta)=\sum_\mu x^\mu\big(L-\tfrac{79}{12}\big)\,\text{units}\cdot(T_1+T_2)(\rho(e_\mu)\eta)-\big(3L-\tfrac32\big)\,\text{units}\cdot s_0-\tfrac52\,\text{units}\cdot u_0+\theta^\eta_V\,\text{units}\cdot S^V(\eta)+\theta^\eta_\perp\,\text{units}\cdot S^\perp(\eta),
\]
with $S^V(\eta)=\sum_\mu\int x^\mu S^V_2(\rho(e_\mu)\eta)$ and $S^\perp(\eta)=\sum_\mu\int x^\mu S^\perp_2(\rho(e_\mu)\eta)+\int s_0$ (Proposition \ref{prop:c11-freexi} (iv)).
\item (iii) Along $M_\omega\in\mathfrak{so}(V)$, with $\rho_\mu=\rho(e_\mu)$ and $S^{(1)}(M_\omega)_{\rm spin}:=\int\omega_{\mu\nu}(A_\mu,A^*_\nu)+\int\big(\tfrac14\omega_{\mu\nu}\rho_\mu\rho_\nu\lambda,\lambda^*\big)$: $\hbar\,\mathrm{Ob}^{(1)}(M_\omega)=\tfrac13$ units times $Q\,S^{(1)}(M_\omega)_{\rm spin}$, $0$ on $(A|_{W_{\mathbb R}^\perp},A|_{W_{\mathbb R}^\perp})$, and $\hbar F_1^{(1)}(M_\omega)=\tfrac13$ units times $S^{(1)}(M_\omega)_{\rm spin}$ plus the three sector parts of $S^{(1)}(M_\omega)$ with coefficients $\theta^M_\perp,\theta^M_\lambda,\theta^M_V$ units.
\item (iv) Along $R\in\mathfrak{so}(W_{\mathbb R}^\perp)$, $I^{CT}_{\mathfrak{psl}(4|4)}[\varepsilon]^{(1)}(R)=0$ and $\mathrm{Ob}^{(1)}(R)=0$: $R.$ commutes with $Q$, $Q^{GF}$ and $H$, preserves the pairing, annihilates $I$, $I^{CT}[\varepsilon]$ and $\hbar F_1^{(0)}$, and $S^{(1)}(R)$ forms no line; the solutions are $\theta^R_\perp$, $\theta^R_\lambda$ units times the sector parts of $S^{(1)}(R)$.
\item (v) Along $K_b$, at a finite endpoint $T$ and with position degree at most two,
\[
\begin{aligned}
\hbar F_1^{(1)}(K_b)=\ &\text{units}\int(b,x)\Big\{-\tfrac{22}3(A,A^*)+16(\lambda,\lambda^*)+\tfrac83(c,c^*)\Big\}+\tfrac23\,\text{units}\int\big\{(\iota_bA,\iota_xA^*)-(\iota_xA,\iota_bA^*)\big\}\\
&-\tfrac13\,\text{units}\int b_\mu x_\nu\big(\lambda,\rho_\mu\rho_\nu\lambda^*\big)+\text{units}\cdot\sum_d\theta^K_d\,S^{(1)}(K_b)_d+\sum_{j=1}^4\kappa^K_j\,G^K_j,
\end{aligned}
\]
with the sector parts $S^{(1)}(K_b)_d$ and four $Q$-closed quadratic functionals $G^K_1,\dots,G^K_4$ of degree $-1$, the second-order symmetries of the free fields built from the conformal Killing tensors, which are neither exact nor of the form $S^{(1)}(z)$ on two legs; the cohomology of $Q$ is four-dimensional in degree $0$ and seven-dimensional in degree $-1$, and the residual has no component in degree $0$. The part linear in $x$ is $-(b,x)$ times the translation-invariant solution along $\Delta$ plus the spin term of the rotation part of the derivative of the vector field, $16-\tfrac13=\tfrac{47}3$ on $\lambda$.
\end{itemize}
\end{proposition}

\begin{proof}
(i) The census of Proposition \ref{prop:c11-freexi} (i) and its degree-$0$ analogue: $\xi\otimes\lambda|_{S_-\otimes W^*}$ contains the $SU(4)$-singlet in $V\otimes\mathbf1$ alone, so every bilinear in $(\xi,\lambda|_{S_-\otimes W^*})$ is a component of $\Gamma_V$, paired with $A|_V$ and two derivatives in $d^*d$ and $dd^*$; $\xi\otimes\lambda|_{S_+\otimes W}$ pairs with $A|_{W_{\mathbb R}^\perp}$ through $\Gamma_\perp$; $(\xi,\lambda^*)$ pairs with $c$ and two derivatives. $Q$ substitutes $A^*\mapsto-\tfrac12d^*dA$, $\lambda^*\mapsto{\not d}\lambda$, $A|_V\mapsto-dc$ (Proposition \ref{prop:c1-Q}), with $(\xi,{\not d}\lambda)=-d^*\Gamma_V$; $Qg_1=Qg_3=0$ and $Qg_2$, $Qg_4$ are proportional and nonzero, so the closed functionals of degree $0$ are spanned by $g_1$, $g_3$, $g_2+g_4$, which is the image of $Q$. The obstruction is the sum of the coefficients of $\varepsilon^0$ of the defect cycles with two legs, evaluated in two independent ways; its components along $g_2$, $g_4$ are fixed by the Wess--Zumino identity of Lemma \ref{lem:c9-sform}, $Q\,\mathrm{Ob}^{(1)}(\xi)=-\{\mathrm{Ob}^{(0)},S^{(1)}(\xi)\}$ with the two-leg arity-zero obstruction $-\tfrac{27}8$ units times $\int(d^*dc,d^*A|_V)$ (Theorem \ref{thm:c4-obstruction}). The bracket $\{\int(d^*A|_V,d^*A|_V),T_1\}$ is $2g_2$ up to the sign fixed by the Wess--Zumino identity, so the residual $\mathrm{Ob}^{(1)}(\xi)-\{F_1^{(0)},S^{(1)}(\xi)\}$ is $(\tfrac{79}{24}-\tfrac12L)(g_1+g_3)$ units, $Q$ of the stated particular solution; the free directions are Proposition \ref{prop:c11-freexi} (i). The coefficient of $\log\varepsilon$ is the singular part of the composite and cubic cycles along $\xi$; its difference with $\{\mathcal N,S^{(1)}(\xi)\}$ units is $Q$-closed on two legs, the coefficient of $\log\varepsilon$ of the counterterm cochain being $s$-closed and its arity-zero part $Q$ of $\mathcal N$ units (Theorem \ref{thm:c4-log}, Lemma \ref{lem:c9-sform}), hence a combination of $S^V_2(\xi)$ and $S^\perp_2(\xi)$, whose two coefficients are computed to vanish; and $\{\mathcal N,-\}$ multiplies each arrow $f\to f'$ of $S^{(1)}(\xi)$ by the difference of the weights $2$, $4$, $-2$, $-4$ of $A$, $\lambda$, $A^*$, $\lambda^*$. (ii) The position part along $b$ of every defect cycle along $\eta$ is the cycle along $\rho(b)\eta$ (Theorem \ref{thm:c11-oddsimplex} (ii)); $Q$ is block triangular for the split into position parts and translation-invariant parts, and the closed directions are those of Proposition \ref{prop:c11-freexi} (iv); the intrinsic residual is computed on a complete basis of the structures $(\eta,\lambda)\,d^*A|_V$, $(\eta,\rho(dA|_V)\lambda)$, $(\eta,\rho(dA|_{W_{\mathbb R}^\perp})\lambda)$, $(\rho(dc)\eta,\lambda^*)$ and their images, in both evaluations. (iii) $M_\omega.$ commutes with $Q$, $Q^{GF}$ and $H$; in the leg sets $(A|_V,A|_V)$, $(c,A^*|_V)$ and $(\lambda,\lambda)$ the closed functionals of degree $0$ linear in $\omega$ are spanned by $Q\int\omega_{\mu\nu}(A_\mu,A^*_\nu)$ and $Q\int(\tfrac14\omega_{\mu\nu}\rho_\mu\rho_\nu\lambda,\lambda^*)$, the structure $\int(\lambda,\omega_{\mu\nu}\rho_{[\mu}\rho_\nu\rho_{\rho]}\partial_\rho\lambda)$ being a total derivative, so the cohomology vanishes there, and the defect cycles give the multiple $\tfrac13$ in both sectors. (iv) The graphs with one vertex $S^{(1)}(R)$ are $R.$ applied to the legs of the flow, and the singular part of the one-loop part is removed by $I^{CT}[\varepsilon]$; $[Q,D_R]$ is the derivation by $(e^{-\varepsilon H}-e^{-TH})R.$, whose action tends to $R.$ of the renormalised flow, zero on the $R$-invariant action. (v) The sector bound of Proposition \ref{prop:c10-bound1} with $\lambda_{K_b}=2$ and position degree two gives the leg sets $(c,A^*|_V)$, $(A|_V,A|_V)$, $(A|_{W_{\mathbb R}^\perp},A|_{W_{\mathbb R}^\perp})$ and $(\lambda,\lambda)$; the cohomology is computed on complete bases in both evaluations, and the defect cycles are evaluated at a finite $T$, their coefficients of $\varepsilon^0$ being independent of $T$ (Theorem \ref{thm:c11-oddsimplex} (iii)). The part linear in $x$ is the position decomposition \eqref{eq:c11-Kdec} applied to the solution.
\end{proof}

\begin{theorem}[The arity-two equations on two legs and the uniform coefficients]
\label{thm:c11-oddar2}
With the solutions of Theorem \ref{thm:c9-quadratic} (iii), Theorem \ref{thm:c9-P1} (i) and Proposition \ref{prop:c11-oddsol}, write $\theta^\Delta_d$ for $\theta,\theta_\lambda,\theta_V$, $\theta^P_d$ for $\theta_P,\theta_{P,\lambda},\theta_{P,V}$, and $\theta^{\xi_\pm}_d$, $\theta^{\eta_\pm}_d$ for the coefficients of Proposition \ref{prop:c11-oddsol} along $\xi_\pm\in S_+\otimes W,\ S_-\otimes W^*$ and $\eta_+\in S_+\otimes W^*$, $\eta_-\in S_-\otimes W$, with $d\in\{V,\perp\}$ along the odd generators and $L=\log\tfrac32$. The classes of the arity-two residuals of Lemma \ref{lem:c9-sform} on two legs, the obstructions being the sums of the coefficients of $\varepsilon^0$ of the cycles with one special edge and two vertices $S^{(1)}(x)$, $S^{(1)}(y)$ or one vertex $S^{(2)}(x,y)$, vanish exactly under the conditions of Table \ref{tab:c11-oddar2}.

\begin{tabular}{l|l}
pair & condition \\
\hline
$(\Delta,P_b),\ (\Delta,\xi_\pm)$ & $\theta^\Delta_\perp=\theta^\Delta_\lambda=\theta^\Delta_V=-\tfrac13$ \\
$(\xi_+,\xi_-)$ & $\theta^P_d=\theta^{\xi_+}_V+\theta^{\xi_-}_V-\tfrac{27}4-2L\ \text{for every }d,\ \theta^{\xi_+}_V+\theta^{\xi_-}_V=\theta^{\xi_+}_\perp+\theta^{\xi_-}_\perp$ \\
$(\xi_\pm,\xi'_\pm)$ & $\theta^{\xi_\pm}_V=\theta^{\xi_\pm}_\perp$ \\
$(\xi_\pm,P_b)$ & $\theta^P_V=\theta^P_\perp=\theta^P_\lambda$ \\
$(\xi_+,\eta_-)$ & $\theta^{\xi_+}_V-\theta^{\xi_+}_\perp=\theta^{\eta_-}_V-\theta^{\eta_-}_\perp$ \\
$(\xi_+,\eta_+)$ & $\theta^R_\perp=\theta^R_\lambda=:\tau,\ \theta^\Delta_d=\theta^M_d=\tau-\tfrac13,\ \theta^{\xi_+}_d+\theta^{\eta_+}_d=\tau+\tfrac{77}{12}+2L\ (d=V,\perp)$ \\
\end{tabular}
\tabcaption{The conditions under which the classes of the two-leg arity-two residuals vanish, $L=\log\tfrac32$.\label{tab:c11-oddar2}}

On $(\xi_\pm,P_b)$ the obstruction is $\tfrac14\log\tfrac43$ times a $Q$-exact functional, and the arity-two generating cochain of two odd generators is $I+cS^{(1)}(\xi)+c'S^{(1)}(\xi')-cc'S^{(2)}(\xi,\xi')$ for odd parameters $c$, $c'$. Consequently:
\begin{itemize}
\item (i) every sector coefficient along $P_b$, $\Sigma$, $\Sigma^*$, $\mathfrak{so}(V)$ and $\mathfrak{so}(W_{\mathbb R}^\perp)$ is uniform, $\theta^z_d=:t_z$, and
\begin{equation}
t_\Delta=-\tfrac13,\qquad t_M=t_\Delta,\qquad t_R=t_\Delta+\tfrac13=0,\qquad t_{\xi_+}+t_{\xi_-}-t_P=\tfrac{27}4+2\log\tfrac32,\qquad t_{\xi_+}+t_{\eta_+}-t_\Delta=\tfrac{27}4+2\log\tfrac32;
\label{eq:c11-trel}
\end{equation}
\item (ii) the solutions form an affine space over the inner direction $t_z\mapsto t_z+c\lambda_z$, $\lambda_z=-1,+1,-2,+2,0$ on $\Sigma$, $\Sigma^*$, $P_b$, $K_b$ and the even rotations and $\Delta$, generated by the field redefinition $\phi^{(0)}=c\,\text{units}\cdot S^{(1)}(\Delta)$, and the outer direction $t_{\xi_\pm}\mapsto t_{\xi_\pm}\pm a$, $t_{\eta_\pm}\mapsto t_{\eta_\pm}\mp a$, the direction $G_B$ of Definition \ref{def:c8-B}; the relations \eqref{eq:c11-trel} and the sector uniformity are invariant under both, and they are the scheme-independent content of the computed pairs on the action side;
\item (iii) in the translation-covariant scheme of Definition \ref{def:c11-free}, $t_P=t_\Delta=-\tfrac13$, hence $t_{\xi_+}+t_{\xi_-}=\tfrac{77}{12}+2\log\tfrac32$ and $t_{\eta_+}=t_{\xi_-}$;
\item (iv) the individual $t_{\xi_\pm}$, $t_{\eta_\pm}$, $t_P$, the coefficients $\theta^K_d$, $\kappa^K_j$, $\kappa_V$, $\kappa_\lambda$ and every particular part with $\log\tfrac32$ are scheme artefacts.
\end{itemize}
The relations \eqref{eq:c11-trel} are the one-loop parts of $[\Delta,P_b]=-2P_b$, $[\Delta,\xi]=-\xi$, $[\xi,P_b]=0$, $[\xi_+,\xi_-]=P_{\Gamma(\xi_+,\xi_-)|_V}$ and of the bracket $[\xi,\eta]$ with its components along $\Delta$, $\mathfrak{so}(V)$ and $\mathfrak{so}(W_{\mathbb R}^\perp)$ (Proposition \ref{prop:c2-special}): with the corrected generators $S^{(1)}(z)+\hbar t_z\,\text{units}\cdot S^{(1)}(z)$ the arity-two equation along $(x,y)$ states that the order-$\hbar$ part of $(-1)^{|x|}\{S^{(1)}(x),S^{(1)}(y)\}+S^{(1)}([x,y])$ closes up to $\Lambda^{(0)}$-exact terms once the computed class of the obstruction is included.
\end{theorem}

\begin{proof}
In \eqref{eq:c9-arity2} the free directions enter through $F_1^{(1)}([x,y])$ and the brackets $\{F_1^{(1)}(x),S^{(1)}(y)\}$, $\{S^{(1)}(x),F_1^{(1)}(y)\}$; for uniform coefficients, $(-1)^{|x|}\{S^{(1)}(x),S^{(1)}(y)\}=-S^{(1)}([x,y])-\Lambda^{(0)}S^{(2)}(x,y)$ (Lemma \ref{lem:c2-components}) gives the change $(t_x+t_y-t_{[x,y]})$ units times $S^{(1)}([x,y])$ modulo exact terms, $t_{[x,y]}S^{(1)}([x,y])$ meaning the sum over the components of $[x,y]$ in the blocks. On two legs the sector parts enter separately through the brackets $\{S^V_2(\xi),S_2(\xi')\}$, $\{S_2(\xi),S^V_2(\xi')\}$ and their analogues, whose classes in the cohomology of $Q$ are combinations of the classes of the sector parts of $S^{(1)}([x,y])$ by the Fierz identities of the arrows $\lambda\to A|_V\to\lambda$ and $\lambda\to A|_{W_{\mathbb R}^\perp}\to\lambda$. The obstructions are the coefficients of $\varepsilon^0$ of the cycles of the statement, evaluated in two independent ways, and their classes are computed on complete bases of the two-leg functionals with the quantum numbers of the pair, the generating cochain of two odd generators having the sign forced by the classical master equation, $QS^{(2)}(\xi,\xi')+S^{(1)}(P_{\Gamma(\xi,\xi')|_V})-\{S^{(1)}(\xi),S^{(1)}(\xi')\}=0$ on two legs; the vanishing of the residual classes gives the table, the rows along $\Delta$ being Theorem \ref{thm:c9-theta2}. (i) $(\xi_\pm,\xi'_\pm)$ and $(\xi_\pm,P_b)$ give the uniformity along $\Sigma$ and $P_b$, then $(\xi_+,\eta_-)$ the uniformity along $\Sigma^*$, and $(\xi_+,\eta_+)$ with $\theta^\Delta_d=-\tfrac13$ gives $\tau=0$, $t_M=t_\Delta$, $t_R=0$ and $t_{\xi_+}+t_{\eta_+}=\tfrac{77}{12}+2L=t_\Delta+\tfrac{27}4+2L$; $(\xi_+,\xi_-)$ gives the fourth relation. (ii) The field redefinition generated by $\phi^{(0)}$ changes $F_1^{(1)}(x)$ by $\{\phi^{(0)},S^{(1)}(x)\}=c\lambda_x$ units times $S^{(1)}(x)$ modulo exact terms (Lemma \ref{lem:c2-components} with $S^{(2)}(\Delta,x)=0$), and $G_B$ is an $s$-closed cochain whose arity-one value along $\xi_\pm$ is $\pm S^{(1)}(\xi_\pm)$ and along $\eta_\pm$ is $\mp S^{(1)}(\eta_\pm)$ (Theorem \ref{thm:c8-obstruction} (iv)); both preserve the solvability of every equation, $\lambda$ and $B$ being derivations of $\mathfrak{psl}(4|4)$. The left sides of \eqref{eq:c11-trel} are invariant: the inner direction shifts $t_{\xi_\pm}$, $t_{\eta_\pm}$, $t_P$ by $-c$, $+c$, $-2c$ and leaves $t_\Delta$, $t_M$, $t_R$ fixed, and the outer direction shifts $t_{\xi_\pm}$ by $\pm a$ and $t_{\eta_\pm}$ by $\mp a$. (iii) The position part of $\hbar F_1^{(1)}(\Delta)$ along $b$ is $2\hbar F_1^{(1)}(P_b)$, so $t_P=t_\Delta$, and \eqref{eq:c11-trel}. (iv) Under the inner and outer directions the individual coefficients move.
\end{proof}

\begin{proposition}[Renormalisation freedom of $Q_1$, $S_1$ and $K_1$]
\label{prop:c11-freedom}
Under the changes of Proposition \ref{prop:c8-laws} the components change as follows; $\{S^{(1)}(x),-\}$ is the classical transformation along $x$.
\begin{itemize}
\item (i) A field redefinition $F_{\mathfrak{psl}(4|4)}\mapsto F_{\mathfrak{psl}(4|4)}+\hbar s\phi$ changes the component along $x$ by $-\{(s\phi)^{(1)}(x),-\}$, a derivation of the letters: it changes the one-body kernels and leaves the two-body kernels unchanged. For the field countings $\phi^{(0)}=\sum_fy_f\int(f,f^*)$ the change of $Q_1$ multiplies each arrow $f\to f'$ of the classical transformation by $y_{f'}-y_f$; for $\phi^{(0)}=t\,S^{(1)}(\Delta)$ it is a multiple of the classical transformation proportional to $\lambda_x$.
\item (ii) The coupling direction leaves $Q_1$, $S_1$, $K_1$ unchanged. The outer direction $G_B$ changes $Q_1$ along $S_+\otimes W$ by $-a$ and along $S_-\otimes W^*$ by $+a$ times the classical transformation, $S_1$ along $S_+\otimes W^*$ by $+a$ and along $S_-\otimes W$ by $-a$ times it, and leaves $K_1$, $D_1$, $P_1$, $M_1$, $R_1$ unchanged.
\item (iii) The free directions of $\hbar F_1^{(1)}(x)$ change the one-body kernels: along $\Sigma$ by Proposition \ref{prop:c11-freexi} (iii); along $\Delta$ and $P_b$ through the sector coefficients of Definition \ref{def:c11-free}, which enter $D_1$, $P_1$ and, through $P_1$, the relations of Proposition \ref{prop:c11-S} and Corollary \ref{cor:c11-K}; along $K_b$ through the $s$-closed directions of degree $-1$ quadratic in the Killing data. $\hbar F_1^{(0)}$ enters the one-body kernels through \eqref{eq:c11-arity1xi} and through the trees with $S^{(1)}(x)$ (Remark \ref{rem:c10-scheme1} (a), (b)).
\item (iv) A finite part $\psi'$ of the operators changes the component along $x$ by $[\{S^{(1)}(x),-\},\psi']$, on $O$ by $\{S^{(1)}(x),\psi'(O)\}-\psi'(\{S^{(1)}(x),O\})$. For $\psi'=\tfrac a2\{S^{(1)}(\Delta),-\}$ this is $\tfrac a2\lambda_x\{S^{(1)}(x),-\}$: $+a$, $-\tfrac a2$, $+\tfrac a2$ times the classical transformation along $K_b$, $\Sigma$, $\Sigma^*$ (Proposition \ref{prop:c9-scheme}), a change of the one-body kernels. A finite part $\psi'$ that vanishes on $O$ changes the component of $O$ along $x$ by $-\psi'(\{S^{(1)}(x),O\})$: the finite parts of the operators of the Weyl weight of $\{S^{(1)}(x),O\}$ enter the two-body kernels, as the finite part of Lemma \ref{lem:c11-cocycle} on the trace bilinear and the shift of the Konishi entry in Remark \ref{rem:c11-twist}.
\item (v) The cutoff shape enters the finite parts of the defect cycles (Proposition \ref{prop:c10-shape}, Lemma \ref{lem:c10-special}), of the special-edge cycles of $L_1^{(0)}$ and of the finite special vertices; since the counterterms are unique up to finite local terms [C07, Thm B], a change of the shape is a change of $\hbar F_1$ and of $\psi$, a combination of (i)--(iv).
\item (vi) At the values $\theta=\theta_\lambda=\theta_V=-\tfrac13$ of Theorem \ref{thm:c9-theta2}, in the translation-covariant scheme, $P_1\cong0$ on the families of Definition \ref{def:c11-families} (Corollary \ref{cor:c11-D1fam}), and the condition $P_1=0$ of Definition \ref{def:c8-radial} holds there with $\psi'=0$: the components along $K_b$, $\Sigma$ and $\Sigma^*$ are those of the scheme. For equal coefficients $\theta_P\ne\theta$ along $P_b$, the finite part with $a=(\theta-\theta_P)$ units gives $P_1\cong0$ and changes the components along $K_b$, $\Sigma$, $\Sigma^*$ by $a$, $-\tfrac a2$, $\tfrac a2$ times the classical transformation (Corollary \ref{cor:c11-theta} (iv)).
\item (vii) Under (i), (iv) and the coupling direction the representation of $\mathfrak{psl}(4|4)$ on the cohomology of $\Lambda^{(0)}$ changes by an isomorphism, under $G_B$ by the automorphism $e^{-a\hbar B}$, and the spectrum of $D_1$ in a fixed classical Weyl weight is invariant (Theorem \ref{thm:c8-invariants}); the class $[Q^\xi_1]$ on the cohomology of $Q^\xi_0$, its kernel, image and rank are invariant under (i), (ii) and (iv) (Theorem \ref{thm:c8-twist}).
\end{itemize}
\end{proposition}

\begin{proof}
(i), (ii) Proposition \ref{prop:c8-laws} (a), (c); the bracket with a local functional is a derivation in its second argument, so it acts letter by letter; for $\phi^{(0)}=t\,S^{(1)}(\Delta)$, $\{S^{(1)}(x),S^{(1)}(\Delta)\}$ is $S^{(1)}$ of $\pm[x,\Delta]=\mp\lambda_xx$ by Lemma \ref{lem:c2-components}, with $S^{(2)}(x,\Delta)=0$ (Definition \ref{def:c2-S2}). (iii) Proposition \ref{prop:c11-freexi}, Definition \ref{def:c11-free}, Remark \ref{rem:c10-scheme1}. (iv) Proposition \ref{prop:c8-laws} (b) and Proposition \ref{prop:c9-scheme}, with $[\{S^{(1)}(x),-\},\psi'](O)=\{S^{(1)}(x),\psi'(O)\}-\psi'(\{S^{(1)}(x),O\})$. (v) Remark \ref{rem:c10-scheme1} (d); the finite counterterm and the finite parts of the operators absorb the change. (vi) Corollaries \ref{cor:c11-D1fam} and \ref{cor:c11-theta} (iv). (vii) Theorems \ref{thm:c8-invariants} and \ref{thm:c8-twist}.
\end{proof}

\begin{lemma}[The trace bilinear in a cocycle scheme]
\label{lem:c11-cocycle}
Let $O=\int\omega\,(\iota_{w_1}A,\iota_{w_2}A)$, $\xi\in S_+\otimes W$ or $\xi\in S_-\otimes W^*$, and let $\mathcal B_1,\dots,\mathcal B_4$ be the four functionals of Corollary \ref{cor:c7-Q},
\[
\begin{aligned}
&\mathcal B_1:=\{S^{(1)}(\xi),O\},\qquad \mathcal B_2:=(w_1,w_2)\int\omega\,\big(\rho(A|_{W_{\mathbb R}^\perp})\xi,\lambda\big),\\
&\mathcal B_3:=(w_1,w_2)\int\omega\,\big(\rho(A|_V)\xi,\lambda\big),\qquad \mathcal B_4:=(w_1,w_2)\int\omega\,\big(c,(\xi,\lambda^*)\big),
\end{aligned}
\]
$\mathcal B_1$ being the first functional listed there (Lemma \ref{lem:c11-lines} (i)). Write the two-body part of $Q_1(O)$ as $a_1\mathcal B_1+a_2\mathcal B_2+a_3\mathcal B_3+a_4\mathcal B_4$.
\begin{itemize}
\item (i) If $L_1^{(0)}=0$ on $O$ and on $\mathcal B_1$, then $a_3=a_4=0$, and the one-body terms of $Q_1(O)$ are multiples of $\mathcal B_1$.
\item (ii) If $L_1^{(0)}(O)=2\kappa_O(w_1,w_2)\int\omega\,(dc,A|_V)$ and the part of $L_1^{(0)}(\mathcal B_1)$ with the two letters $c$, $\lambda$ is $-\kappa_{\mathcal B}\,s\mathcal B_3$, the finite part $\psi(O)=\kappa_O(w_1,w_2)\int\omega\,(A|_V,A|_V)$, $\psi(\mathcal B_1)=\kappa_{\mathcal B}\mathcal B_3$ removes both and changes $Q_1(O)$ by $(2\kappa_O-\kappa_{\mathcal B})\mathcal B_3$. At $\psi=0$ the part with the letters $c$, $\lambda$ of (R1) of Proposition \ref{prop:c11-relations} is the identity $a_3+2\kappa_O-\kappa_{\mathcal B}=0$.
\end{itemize}
\end{lemma}

\begin{proof}
(i) By (R1), $s\,Q_1(O)=0$. $s\mathcal B_1=0$, since $[s,\{S^{(1)}(\xi),-\}]=\{\{S,S^{(1)}(\xi)\},-\}=0$ at tree level and $sO=0$; $s\mathcal B_2=0$, since $Q$ annihilates the coordinate functionals of $A|_{W_{\mathbb R}^\perp}$ and $\lambda$ and the cubic parts of $s$ on them are the gauge transformation with parameter $c$, which leaves the invariant pairing fixed. In the sector with the two letters $c$, $\lambda$, $s\mathcal B_3=-(w_1,w_2)\int\omega\,(\rho(dc)\xi,\lambda)$ and $s\mathcal B_4=(w_1,w_2)\int\omega\,(c,(\xi,{\not d}\lambda))$, whose densities carry the derivative on different letters and are linearly independent for arbitrary $\omega$; hence $a_3=a_4=0$. The one-body terms are $-\{\hbar F_1^{(1)}(\xi),O\}$; a term of $\hbar F_1^{(1)}(\xi)$ acting on the letters $A|_{W_{\mathbb R}^\perp}$ of $O$ has the leg $A^*|_{W_{\mathbb R}^\perp}$, and by Proposition \ref{prop:c11-freexi} (i) it is a multiple of $T_2$, which gives a multiple of $\mathcal B_1$. (ii) $Q\int\omega\,(A|_V,A|_V)=-2\int\omega\,(dc,A|_V)$ (Remark \ref{rem:c5-cocycle}), and the cubic part of $s\int\omega\,(A|_V,A|_V)$ vanishes by the invariance of $(-,-)$; $L_1$ changes by $[s,\psi]$ (Lemma \ref{lem:c10-psi}), the component along $\xi$ by $[\{S^{(1)}(\xi),-\},\psi]$ (Proposition \ref{prop:c8-laws} (b)), and $\{S^{(1)}(\xi),\int\omega\,(A|_V,A|_V)\}=2\int\omega\,(\Gamma(\xi,\lambda)|_V,A|_V)$ with $(\Gamma(\xi,\lambda)|_V,a)=(\rho(a)\xi,\lambda)$ [ESW, (28)]. In the shifted scheme (i) gives $a_3+2\kappa_O-\kappa_{\mathcal B}=0$; at $\psi=0$ the part with the letters $c$, $\lambda$ of (R1), $s\,Q_1(O)+L_1^{(0)}(\mathcal B_1)+\{S^{(1)}(\xi),L_1^{(0)}(O)\}=0$, is the same equation, since there $s\,Q_1(O)=a_3\,s\mathcal B_3$, $L_1^{(0)}(\mathcal B_1)=-\kappa_{\mathcal B}\,s\mathcal B_3$ and $\{S^{(1)}(\xi),L_1^{(0)}(O)\}=2\kappa_O\,s\mathcal B_3$.
\end{proof}

\begin{theorem}[$Q_1$ on the scalar bilinears]
\label{thm:c11-bilinears}
Let $O=\int\omega\,(\iota_{w_1}A,\iota_{w_2}A)$, $\xi\in S_+\otimes W$ or $\xi\in S_-\otimes W^*$, and $\mathcal B_1,\dots,\mathcal B_4$ as in Lemma \ref{lem:c11-cocycle}. At $\psi=0$:
\begin{itemize}
\item (i) the two-body part of $Q_1(O)$ is
\begin{equation}
-2L\,\mathcal B_1+4L\,\mathcal B_2+4L\,\mathcal B_3+0\cdot\mathcal B_4\ \text{units},\qquad L=\log\tfrac32,
\label{eq:c11-Q1O}
\end{equation}
without poles and the same for both chiralities; its trace part is $4L$ units times $(w_1,w_2)\int\omega\,(\rho(A)\xi,\lambda)$, with the whole connection $A=A|_V+A|_{W_{\mathbb R}^\perp}$;
\item (ii) on the traceless part, $(w_1,w_2)=0$,
\[
Q_1(O)=c_Q\,\{S^{(1)}(\xi),O\},\qquad c_Q=-2\log\tfrac32-t_\perp\ \text{units},
\]
with $t_\perp$ units the coefficient of $T_2$ in $\hbar F_1^{(1)}(\xi)$, which the two-leg solution of \eqref{eq:c11-arity1xi} gives as $t_\perp=\log\tfrac32-\tfrac{79}{12}+\theta^\xi_\perp$ (Proposition \ref{prop:c11-oddsol} (i)), so that $c_Q=\tfrac{79}{12}-3\log\tfrac32-\theta^\xi_\perp$ units; $c_Q$ moves by $-1$ times the change of the free direction $\theta^\xi_\perp$, by $\mp a$ under $G_B$ along $S_+\otimes W$ and $S_-\otimes W^*$, by $-\tfrac a2$ under $\psi'=\tfrac a2\{S^{(1)}(\Delta),-\}$, and under the field countings as in Proposition \ref{prop:c11-freedom} (i);
\item (iii) $L_1^{(0)}(O)=-(2+4L)$ units times $(w_1,w_2)\int\omega\,(dc,A|_V)$, and the part of $L_1^{(0)}(\mathcal B_1)$ with the letters $c$, $\lambda$ is $2$ units times $s\mathcal B_3$; Lemma \ref{lem:c11-cocycle} (ii) applies with $\kappa_O=-(1+2L)$ units and $\kappa_{\mathcal B}=-2$ units, and in the cocycle scheme
\[
\psi(O)=-(1+2L)\,\text{units}\cdot(w_1,w_2)\int\omega\,(A|_V,A|_V),\qquad \psi(\mathcal B_1)=-2\,\text{units}\cdot\mathcal B_3,
\]
the two-body part of $Q_1(O)$ is $-2L\,\mathcal B_1+4L\,\mathcal B_2$ units, $a_3=4L+2\kappa_O-\kappa_{\mathcal B}=0$;
\item (iv) at $\psi=0$, for $\xi\in S_+\otimes W$, the part with the letters $c$, $\lambda|_{S_-\otimes W^*}$ of (R1) is the identity
\[
\begin{aligned}
&s\,Q_1(O)+L_1^{(0)}\big(\{S^{(1)}(\xi),O\}\big)+\big\{S^{(1)}(\xi),L_1^{(0)}(O)\big\}\\
&\qquad=\big(-4L-2+(2+4L)\big)\,\text{units}\cdot(w_1,w_2)\int\omega\,\big(\rho(dc)\xi,\lambda\big)=0 ;
\end{aligned}
\]
\item (v) at $\psi=0$, with the uniform coefficient $t_\xi=\theta^\xi_V=\theta^\xi_\perp$ forced by the arity-two equations along $(\xi_\pm,\xi'_\pm)$ (Theorem \ref{thm:c11-oddar2}), $Q_1(O)=\big(\tfrac{79}{12}-3L-t_\xi\big)\mathcal B_1+4L\,\mathcal B_2+4L\,\mathcal B_3$ units: the two-body part \eqref{eq:c11-Q1O} and the one-body part $-\big(L-\tfrac{79}{12}+t_\xi\big)\mathcal B_1$ units.
\end{itemize}
\end{theorem}

\begin{proof}
(i) By Theorem \ref{thm:c11-classes} (iii) the two-body part is $\Upsilon_\xi$ on the single pair $(A|_{W_{\mathbb R}^\perp},A|_{W_{\mathbb R}^\perp})$ of $O$, whose leg sets are those of Table \ref{tab:c11-defectpairs}: $\lambda|_{S_-\otimes W^*}$, $A|_V\,\lambda|_{S_-\otimes W^*}$, $A|_{W_{\mathbb R}^\perp}\lambda|_{S_+\otimes W}$ and $c\,\lambda^*|_{S_+\otimes W^*}$ for $\xi\in S_+\otimes W$. The single output $Y$ has the colour factor $(1\otimes\mathrm{ad}_Y)(-,-)^{-1}$, which the colour pairing of $O$ contracts to $\operatorname{tr}_{\mathfrak g}\mathrm{ad}_Y=0$. The other three leg sets have $D=0$, so they carry no poles and no finite special vertex (Proposition \ref{prop:c11-oddbound}), and they give $\mathcal B_1$, $\mathcal B_2$ (outputs $A|_{W_{\mathbb R}^\perp}\lambda$), $\mathcal B_3$ (outputs $A|_V\lambda$) and $\mathcal B_4$ (outputs $c\,\lambda^*$). Their chains have three edges and, at $D=0$, constant integrands, so each chain gives a rational multiple of $\tfrac12\log\tfrac32$ (Theorem \ref{thm:c11-oddsimplex} (iv)); the exact evaluation of the chains in all placements, in both ways of Remark \ref{rem:c11-oddmethod}, gives \eqref{eq:c11-Q1O}, with the same coefficient for $\mathcal B_2$ and $\mathcal B_3$. (ii) By Corollary \ref{cor:c7-Q} the functionals at $(w_1,w_2)=0$ are the multiples of $\mathcal B_1$; the one-body term is $-t_\perp\mathcal B_1$ units (proof of Lemma \ref{lem:c11-cocycle} (i)); the changes are Propositions \ref{prop:c11-freexi} and \ref{prop:c11-freedom}. (iii) $L_1^{(0)}$ is given by the arity-zero two-body kernels, the cycles with one special edge of the proof of Theorem \ref{thm:c10-brst}, of the pairs $(A|_{W_{\mathbb R}^\perp},A|_{W_{\mathbb R}^\perp})$ of $O$ and $(A|_{W_{\mathbb R}^\perp},\lambda)$ of $\mathcal B_1$, evaluated in the same way; then Lemma \ref{lem:c11-cocycle} (ii). (v) By (i) and (ii), $T_2$ being the only term of the solution of Proposition \ref{prop:c11-oddsol} (i) with the leg $A^*|_{W_{\mathbb R}^\perp}$. (iv) With $s\mathcal B_3=-(w_1,w_2)\int\omega\,(\rho(dc)\xi,\lambda)$ in this sector, the three terms are $4L\,s\mathcal B_3$, $2\,s\mathcal B_3$ and $2\kappa_O\,s\mathcal B_3=-(2+4L)\,s\mathcal B_3$ units (proof of Lemma \ref{lem:c11-cocycle} (ii)).
\end{proof}

\begin{proposition}[$S_1$ and $K_1$ on the scalar bilinears]
\label{prop:c11-SKbil}
Let $O=\int\omega\,(\iota_{w_1}A,\iota_{w_2}A)$ and $\psi=0$.
\begin{itemize}
\item (i) $\Upsilon^0_\eta(A|_{W_{\mathbb R}^\perp},A|_{W_{\mathbb R}^\perp})$ has the single leg set $\lambda$, which vanishes on the colour singlet, and $s_0$ has zero bracket with a letter $A|_{W_{\mathbb R}^\perp}$. Hence the two-body part of $S_1(O)$ at the point $x$ of the density is the two-body part of $Q_1(O)$ along the Killing spinor $\rho(x)\eta$, given by \eqref{eq:c11-Q1O} along $S_-\otimes W^*$ for $\eta\in S_+\otimes W^*$ and along $S_+\otimes W$ for $\eta\in S_-\otimes W$; the term linear in ${\not d}(\rho(x)\eta)=4\eta$ is absent (Corollary \ref{cor:c7-S}). With the solution of Proposition \ref{prop:c11-oddsol} (ii), whose position part is $x^\mu\hbar F_1^{(1)}(\rho(e_\mu)\eta)$ with the uniform coefficient $t_\eta$ and whose terms $s_0$, $u_0$ have no leg $A^*$, the component is
\[
S_1(O)=\big(\tfrac{79}{12}-3\log\tfrac32-t_\eta\big)\,\text{units}\cdot\{S^{(1)}(\eta),O\}+x^\mu\Big(4\log\tfrac32\,\mathcal B_2+4\log\tfrac32\,\mathcal B_3\Big)\big(\rho(e_\mu)\eta\big)\ \text{units},
\]
at the point $x$ of the density, with $\mathcal B_2$, $\mathcal B_3$ of Lemma \ref{lem:c11-cocycle} at the spinor $\rho(e_\mu)\eta$; its intrinsic part vanishes.
\item (ii) $\Upsilon^0_{K_b}(A|_{W_{\mathbb R}^\perp},A|_{W_{\mathbb R}^\perp})$ has the leg sets $\emptyset$, of superficial degree $1$, and $A|_V$, both zero on the colour singlet, and $\Upsilon^0_{M_\omega}$ vanishes there. By \eqref{eq:c11-Kdec} the two-body part of $K_1(O)$ is $v(x)^\mu\,\Upsilon_{P_{e_\mu}}-(b,x)\,\Upsilon^0_\Delta$, and the one-body part is the bracket with the solution of Proposition \ref{prop:c11-oddsol} (v), whose terms with the leg $A^*|_{W_{\mathbb R}^\perp}$ are the part $-\tfrac{22}3$ units times $\int(b,x)(A,A^*)$, the sector part $\theta^K_\perp$ units times $S^{(1)}(K_b)_\perp$ and the directions $G^K_j$. On the part with two legs $A|_{W_{\mathbb R}^\perp}$, with $\Upsilon^0_\Delta=\tfrac{40}3O$, $-\tfrac{104}3O$ units at the point of the density and $\Upsilon_{P_{e_\mu}}=-\tfrac13$ units times $\partial_\mu O$ (Corollary \ref{cor:c11-DPdefect}),
\[
\begin{aligned}
&K_1(O)=-\big(\tfrac13+\theta^K_\perp\big)\,\text{units}\cdot\{S^{(1)}(K_b),O\}+\begin{cases}0&\text{(traceless part)},\\ 48\,\text{units}\cdot(b,x)\,O&\text{(trace)},\end{cases}\ -\ \sum_{j=1}^4\kappa^K_j\{G^K_j,O\},\\
&\{S^{(1)}(K_b),O\}=v^\mu\partial_\mu O-4(b,x)\,O,
\end{aligned}
\]
so that at $\theta^K_\perp=-\tfrac13$ and $\kappa^K_j\{G^K_j,O\}=0$, $K_1(O)=-(b,x)D_1(O)$: $0$ on the traceless bilinears and $48$ units times $(b,x)\,O$ on the trace, the relation of Proposition \ref{prop:c8-primary} with $\kappa=-24$ units. The coefficients $\theta^K_d$ and $\kappa^K_j$ enter the arity-two equations along $(K_b,P_{b'})$ and $(\eta,\eta')$ at a finite endpoint.
\end{itemize}
\end{proposition}

\begin{proof}
(i) Table \ref{tab:c11-defectcounts} and the evaluation of the leg set: a single output $Y$ gives $\operatorname{tr}_{\mathfrak g}\mathrm{ad}_Y=0$ on the colour singlet; $s_0$ has the legs $A|_{W_{\mathbb R}^\perp}$, $\lambda^*$, and its bracket with a letter needs the leg $A^*|_{W_{\mathbb R}^\perp}$. Theorem \ref{thm:c11-oddsimplex} (ii) gives the position part. (ii) The leg set $A|_V$ has a single output; the empty leg set is homogeneous in $\varepsilon$ of degree $-\tfrac12$ (Lemma \ref{lem:c10-self}). Theorem \ref{thm:c11-oddsimplex} (iii) gives the two-body parts; by Corollary \ref{cor:c7-K} the term in the second derivatives of $v$ and the term in the antisymmetric part of $dv$ vanish on a scalar bilinear. The one-body part: $-\{-\tfrac{22}3\,\text{units}\int(b,x)(A,A^*),O\}=\tfrac{44}3(b,x)\,O$ units, the rotation and spin terms of the solution acting on $A|_V$ and $\lambda$ alone, and $-\theta^K_\perp$ units times $\{S^{(1)}(K_b),O\}$. With $v(x)^\mu\partial_\mu O$ and $\{S^{(1)}(K_b),O\}=v^\mu\partial_\mu O+\tfrac24(\operatorname{div}v)O$ (Definition \ref{def:c7-vdot}), $-\tfrac13v^\mu\partial_\mu O-\tfrac{40}3(b,x)O+\tfrac{44}3(b,x)O=-\tfrac13\{S^{(1)}(K_b),O\}$ on the traceless part, and on the trace $-(b,x)\Upsilon^0_\Delta$ adds $(\tfrac{104}3+\tfrac{40}3)(b,x)O=48(b,x)O$, the one-body part being the same. Proposition \ref{prop:c8-primary}: $K_1(O)=-2(b,x)\kappa\,O$ for $D_1(O)=2\kappa\,O$.
\end{proof}

\begin{corollary}[The dilatation and the translations by the defect cycles]
\label{cor:c11-DPdefect}
On every local functional $O$ with $sO=0$, with the Koszul signs of \eqref{eq:c10-decomp},
\[
\begin{aligned}
&\sum_{a<b}\pm\Upsilon_\Delta(X_a,X_b)=\sum_{a<b}\pm\big(4\,\mathrm{lc}(X_aX_b)+\text{composite pair kernel along }\Delta\big),\\
&\sum_{a<b}\pm\Upsilon_{P_b}(X_a,X_b)=\sum_{a<b}\pm\big(\text{composite pair kernel along }P_b\big),
\end{aligned}
\]
and the one-body terms of the two forms agree. On the part with two legs $A|_{W_{\mathbb R}^\perp}$ of the scalar bilinears, at the point of the density, $\Upsilon^0_\Delta$ is $\tfrac{40}3$ units times $O$ on the traceless part and $-\tfrac{104}3$ units times $O$ on the trace, and $\Upsilon_{P_b}$ is $-\tfrac13$ units times $\partial_bO$; with the position part $2x^\mu\Upsilon_{P_{e_\mu}}$ this is $\tfrac{44}3O-\tfrac13\{S^{(1)}(\Delta),O\}$ and $-\tfrac{100}3O-\tfrac13\{S^{(1)}(\Delta),O\}$ units, the first row of the table of Theorem \ref{thm:c9-D1}, and the value $-\tfrac13$ of Theorem \ref{thm:c9-P1}.
\end{corollary}

\begin{proof}
Both forms compute the unique local representative (Theorem \ref{thm:c8-L} (i)): for $\Delta$ this is Proposition \ref{prop:c11-routes}, whose route of the special edges is \eqref{eq:c11-decompx}, and for $P_b$ the same argument applies with Corollary \ref{cor:c11-D1local} and the one-body term $-\{\hbar F_1^{(1)}(P_b),-\}$ in both forms. On the bilinears, the defect cycles with $S^{(1)}(\Delta)$ and $S^{(1)}(P_b)$ in all their placements, evaluated by Theorem \ref{thm:c11-oddsimplex}, give the stated values, with $\{S^{(1)}(\Delta),O\}=2x^\mu\partial_\mu O+4O$ at the point of the density: $\tfrac{40}3+\tfrac43=\tfrac{44}3$ and $-\tfrac{104}3+\tfrac43=-\tfrac{100}3$.
\end{proof}

\begin{proposition}[The dilatation and the translations on the families from the defect cycles]
\label{prop:c11-masterfam}
With the master formula \eqref{eq:c11-decompx} along $\Delta$ and $P_n$, the sector coefficients $-\tfrac13$ of Theorem \ref{thm:c9-theta2} and the families and kinematics of Definition \ref{def:c11-families}, on the parts with the fewest fields:
\begin{itemize}
\item (i) on the $\mathfrak{su}(2)$ sector, $L\le8$, $D_1=-8\,\text{units}\cdot\sum_l(1-P_{l,l+1})$, $P_{l,l+1}$ the exchange of the letters at the sites $l$, $l+1$, which is $\frac{g^2N}{4\pi^2}$ times the Heisenberg Hamiltonian; on the $\mathfrak{su}(1|1)$ sector, $L\le6$, on all twenty-two graded cyclic sectors, $D_1=-8\,\text{units}\cdot\sum_l(1-\Pi_{l,l+1})$ with the graded exchange $\Pi_{l,l+1}$;
\item (ii) on twist two, $S'\le6$, the spectrum of $D_1$ is $\{-32\,h(j)\ \text{units}: j\le S'\ \text{even}\}$, and on the fermion $\mathfrak{sl}(2)$ sector, $S'\le5$, it is $\{-32\,h(j+1)\ \text{units}: j\le S'\ \text{odd}\}$;
\item (iii) on the family of two gaugino components with $Z'F_{11}$, $S'\le3$, the eigenvalues are $0$, $-48$, $-\tfrac{200}3$ units, on the gluon $\mathfrak{sl}(2)$ sector, $S'\le2$, $-48$, $-\tfrac{200}3$ units, and on $(F_{11},\psi)$, $S'\le4$, $-48$, $-\tfrac{200}3$, $-\tfrac{392}5$ units: $-8$ units times the eigenvalues of Proposition \ref{prop:c11-targets};
\item (iv) the sum over the pairs of $\Upsilon_{P_n}$ is $-\tfrac13$ units times $\partial_n$ on twist two and on the fermion $\mathfrak{sl}(2)$ sector, and on the families of (iii), and $D_1(\partial_nO)=\partial_nD_1(O)$ on twist two and the families of (ii), (iii), as an identity of matrices from $S'$ to $S'+1$;
\item (v) pair by pair, $\Upsilon_\Delta(X_a,X_b)=4\,\mathrm{lc}(X_aX_b)+\big(\text{composite pair kernel along }\Delta\big)$, $\Upsilon_{P_b}$ and $\Upsilon_{M_\omega}$ equal to the composite pair kernels along $P_b$ and $M_\omega$, and $\Upsilon_R=0$, on every pair of slots with defect cycles and every leg set with at most one letter derivative on which both evaluations are complete, and with at most two letter derivatives in the simplex evaluation.
\end{itemize}
The spectra of (i)--(iii) are those of Table \ref{tab:c11-results}, obtained here without the composite term.
\end{proposition}

\begin{proof}
The two-body kernels are the coefficients of $\varepsilon^0$ of the defect cycles of the pairs of letters with $S^{(1)}(\Delta)$ or $S^{(1)}(P_n)$ at the special edge or on a line, evaluated exactly by Theorem \ref{thm:c11-oddsimplex} with the colour words reduced to the planar adjacent form of Lemma \ref{lem:c10-colour} (iv), and the one-body terms are the brackets with \eqref{eq:c9-F1Delta} and with $\hbar F_1^{(1)}(P_n)$; the operators are diagonalised exactly on bases of the families closed under them, with the projection to cohomology of Remark \ref{rem:c11-method}. Items (i), (ii), (iv) on twist two and the fermion sector and (v) at one letter derivative are obtained in both evaluations of Remark \ref{rem:c11-oddmethod}; the families of (iii), the matrix identity of (iv) on them and (v) at two letter derivatives in the simplex evaluation. Item (v) summed over the pairs of a cocycle is Corollary \ref{cor:c11-DPdefect}.
\end{proof}

\begin{corollary}[The components on the scalar bilinears: representatives and invariant content]
\label{cor:c11-F0}
Let $O=\int\omega\,(\iota_{w_1}A,\iota_{w_2}A)$, $\psi=0$, $L=\log\tfrac32$, $\mathcal B_1,\dots,\mathcal B_4$ as in Lemma \ref{lem:c11-cocycle}, and let the action-side solutions be those of Theorem \ref{thm:c9-quadratic} (iii), Theorem \ref{thm:c9-P1} (i) and Proposition \ref{prop:c11-oddsol}, with the uniform coefficients $t_\Delta$, $t_P$, $t_M$, $t_R$, $t_{\xi_\pm}$, $t_{\eta_\pm}$ of Theorem \ref{thm:c11-oddar2}. The representatives of the components given by the master formula on the part with two legs $A|_{W_{\mathbb R}^\perp}$ are, in units,
\[
\begin{aligned}
&Q_1(O)=\big(\tfrac{79}{12}-3L-t_\xi\big)\mathcal B_1+4L\,\mathcal B_2+4L\,\mathcal B_3+0\cdot\mathcal B_4,\qquad \hbar L_1(O)^{(0)}=-(2+4L)\,(w_1,w_2)\int\omega\,(dc,A|_V),\\
&S_1(O)=\big(\tfrac{79}{12}-3L-t_\eta\big)\{S^{(1)}(\eta),O\}+x_O^\mu\big(4L\,\mathcal B_2+4L\,\mathcal B_3\big)\big(\rho(e_\mu)\eta\big),\\
&D_1(O)=-\big(4t_\Delta+\tfrac43\big)\,O\ \text{(traceless)},\qquad D_1(O)=-\big(4t_\Delta+\tfrac{148}3\big)\,O\ \text{(trace)},\quad\text{at the point }x_O\text{ of the density},\\
&P_1(O)=-\big(\tfrac13+t_P\big)\,\partial_bO,\qquad M_1(O)=-\big(\tfrac13+t_M\big)\,\{S^{(1)}(M_\omega),O\},\qquad R_1(O)=0,\\
&K_1(O)=-\big(\tfrac13+\theta^K_\perp\big)\{S^{(1)}(K_b),O\}+\begin{cases}0&\text{(traceless)},\\ 48\,(b,x)\,O&\text{(trace)},\end{cases}\ -\ \sum_{j=1}^4\kappa^K_j\big\{G^K_j,O\big\},
\end{aligned}
\]
with $t_\xi$, $t_\eta$ the coefficients of the chirality of $\xi$, $\eta$, the spinor $\rho(e_\mu)\eta$ in $\mathcal B_2$, $\mathcal B_3$ in place of $\xi$, and the intrinsic parts of $S_1$ and $K_1$ equal to $0$. Hence:
\begin{itemize}
\item (i) the scheme-independent content is $D_1=0$ on the twenty traceless bilinears and $D_1=-48$ units times $O$ on the trace, $\frac{3g^2N}{2\pi^2}O$ for $\mathfrak g=\mathfrak{su}(N)$, twice the Konishi anomalous dimension: these are the eigenvalues of the induced operator on the cohomology of $\Lambda^{(0)}$, at the value $t_\Delta=-\tfrac13$ fixed by the arity-two equations, invariant under the inner and outer directions and under the finite parts $\psi'$ (Theorem \ref{thm:c8-invariants} (iii));
\item (ii) as functions of the action-side free directions at $\psi=0$, with $c_Q$, $c_S$ the coefficients of $\mathcal B_1$ and of $\{S^{(1)}(\eta),O\}$ and $p_1=-(\tfrac13+t_P)$, on the traceless bilinears
\[
c_Q(\xi_+)+c_Q(\xi_-)-p_1=c_Q(\xi_+)+c_S(\eta_+)-\tfrac14D_1=\tfrac{27}4-8\log\tfrac32,
\]
independent of $t_\Delta$, $t_P$, $t_\xi$, $t_\eta$ and of the inner and outer directions by \eqref{eq:c11-trel}; a finite part $\psi'$ of the operators shifts $c_Q$, $c_S$ and $p_1$ by the coboundaries of Proposition \ref{prop:c11-freedom} (iv), with weight-dependent coefficients on the levels of Weyl weight $2$, $\tfrac52$ and $3$, so these combinations become scheme-independent together with the one-loop components on the weight-$\tfrac52$ level that enter (R2) and (R3) of Proposition \ref{prop:c11-relations};
\item (iii) $P_1$, $M_1$ and $R_1$ vanish at the values $t_P=t_M=t_\Delta=-\tfrac13$ of the translation-covariant scheme, and $L_1^{(0)}(O)$ is removed by the finite part $\psi(O)=-(1+2L)$ units times $(w_1,w_2)\int\omega\,(A|_V,A|_V)$ (Theorem \ref{thm:c11-bilinears} (iii));
\item (iv) the scheme artefacts are $c_Q$ and $c_S$ separately, every coefficient with $\log\tfrac32$, $P_1$, $L_1^{(0)}(O)$, $M_1$, $R_1$, the one-body part of $K_1$, which depends on $\theta^K_\perp$ and the four directions $\kappa^K_j$ left free on two legs, and the individual two-body kernels.
\end{itemize}
\end{corollary}

\begin{proof}
The two-body parts are Theorem \ref{thm:c11-bilinears} (i), Proposition \ref{prop:c11-SKbil} and Corollary \ref{cor:c11-DPdefect}: $\Upsilon^0_\Delta=\tfrac{40}3O$, $-\tfrac{104}3O$, $\Upsilon_{P_b}=-\tfrac13\partial_bO$, $\Upsilon^0_{M_\omega}=0$, $\Upsilon_R=0$, the intrinsic parts along $\eta$ and $K_b$ vanishing on the colour singlet. The one-body parts are the brackets of $O$ with the solutions; on $O$ only the terms with the leg $A^*|_{W_{\mathbb R}^\perp}$ act. Along $\xi$, $T_2$ is the only such term of $\hbar F_1^{(1)}(\xi)$ and $\{T_2,O\}=\mathcal B_1$, which gives $-(L-\tfrac{79}{12}+t_\xi)\mathcal B_1$; along $\eta$ the position part of the solution is $x^\mu\hbar F_1^{(1)}(\rho(e_\mu)\eta)$ with the uniform $t_\eta$, and $s_0$, $u_0$ have no leg $A^*$; along $\Delta$, $-\{\tfrac{22}3\int(A,A^*),O\}=-\tfrac{44}3O$ and $-\{t_\Delta S^{(1)}(\Delta)_\perp,O\}=-4t_\Delta O$ at the point of the density; along $P_b$, $M_\omega$, $R$ the sector parts in the sector of $A|_{W_{\mathbb R}^\perp}$, with $\{S^{(1)}(M_\omega),O\}=(\omega x)^\mu\partial_\mu O$ on a scalar bilinear and the position part $(\omega x_O)^\mu\Upsilon_{P_{e_\mu}}$ of $\Upsilon_{M_\omega}$; along $K_b$, the part linear in $x$ of Proposition \ref{prop:c11-oddsol} (v) gives $\tfrac{44}3(b,x)O$, which with $-(b,x)\Upsilon^0_\Delta$ and $v(x)^\mu\Upsilon_{P_{e_\mu}}$ is $-\tfrac13\{S^{(1)}(K_b),O\}$ on the traceless part and the same plus $48(b,x)O$ on the trace, $\{S^{(1)}(K_b),O\}=v^\mu\partial_\mu O-4(b,x)O$. (i) Theorem \ref{thm:c9-D1} and Theorem \ref{thm:c11-oddar2}. (ii) $c_Q(\xi_+)+c_Q(\xi_-)-p_1=\tfrac{79}6-6L+\tfrac13-(t_{\xi_+}+t_{\xi_-}-t_P)$ and $c_Q(\xi_+)+c_S(\eta_+)-\tfrac14D_1=\tfrac{79}6-6L+\tfrac13-(t_{\xi_+}+t_{\eta_+}-t_\Delta)$, both $\tfrac{27}2-6L-\tfrac{27}4-2L$ by \eqref{eq:c11-trel}. (iii) Theorem \ref{thm:c11-oddar2} (i), (iii). (iv) Theorem \ref{thm:c11-oddar2} (ii), (iv) and Proposition \ref{prop:c11-freedom}.
\end{proof}

\begin{proposition}[The special supercharges and the supercharges]
\label{prop:c11-S}
Let $\eta\in\Sigma^*$ and $b\in V$, so that $\rho(b)\eta\in\Sigma$.
\begin{itemize}
\item (i) Graph by graph, the position part along $b$ of the two-body kernel of $(X_a,X_b)$ along $\eta$ is the two-body kernel of $(X_a,X_b)$ along $\rho(b)\eta$; its translation-invariant part is computed with $S^{(1)}(\eta)$ at the spinor $\rho(x-x_O)\eta$, $x_O$ the point of the density, together with the term $-\int(\rho(A|_{W_{\mathbb R}^\perp})\eta,\lambda^*)$.
\item (ii) Under the hypotheses of Proposition \ref{prop:c7-select} and Corollary \ref{cor:c8-homotopy}, on cohomology, the position part along $b$ of $S_1$ along $\eta$ is $Q_1$ along $\rho(b)\eta$ minus the commutator of $\{S^{(1)}(\eta),-\}$ with $P_1$.
\item (iii) On a family on which $P_1\cong0$, in particular on the families of Definition \ref{def:c11-families} in the translation-covariant scheme at the values $\theta=\theta_\lambda=\theta_V=-\tfrac13$ of Theorem \ref{thm:c9-theta2}, the position part of $S_1$ along $\eta$ is $Q_1$ along $\rho(b)\eta$ on cohomology, and $S_1$ along $\eta$ is fixed by $Q_1$ along $\Sigma$ up to the term of Proposition \ref{prop:c7-select} linear in ${\not d}(\rho(x)\eta)=4\eta$. On the scalar bilinears this term is absent (Corollary \ref{cor:c7-S}), and $S_1$ along $\eta$ is, on cohomology, the combination of the functionals of Corollary \ref{cor:c7-S} with the coefficients of $Q_1$ along $\Sigma$ in Corollary \ref{cor:c7-Q}, given by Theorem \ref{thm:c11-bilinears} and Proposition \ref{prop:c11-SKbil} (i); the component $S_1$ computed from its own defect cycles and from the action-side solution along $\eta$ (Proposition \ref{prop:c11-SKbil} (i)) satisfies this relation.
\item (iv) In a reflection-compatible scheme, on the family formed by a traceless scalar bilinear $O$ at the origin and its descendants $\{S^{(1)}(\xi),O\}$, on which $Q_1=c_Q\{S^{(1)}(\xi),-\}$ along the unit generator $Q^i_\alpha$ (Theorem \ref{thm:c11-bilinears} (ii)) and $S_1=c_S\{S^{(1)}(\eta),-\}$ along $S^\alpha{}_i$, Corollary \ref{cor:c8-gram} reads
\[
c_S=\overline{c_Q}+g_{5/2}-g_2,
\]
with $g_w$ the eigenvalue of $\mathrm{Gram}_{\rm tree}^{-1}\mathrm{Gram}_{\rm 1}$ in Weyl weight $w$ on the family; the outer direction changes $c_Q$ by $-a$ and $c_S$ by $+a$ and preserves the relation for $\bar a=-a$.
\end{itemize}
\end{proposition}

\begin{proof}
(i) By \eqref{eq:c11-Seta}, $S^{(1)}(\eta)=\int x^\mu s_\mu+\int s_0$, where $\int s_\mu$ is $S^{(1)}$ along the constant spinor $\rho(e_\mu)\eta\in\Sigma$. Insert $x^\mu=x_O^\mu+(x-x_O)^\mu$ and argue as in the proof of Proposition \ref{prop:c11-position} (i): on plane waves the factor $x^\mu$ is $-i\partial/\partial r^\mu$ of the momentum injected at the vertex, and its action on the momentum delta function is $-i$ times the graph with $\int s_\mu$; the same holds for the special edges, the vertices $S^{(1)}(\eta)$ on the cycle and the counterterm vertices, which are the defect insertions of Proposition \ref{prop:c10-arity1} (Theorem \ref{thm:c11-oddsimplex} (ii)). (ii) Corollary \ref{cor:c8-homotopy} (ii) at $(\eta,P_b)$: by Proposition \ref{prop:c2-special}, $[P_b,\eta]=\rho(b)\eta$, so $[\eta,P_b]=-\rho(b)\eta$, and $S^{(2)}(\eta,P_b)=0$ (Definition \ref{def:c2-S2}). At order $\hbar$ the identity states that the commutator of $S_1$ along $\eta$ with $\{S^{(1)}(P_b),-\}$, plus the commutator of $\{S^{(1)}(\eta),-\}$ with $P_1$, minus $Q_1$ along $\rho(b)\eta$, is a graded commutator with $\Lambda^{(0)}$, which induces $0$ on cohomology; at order $\hbar^0$ the same identity is the statement of (i) for the classical transformation, which fixes the signs. By Proposition \ref{prop:c10-bound1}, $S_1$ along $\eta$ is at most linear in the position, and Lemma \ref{lem:c11-canonical} identifies the first commutator with its position part. (iii) By Proposition \ref{prop:c7-select} the representative of $S_1$ along $\eta$ at a point $x_0$ is a term linear in the value $\rho(x_0)\eta$ plus a term linear in $4\eta$; the first has the position part given by its value at $\rho(b)\eta$, and as $\eta$ and $b$ vary, $\rho(b)\eta$ spans $\Sigma$, so by (ii) it is the representative of $Q_1$ evaluated at $\rho(x_0)\eta$ on cohomology. (iv) Write $X_0+\hbar X_1$ for the matrix of $-\Lambda^{(1)}(Q^i_\alpha)$ and $Y_0+\hbar Y_1$ for that of $-\Lambda^{(1)}(-S^\alpha{}_i)$ (Theorem \ref{thm:c8-dag} (i)): $\hbar X_1=c_QX_0$ and $\hbar Y_1=c_SY_0$, and $Y_0$ maps Weyl weight $\tfrac52$ to Weyl weight $2$. Corollary \ref{cor:c8-gram} gives $\hbar Y_1=\overline{c_Q}\,X_0^{\rm adj}+\hbar[Y_0,\mathrm{Gram}_{\rm tree}^{-1}\mathrm{Gram}_{\rm 1}]=\big(\overline{c_Q}+g_{5/2}-g_2\big)Y_0$ on the descendants, with $Y_0=X_0^{\rm adj}$. The outer direction: Proposition \ref{prop:c11-freedom} (ii) and Corollary \ref{cor:c8-gram}.
\end{proof}

\begin{corollary}[$K_1$ in the translation-covariant scheme]
\label{cor:c11-K}
Under the hypotheses of Theorems \ref{thm:c8-adjoint} and \ref{thm:c11-arity2}, on a family of Definition \ref{def:c11-families} at the origin of one length, in the translation-covariant scheme at the values $\theta=\theta_\lambda=\theta_V=-\tfrac13$ of Theorem \ref{thm:c9-theta2}, $K_1$ along $K_{\bar b}$ equals $\hbar\big[\{S^{(1)}(K_{\bar b}),-\},\mathrm{Gram}_{\rm tree}^{-1}\mathrm{Gram}_{\rm 1}\big]$, the statement $K=P^\dagger$ of Corollary \ref{cor:c8-gram}. For equal coefficients $\theta_P$ along $P_b$, $K_1$ along $K_{\bar b}$ is $(\theta-\theta_P)$ units times $\{S^{(1)}(K_{\bar b}),-\}$ plus the same commutator, with the radial Gram matrix of Definition \ref{def:c8-inner} in that scheme. On a scalar bilinear at the origin both sides vanish: the defect cycles give $K_1(O)=0$ there (Proposition \ref{prop:c11-SKbil} (ii)), and the commutator vanishes on a primary at the origin on which $\mathrm{Gram}_{\rm tree}^{-1}\mathrm{Gram}_{\rm 1}$ is diagonal. On its descendants $\partial_\mu O$ the relation equates the intrinsic kernels $\Upsilon^0_{K_b}$ of the descendant pairs with the difference of the two-point normalisations in Weyl weights $w_O+1$ and $w_O$ times the classical transformation. The component $K_1$ itself is computed from its own defect cycles at a finite endpoint and from the action-side solution along $K_b$ (Proposition \ref{prop:c11-SKbil} (ii)); this corollary is a cross-check of it through the inner product and enters no number.
\end{corollary}

\begin{proof}
Corollary \ref{cor:c11-D1fam} gives $P_1\cong0$, resp. $P_1\cong(\theta-\theta_P)$ units times $\{S^{(1)}(P_b),-\}$ with a real coefficient; Corollary \ref{cor:c8-gram} gives the two relations. At the origin $v(0)=0$ and $(b,0)=0$, so the formula of Proposition \ref{prop:c11-SKbil} (ii) vanishes; $\{S^{(1)}(K_{\bar b}),-\}$ lowers the Weyl weight by one and annihilates a primary at the origin.
\end{proof}

\begin{proposition}[The relations at order $\hbar$]
\label{prop:c11-relations}
Write $\Lambda^{(1)}(x)=\Lambda^{(1)}_0(x)+\hbar\Lambda^{(1)}_1(x)$ with $\Lambda^{(1)}_0(x)=\pm\{S^{(1)}(x),-\}$ (Theorem \ref{thm:c8-L} (iv)), so that $\hbar\Lambda^{(1)}_1(x)O$ is the component of $O$ along $x$, and $\Lambda^{(0)}=s+\hbar L_1^{(0)}$. Let $\mathfrak F$ be a family of colour-singlet local functionals of degree $0$, stable under $\Lambda$ and under the classical transformations, with $L_1^{(0)}=0$ on $\mathfrak F$ for a finite part $\psi$ as in Definition \ref{def:c10-hyp} (ii), and let $O\in\mathfrak F$ with $sO=0$.
\begin{itemize}
\item (R1) $s\big(\hbar\Lambda^{(1)}_1(x)O\big)=0$ for every generator $x$; along $\xi$, $s\,Q_1(O)=0$, an identity of local functionals.
\item (R2) For $\xi,\xi'\in\Sigma$ with $b=\Gamma(\xi,\xi')|_V$, $[\Lambda^{(1)}_1(\xi),\Lambda^{(1)}_0(\xi')]O+[\Lambda^{(1)}_0(\xi),\Lambda^{(1)}_1(\xi')]O+\Lambda^{(1)}_1(P_b)O=\Lambda^{(0)}(\chi)$, with $\chi$ local of degree $-1$, built from the arity-two component of $O$ along $(\xi,\xi')$; for $\xi,\xi'\in S_+\otimes W$, $b=0$.
\item (R3) For $\xi\in\Sigma$ and $\eta\in\Sigma^*$ with $[\xi,\eta]=\tfrac12(\xi,\eta)\Delta+M_{\xi\eta}+R_{\xi\eta}$ (Proposition \ref{prop:c2-special}), $M_{\xi\eta}\in\mathfrak{so}(V)$, $R_{\xi\eta}\in\mathfrak{so}(W_{\mathbb R}^\perp)$,
\[
[\Lambda^{(1)}_1(\xi),\Lambda^{(1)}_0(\eta)]O+[\Lambda^{(1)}_0(\xi),\Lambda^{(1)}_1(\eta)]O+\tfrac12(\xi,\eta)\,D_1(O)+M_1(O)+R_1(O)=\Lambda^{(0)}(\chi'),
\]
the components $M_1$, $R_1$ taken along $M_{\xi\eta}$, $R_{\xi\eta}$.
\item (R4) For a rank-one $\xi$, $Q^\xi_0Q^\xi_1+Q^\xi_1Q^\xi_0=0$ on every local functional (Theorem \ref{thm:c8-twist}).
\item (R5) With $[K_b,P_{b'}]=(b,b')\Delta+M_{bb'}$ (Remark \ref{rem:c7-KDelta}), $[\Lambda^{(1)}_1(K_b),\Lambda^{(1)}_0(P_{b'})]+[\Lambda^{(1)}_0(K_b),\Lambda^{(1)}_1(P_{b'})]+(b,b')\hbar\Lambda^{(1)}_1(\Delta)+\hbar\Lambda^{(1)}_1(M_{bb'})\cong0$ on cohomology; with Theorem \ref{thm:c11-oddsimplex} (iii) it states that the parts of $K_1$ multiplying the position are those of \eqref{eq:c11-Kdec}.
\item (R6) For $\eta_1,\eta_2\in\Sigma^*$ with $[\eta_1,\eta_2]=-K_b$, $b=\Gamma(\eta_1,\eta_2)|_V$, at a finite $T$, $K_1$ along $K_b$ is, on cohomology, $[\Lambda^{(1)}_1(\eta_1),\Lambda^{(1)}_0(\eta_2)]+[\Lambda^{(1)}_0(\eta_1),\Lambda^{(1)}_1(\eta_2)]$.
\end{itemize}
In (R2), (R3), (R5) and (R6) the right sides vanish on cohomology. For $O$ without letters $\lambda$ and $c$, such as the scalar bilinears and the families of scalars, the parts of order $\hbar^0$ of $\chi$, $\chi'$ vanish, the right sides are $s$-exact, and their parts with two field letters and no antifield vanish on free on-shell plane waves, $|p|^2=0$ on the scalars and the gauge field, $\rho(p)\lambda=0$ and $(p,a)=0$ on the polarisations $a$ of $A|_V$. On the scalar bilinears at $\theta=-\tfrac13$, (R3) contains $D_1(O)=0$ on the traceless part and $-48$ units times $O$ on the trace (Corollary \ref{cor:c9-ident}), the Konishi relation between $Q_1$, $S_1$ and the anomalous dimension, and (R2) with $\xi\in S_+\otimes W$, $\xi'\in S_-\otimes W^*$ contains $P_1(O)=-(\tfrac13+\theta_P)$ units times $\partial_bO$ (Theorem \ref{thm:c9-P1}).
\end{proposition}

\begin{proof}
Corollary \ref{cor:c8-homotopy} (i) at order $\hbar$ reads $[s,\Lambda^{(1)}_1(x)]+[L_1^{(0)},\Lambda^{(1)}_0(x)]=0$; on $O$ with $sO=0$, $L_1^{(0)}O=0$ and $L_1^{(0)}\{S^{(1)}(x),O\}=0$ (stability of $\mathfrak F$) this is (R1). Corollary \ref{cor:c8-homotopy} (ii) at order $\hbar$ for the pairs $(\xi,\xi')$, $(\xi,\eta)$, $(K_b,P_{b'})$, $(\eta_1,\eta_2)$, with the brackets of Proposition \ref{prop:c2-special} and Remark \ref{rem:c7-KDelta}, has the right side $[\Lambda^{(0)},H]$ with $H$ built from $\Lambda^{(2)}$, which on $O$ with $\Lambda^{(0)}O=0$ is $\Lambda^{(0)}(HO)$; the part of order $\hbar^0$ of $H$ is $\pm\{S^{(2)}(x_1,x_2),-\}$, and $S^{(2)}$ has the legs $(\lambda^*,\lambda^*)$ and $(A,c^*)$ (Definition \ref{def:c2-S2}), whose bracket with a functional of the fields $A$ alone vanishes. (R4) is Theorem \ref{thm:c8-twist}. For the vanishing on shell: the part of $s\chi$ with two field letters is $Q\chi_2+\{I_3,\chi_1\}$, with $\chi_k$ the part of $\chi$ with $k$ letters; $\chi_1$ is a single adjoint letter contracted with a constant tensor, zero for a colour singlet; $\chi_2$ of degree $-1$ has one antifield and one field, or $c^*$ and $c$, and $Q$ substitutes $A^*\mapsto-\tfrac12d^*dA$, $\lambda^*\mapsto{\not d}\lambda$, $c^*\mapsto d^*(A^*|_V)$, $A\mapsto-dc$ (Proposition \ref{prop:c1-Q}), so the part of $Q\chi_2$ without antifields and ghosts is a free field equation times a field. The values on the bilinears are Corollary \ref{cor:c9-ident} and Theorem \ref{thm:c9-P1}, with $R_1=0$ on them (Remark \ref{rem:c7-internal}).
\end{proof}

\begin{proposition}[The chiral-primary bilinears along a rank-one supercharge]
\label{prop:c11-plane}
Let $\xi\in S_+\otimes W$ be of rank one, $\xi_\alpha{}^i=u_\alpha z^i$, and let $w_1,w_2\in W_{\mathbb R}^\perp\otimes\mathbb C$ satisfy $\rho(w_1)\xi=\rho(w_2)\xi=0$. Then $(w_1,w_2)=0$, and under the hypotheses of Proposition \ref{prop:c7-select} the component along $\xi$ of $O=\int\omega\,(\iota_{w_1}A,\iota_{w_2}A)$ vanishes at every order at which they hold: its classical part $\{S^{(1)}(\xi),O\}$ and $Q_1(O)$ are $0$, in every scheme in which the counterterms are equivariant for translations, $\mathfrak{so}(V)$, $\mathfrak{so}(W_{\mathbb R}^\perp)$ and $\mathfrak g$. At one loop the order-$\hbar$ part of the twisted differential $Q^\xi_0+\hbar Q^\xi_1$ (Definition \ref{def:c8-twisted}) vanishes on $O$ identically: the arity-one component is $Q_1(O)=0$, and the arity-two component along $(\xi,\xi)$ has no output with two letters. This holds for every free direction of the one-loop quantisation, for every finite part $\psi$ and independently of the normalisation of the transport map $\Phi$ of Theorem \ref{thm:c11-twistcheck}; through $\Phi$, $O$ is the untwisted form of $\operatorname{tr}\gamma^{(I}\gamma^{J)}$, and the statement is the case $n=2$ of [BK, Sec. 3.4.1, (3.77)] at the level of classes.
\end{proposition}

\begin{proof}
$\rho(w_1)\rho(w_2)+\rho(w_2)\rho(w_1)=2(w_1,w_2)$ applied to $\xi\ne0$ gives $(w_1,w_2)=0$. By Corollary \ref{cor:c7-Q}, for $(w_1,w_2)=0$ the component is a multiple of $\int\omega\,\big[(\iota_{w_1}A,(\rho(w_2)\xi,\lambda))+(\iota_{w_2}A,(\rho(w_1)\xi,\lambda))\big]$, which vanishes; the classical component is the same functional (proof of Corollary \ref{cor:c7-Q}). In the evaluation of Theorem \ref{thm:c11-bilinears} the two-body part on these bilinears, $-2L\,\mathcal B_1+4L(\mathcal B_2+\mathcal B_3)$ units, vanishes term by term, with $\mathcal B_1=\{S^{(1)}(\xi),O\}=0$ and $(w_1,w_2)=0$, and the one-body term is a multiple of $\mathcal B_1$; by Corollary \ref{cor:c11-F0} the same holds with the computed solution and every free direction, the coefficients of $\mathcal B_1,\dots,\mathcal B_4$ multiplying functionals that vanish here, and $\hbar L_1(O)^{(0)}$ vanishes with $(w_1,w_2)$. The component along $(\xi,\xi)$ of a functional of degree $0$ and Weyl weight $2$ has degree $-1$ and Weyl weight $3$ (Proposition \ref{prop:c7-select}, applied once for each argument), so an output with two letters contains an antifield, of Weyl weight at least $\tfrac52$, and a field of weight at least $1$, of total weight at least $\tfrac72>3$; the one-body term $-\{\hbar F_1^{(2)}(\xi,\xi),O\}$ vanishes, a quadratic functional of degree $-2$ and Weyl weight $5$ having the legs $(\lambda^*,\lambda^*)$ or $(A,c^*)$ and no leg $A^*$. $\Phi$ maps classes to classes, so a vanishing differential stays zero under every normalisation.
\end{proof}

\begin{remark}[The pair entries of the holomorphic twist]
\label{rem:c11-twist}
The twisted one-loop differential is $\hbar Q^\xi_1(O)=\hbar L_1(O)^{(0)}+Q_1(O)-\tfrac12\hbar L_1(O)^{(2)}(\xi,\xi)-\cdots$ (Definition \ref{def:c8-twisted}), so its pair rule [BK, (3.7)--(3.8)] corresponds to the two-body kernels of $Q_1$ along $\xi$ and of the arity-two component along $(\xi,\xi)$. By Proposition \ref{prop:c7-select}, applied once for each argument, a component of arity $k$ along $\xi$ of an operator of Weyl weight $w$ and degree $d$ has representatives of Weyl weight $w+\tfrac k2$ and degree $d+1-k$. [BK, (3.51)] relates the holomorphic letters to the fields up to field redefinitions; with $\xi_\alpha{}^i=u_\alpha z^i$, $u'\in S_+$ with $u'^\alpha u_\alpha=1$, and the representatives $\gamma\leftrightarrow A_{ij}z^j$ (Weyl weight $1$, degree $0$), $\beta\leftrightarrow u'^\alpha\lambda_\alpha{}^i$ modulo $z^i$ $(\tfrac32,0)$, $\partial_{\dot\alpha}c\leftrightarrow\bar\lambda_{\dot\alpha i}z^i$ $(\tfrac32,0)$, $c\leftrightarrow c$ $(0,1)$, $b\leftrightarrow u'^\alpha F_{mn}(\sigma^m\bar\sigma^n)_\alpha{}^\beta u'_\beta$ $(2,0)$ or a component of $\lambda^*$ $(\tfrac52,-1)$, and $\partial_{\dot\alpha}\leftrightarrow u'^\alpha\sigma^m_{\alpha\dot\alpha}D_m$ $(1,0)$, the pair entries of [BK, (3.64)--(3.69)] lie in the following arities:

\begin{tabular}{l|c|c|c}
entry & Weyl weight and degree of the pair & of the output & arity\\
\hline
$b\,c$ [BK, (3.64)] & $(2,1)$ & $(3,0)$ & $2$\\
$b\,b$ [BK, (3.65)] & $(4,0)$ & $(\tfrac92,0)$ & $1$\\
$\beta\,\gamma$ [BK, (3.66)] & $(\tfrac52,0)$ & $(3,0)$ & $1$\\
$\beta\,\beta$ [BK, (3.67)] & $(3,0)$ & $(\tfrac72,0)$ & $1$\\
$b\,\gamma$ [BK, (3.68)] & $(3,0)$ & $(\tfrac72,0)$ & $1$\\
$\beta\,b$ [BK, (3.69)] & $(\tfrac72,0)$ & $(4,0)$ & $1$\\
\end{tabular}
\tabcaption{The arities of the one-loop pair entries of the holomorphic twist, with $b$ represented by the self-dual field strength; with $b$ represented by a component of $\lambda^*$ the entry $b\,c$ has arity $1$.\label{tab:c11-twist}}

The entry $\beta\gamma$, the Konishi entry, is the two-body kernel of $Q_1$ along $\xi$ on the pair $(\lambda|_{S_+\otimes W},A|_{W_{\mathbb R}^\perp})$ with two outputs $\lambda|_{S_-\otimes W^*}$: the coefficient of $\int\omega\,(\bar\lambda_{\dot\alpha i}z^i,\bar\lambda^{\dot\alpha}{}_jz^j)$ in $Q_1$ of $\int\omega\,(u'^\alpha\lambda_\alpha{}^i,A_{ij}z^j)$. The classical component of this pair along $\xi$ is a number $\kappa$ times a cubic functional of the scalars, and a finite part of the operators of Weyl weight $3$ that adds $y\int\omega\,(\bar\lambda_{\dot\alpha i}z^i,\bar\lambda^{\dot\alpha}{}_jz^j)$ to that cubic functional changes the entry by $-\kappa y$ (Proposition \ref{prop:c11-freedom} (iv)), the two functionals having the same Weyl weight and quantum numbers. The pair entries are therefore compared through the class $[Q^\xi_1]$ on the cohomology of $Q^\xi_0$ (Theorem \ref{thm:c8-twist}), transported to the holomorphic letters through [ESW, Thm 10.9] (remark after Theorem \ref{thm:c8-twist}). Proposition \ref{prop:c11-plane} is the exact untwisted form, for two letters, of [BK, (3.77)], $Q_1(\gamma^I\gamma^J)=0$: for $\rho(w_a)\xi=0$, $\iota_{w_a}A$ is a combination of the components $A_{ij}z^j$, the letters $\gamma$. Theorem \ref{thm:c11-twistentries} gives the arity-one entries at $\psi=0$, and Theorem \ref{thm:c11-twistcheck} states what their comparison with [BK] tests.
\end{remark}

\begin{theorem}[The pair entries along a rank-one supercharge]
\label{thm:c11-twistentries}
Let $\xi=u\otimes z\in S_+\otimes W$ be of rank one, with $u'$, the twisted letters and their representatives of Remark \ref{rem:c11-twist}, and let $\psi=0$. Normalise the letters as the components along the weight vectors of a Cartan subalgebra of $\mathfrak{so}(V)\oplus\mathfrak{so}(W_{\mathbb R}^\perp)$ for which $u$, $u'$ and $z$ are weight vectors, read the outputs on the dual weight vectors, take $b$ as the component $F(d_1,d_2)$ of the field strength on the two null vectors $d_1$, $d_2$ of $V\otimes\mathbb C$ along which $\partial_{\dot\alpha}$ differentiates, and expand the output momenta in the null coordinates; this fixes each entry up to one nonzero factor for each letter and each output. Reduce modulo $\operatorname{Im}Q^\xi_0$, total derivatives kept: the derivatives $u^\alpha\sigma^m_{\alpha\dot\alpha}\partial_m$ are $Q^\xi_0$-exact on the letters, and so is, up to terms with antifields, the holomorphic Dirac component of a derivative on an output $\partial_{\dot\alpha}c$. Then the arity-one pair entries of $Q_1$ are those of Table \ref{tab:c11-twistentries}, and:
\begin{itemize}
\item (i) the entry $\beta\gamma$, the kernel of $(\lambda|_{S_+\otimes W},A|_{W_{\mathbb R}^\perp})$ on the leg set $(\lambda|_{S_-\otimes W^*})^2$, vanishes for every $\xi\in S_+\otimes W$, the placements on a line and at the special edge cancelling chain by chain; the one-body terms and $L_1^{(0)}$ have outputs of other Weyl weights or with a spectator letter, so the Konishi entry of Remark \ref{rem:c11-twist} is $0$ at $\psi=0$; the entry $bc$ with $b$ represented by $(\xi,\lambda^*)$ and the entry $\gamma\gamma$ vanish as well;
\item (ii) the entries $\beta\beta$ and $\beta b\to\gamma\gamma$ carry the colour and derivative structure of [BK, (3.67)] and of the last term of [BK, (3.69)]: the output $\gamma^K$, resp. the outputs $\gamma^J$, $\gamma^K$, with the indices complementary to those of the letters $\beta$, the same value in the two colour orders, and, for $\beta b\to\gamma\gamma$, the antisymmetric combination of the two holomorphic derivatives;
\item (iii) the entry $b\gamma$ has different values in the two colour orders, each with a term in $\log\tfrac32$, and the entry $\beta b\to\beta\,\partial c$ is nonzero in one colour order;
\item (iv) every derivative in the entries is holomorphic.
\end{itemize}
Statements (i)--(iv) are independent of the normalisation factors; the values of the individual entries depend on the scheme (Proposition \ref{prop:c11-freedom} (iv), Remark \ref{rem:c11-twist}).
\end{theorem}

\begin{tabular}{l|l|l|l|l}
entry & pair $\to$ outputs & order $(Y_1,Y_2)$ & order $(Y_2,Y_1)$ & Dirac components\\
\hline
$\beta\gamma$ [BK, (3.66)] & $(\beta_I,\gamma^J)\to\partial c\,\partial c$ & $0$ & $0$ & $0$, $0$\\
$bc$ [BK, (3.64)], $b\leftrightarrow(\xi,\lambda^*)$ & $(c,b)\to\partial c\,\partial c$ & $0$ & $0$ & $0$, $0$\\
$\beta\beta$ [BK, (3.67)] & $(\beta_I,\beta_J)\to\gamma^K\,\partial c$ & $\tfrac23$ & $\tfrac23$ & $\tfrac13-2L$, $\tfrac13-2L$\\
$b\gamma$ [BK, (3.68)] & $(b,\gamma^I)\to\gamma^I\,\partial c$ & $\tfrac{64}3-128L$ & $32+64L$ & $-\tfrac{64}3-128L$, $-\tfrac{32}3+64L$\\
$\beta b\to\beta\,\partial c$ [BK, (3.69)] & $(\beta_I,b)\to\beta_I\,\partial c$ & $-\tfrac{64}3$ & $0$ & $0$, $0$\\
$\beta b\to\gamma\gamma$ [BK, (3.69)] & $(\beta_I,b)\to\gamma^J\gamma^K$ & $\tfrac83\big(\partial_2\gamma^J\partial_1\gamma^K-\partial_1\gamma^J\partial_2\gamma^K\big)$ & the same & \\
$\gamma\gamma$ [BK, (3.77)] & $(\gamma^I,\gamma^J)\to$ twisted letters & $0$ & $0$ & \\
\end{tabular}
\tabcaption{The arity-one pair entries of $Q_1$ along $\xi=u\otimes z$ at $\psi=0$, reduced modulo $\operatorname{Im}Q^\xi_0$, in the normalisation of Theorem \ref{thm:c11-twistentries}: coefficients of $\frac{\hbar}{16\pi^2}$ times the colour factor $(1\otimes\mathrm{ad}_{Y_2}\mathrm{ad}_{Y_1})(-,-)^{-1}$ of the colour order $(Y_1,Y_2)$, $Y_1$ the first output listed, with $L=\log\tfrac32$. The third and fourth columns give the coefficient of the holomorphic derivative on the output other than $\partial c$, and for $\beta b\to\gamma\gamma$ the derivative structure, $\partial_1$ and $\partial_2$ being the derivatives along $d_1$ and $d_2$; the last column gives, in the two orders, the coefficient of the holomorphic Dirac component on the output $\partial c$, which is $Q^\xi_0$-exact up to terms with antifields; $K$ is complementary to $I,J$, and $J,K$ to $I$.\label{tab:c11-twistentries}}

\begin{proof}
The entries are the kernels on the leg sets $(\lambda|_{S_-\otimes W^*})^2$ of $(\lambda|_{S_+\otimes W},A|_{W_{\mathbb R}^\perp})$; $A|_{W_{\mathbb R}^\perp}\lambda|_{S_-\otimes W^*}$ of $(\lambda|_{S_+\otimes W},\lambda|_{S_+\otimes W})$ and of $(A|_V,A|_{W_{\mathbb R}^\perp})$ with one derivative on $A|_V$; $\lambda|_{S_+\otimes W}\lambda|_{S_-\otimes W^*}$ and $(A|_{W_{\mathbb R}^\perp})^2$ of $(\lambda|_{S_+\otimes W},A|_V)$ with one derivative on $A|_V$; and $(\lambda|_{S_-\otimes W^*})^2$ of $(c,\lambda^*|_{S_+\otimes W^*})$ (Tables \ref{tab:c11-defectpairs} and \ref{tab:c11-defectcounts}); all have superficial degree $0$, so they carry no poles and no finite special vertex (Proposition \ref{prop:c11-oddbound}), and Theorem \ref{thm:c11-oddsimplex} evaluates them exactly, with the letter vectors contracted exactly in the weight basis (Remark \ref{rem:c11-oddmethod}). The colour factors are those of Lemma \ref{lem:c10-colour} with the output of a term $T_5$ read as an output of a cubic vertex. (i) The vanishing of $\beta\gamma$ for general $\xi$ and of $bc$ holds in both evaluations. The outputs $(\lambda|_{S_-\otimes W^*})^2$ have Weyl weight $3$, while $L_1^{(0)}$ preserves the Weyl weight $\tfrac52$ of the Konishi descendant, a one-body term keeps the other letter as a spectator, and the arity-two kernels have outputs of degree $-1$; hence the coefficient of $\int\omega\,(\bar\lambda_{\dot\alpha i}z^i,\bar\lambda^{\dot\alpha}{}_jz^j)$ in $\hbar Q^\xi_1$ of the Konishi descendant is the entry $\beta\gamma$. $\gamma\gamma$ is Proposition \ref{prop:c11-plane}. (ii)--(iv) are read off the table. Rescaling a weight vector rescales its letter by the inverse factor and its output by the factor, which multiplies an entry by one product of factors, the same for the two colour orders, for the reduced and the Dirac coefficients, and for the rational part and the part in $\log\tfrac32$; the vanishing, the equality of the two orders, the indices and the derivative structure are therefore independent of the normalisation.
\end{proof}

\begin{theorem}[What the comparison with the holomorphic twist tests]
\label{thm:c11-twistcheck}
Let $\xi=u\otimes z$ be of rank one, let $\Phi$ be the quasi-isomorphism from the cohomology of $Q^\xi_0$ on local operators to that of the free holomorphic differential on the letters $c$, $\gamma$, $\beta$, $b$, $\partial_{\dot\alpha}$ of [BK], transported through [ESW, Thm 10.9] (remark after Theorem \ref{thm:c8-twist}), and let $D$ be the difference, read through $\Phi$, of the one-loop quantisation $r_\xi(\hbar F_1)$ of the twisted theory restricted from the equivariant one and the holomorphic quantisation of [BK, (2.42), footnote 22] and [W18, Thm 3.4] (B. R. Williams, Renormalization for holomorphic field theories, arXiv:1809.02661). If $D\in\mathbb C\,S^\xi+\operatorname{Im}Q^\xi_0$, then $\Phi\circ[Q^\xi_1]\circ\Phi^{-1}$ is the class of the one-loop differential of [BK, (3.7), (3.64)--(3.69)], and:
\begin{itemize}
\item (a) the one-body terms of $Q^\xi_1$ on the twisted letters are $a\,Q^\xi_0+[Q^\xi_0,N]$ for a number $a$ and a derivation $N$ of the letters, zero on cohomology;
\item (b) on the pairs from $\{c,\gamma\}$ and on $(\beta,c)$ the two-body kernels lie, on the twisted letters, in $\operatorname{Im}Q^\xi_0$ plus terms cancelled on cohomology by the one-body terms, as in [BK, (3.11)]; on $(\gamma,\gamma)$ they vanish (Proposition \ref{prop:c11-plane});
\item (c) on the pairs $\beta\gamma$, $\beta\beta$, $b\gamma$, $\beta b$, $bb$ and $bc$ the kernels, reduced modulo $\operatorname{Im}Q^\xi_0$ as in Theorem \ref{thm:c11-twistentries}, have the colour factors $(1\otimes\mathrm{ad}_{Y_2}\mathrm{ad}_{Y_1})(-,-)^{-1}$, which are $\tfrac12f_{ACD}f_{BCE}$ of [BK] up to normalisation, and holomorphic derivatives $\partial_{\dot\alpha}=u'^\alpha\sigma^m_{\alpha\dot\alpha}\partial_m$ on the outputs;
\item (d) their coefficients satisfy, modulo the finite parts of Proposition \ref{prop:c11-freedom} (iv), the relations [BK, (3.59)--(3.63)], which are the order-$\hbar$ part of Corollary \ref{cor:c8-homotopy} (ii) at $(\xi,R)$ for the elements $R$ of the tree-level commutant $\mathfrak{psl}(3|3)$ of $\xi$, restricted by $r_\xi$;
\item (e) every statement of the holomorphic twist about $[Q_1]$ on classes of several letters is the image under $\Phi$ of $[Q^\xi_1]$, computed by the pair rule from the kernels of (c), with the normalisation of $\hbar$ relative to the unit fixed by $\Phi$.
\end{itemize}
For a general $D$ the two class maps differ by $[\{D,-\}]$, and (a)--(e) compute that class. On the chiral bilinears the comparison is complete at the level of classes: by Proposition \ref{prop:c11-plane} the order-$\hbar$ twisted differential vanishes on them for every free direction of the equivariant quantisation and every normalisation of $\Phi$, which is the case $n=2$ of [BK, (3.77)].
\end{theorem}

\begin{proof}
Two quantisations of the twisted theory differing by a $Q^\xi_0$-closed $D$ of degree $0$ have class maps differing by $[\{D,-\}]$, which vanishes for $D\in\mathbb C\,S^\xi+\operatorname{Im}Q^\xi_0$, since $\{S^\xi,-\}=Q^\xi_0$ and $\{Q^\xi_0\chi,-\}=\pm[Q^\xi_0,\{\chi,-\}]$ (Theorem \ref{thm:c8-twist} and the remark after it). (a) The one-body terms are brackets with local functionals (Theorem \ref{thm:c11-classes} (iii)), and the condition on $D$ makes them $aQ^\xi_0$ plus commutators with $Q^\xi_0$ on the twisted letters. (b), (c) The pair rule [BK, (3.7)--(3.11)], the colour factors of Lemma \ref{lem:c10-colour} with two outputs, and the exactness of the antiholomorphic derivatives, $\{Q,\bar Q^1_{\dot\alpha}\}=P_{-\dot\alpha}$ [BK, (2.5)]. (d) Corollary \ref{cor:c8-homotopy} (ii) at $(\xi,R)$ with $[\xi,R]=0$, at order $\hbar$, mapped by $r_\xi$ (Lemma \ref{lem:c8-ev}). (e) $[Q^\xi_1]$ is determined by its values on representatives, polynomials in the letters, through the pair rule.
\end{proof}

\begin{remark}[Method of computation]
\label{rem:c11-oddmethod}
Each defect cycle is evaluated exactly as in Remark \ref{rem:c11-method}: the Gaussian loop integral and the integrals of \eqref{eq:c11-oddsimplex} over the ordered subsimplices of the region where $\alpha_\sigma$ is the smallest coordinate, the letter derivatives through Proposition \ref{prop:c11-genfun} with $n+1$ edges, the position along $\eta$ and $K_b$ through the injected momentum $r$, the fibre algebra in an explicit realisation of $V_{10}$ and $\Sigma$ over $\mathbb Q(i)$ with the vertices $S^{(1)}(x)$ of Chapter 2 and the sign of Definition \ref{def:c11-defect} (ii), letter vectors with fractional components contracted exactly, and the colour words by Lemma \ref{lem:c10-colour}; the twisted letters are built from the weights of a Cartan subalgebra. A second evaluation works in momentum space with an independently written realisation of $S^{(1)}(x)$ in the components of Definition \ref{def:c8-twocomp}, the special edge as the factor $e^{-\varepsilon|p|^2}-1$, Feynman parameters for the other edges and numerical integration at several values of $\varepsilon$ with thirty digits, extrapolated to $\varepsilon^0$ with the poles fitted; along $K_b$ it evaluates the cycles at the endpoints $T=1,4,16$. The two evaluations give the same operators $C_\xi$, $Y_\eta$, $C_K$ and the same arrows of Lemma \ref{lem:c11-lines}, the same composite line along $\xi$, the same free directions of Proposition \ref{prop:c11-freexi}, the same kernels $\Upsilon_\xi$ on every leg set with at most two outputs of the pairs of field and ghost letters without derivatives and with one letter derivative, the same coefficients of Theorem \ref{thm:c11-bilinears} for both chiralities, the same entries $\beta\gamma$, $\beta\beta$ and $bc$, and the same translation-invariant kernels along $\eta$ and $K_b$ where both are evaluated; both reproduce Example \ref{ex:c11-cubic}, Proposition \ref{prop:c11-plane} and the value along $P_b$ of Corollary \ref{cor:c11-DPdefect}. The first evaluation gives the values along $\Delta$ of Corollary \ref{cor:c11-DPdefect}, extends $\Upsilon_\xi$ to three letter derivatives, and gives the entries of Table \ref{tab:c11-twistentries} and the identity (iv) of Theorem \ref{thm:c11-bilinears}; the second gives the closed forms of Theorem \ref{thm:c11-oddsimplex} (iv) for $n_t=2,\dots,6$, the coefficients of $\varepsilon^{-m}$ of the defect cycles on the pairs of gaugino letters and on the pair of Example \ref{ex:c11-cubic}, against the right side of \eqref{eq:c11-pole}, the position identity of Theorem \ref{thm:c11-oddsimplex} (ii), and, along $K_b$, $T\partial_T$ of the defect cycles of the intrinsic vertex $S^{(1)}(K_b)$ on $(A|_V,A|_V)$ with the output $A|_V$, with the vertex on a line and at the special edge, decreasing at least as $\varepsilon$ at the three endpoints, as Theorem \ref{thm:c11-oddsimplex} (iii) states. The action-level obstructions of Proposition \ref{prop:c11-oddsol} and the arity-two classes of Theorem \ref{thm:c11-oddar2} are computed in the same two ways, on cycles of vertices of the action with one special edge, and the cohomology of $Q$ on complete bases of the two-leg functionals fixed by representation counts. Every integral is a four-dimensional Schwinger integral with the cutoff $\varepsilon$ on every edge.
\end{remark}

\begin{remark}[The data that complete the one-loop components]
\label{rem:c11-next}
The next computations are the following. On the action side: the arity-two equations along $(K_b,P_{b'})$ and $(\eta,\eta')$ at a finite endpoint, which contain the coefficients $\theta^K_d$ and $\kappa^K_j$ of Proposition \ref{prop:c11-oddsol} (v) and give the one-loop parts of $[K,P]=D+M$ and $\{S,S\}=K$, and along $(\xi_-,\eta_\pm)$; the three- and four-leg parts of the arity-one equations, which decide whether the sector parts of the generators extend to symmetries of the interacting theory, fix the differences of the sector coefficients (Proposition \ref{prop:c11-arity1} (iii), Proposition \ref{prop:c11-freexi} (ii)) and the couplings $\kappa_k$ of Theorem \ref{thm:c9-F3}; and the arity-three component of $s(\hbar F_1)=\hbar\,\mathrm{Ob}$. On the operator side: $Q_1$, $S_1$, $M_1$, $K_1$ and $P_1$ along generic directions on the families of Definition \ref{def:c11-families}; the arity-two operator kernels, with which the relations (R2)--(R6) of Proposition \ref{prop:c11-relations} become identities of local functionals before the projection to cohomology; and on the scalar bilinears the one-loop components on the level of Weyl weight $\tfrac52$, with which (R2) and (R3), the Konishi relation between $Q_1$, $S_1$ and $D_1=-48$ units, make the combinations of Corollary \ref{cor:c11-F0} (ii) scheme independent. The class-level comparison of Theorem \ref{thm:c11-twistcheck} beyond the chiral bilinears uses the transport map $\Phi$ built from [ESW, Thm 10.9] with the normalisation of the letters $b$, $\beta$, $\gamma$, $\partial_{\dot\alpha}c$ of [BK], the condition on $D$, the arity-two kernels along $(\xi,\xi)$, which give the entry $bc$ with $b$ represented by the self-dual field strength and the entry $bb$ on the pair $(A|_V,A|_V)$ with two derivatives, and the relations [BK, (3.59)--(3.63)] applied to the entries of Table \ref{tab:c11-twistentries}. The inner-product relations of Proposition \ref{prop:c11-S} (iv) and Corollary \ref{cor:c11-K} use $\mathrm{Gram}_{\rm 1}$ from the two-point graphs of Definition \ref{def:c8-2pt}.
\end{remark}

\begin{remark}[Two-component forms]
\label{rem:c11-twocomp}
With the components of Definition \ref{def:c8-twocomp}:
\begin{itemize}
\item Letters. $A|_V$, $A|_{W_{\mathbb R}^\perp}$, $\lambda|_{S_+\otimes W}$, $\lambda|_{S_-\otimes W^*}$ are $A_m$, $A_{ij}$ (the $\mathbf 6$ of $SU(4)_R$), $\lambda_\alpha{}^i$ (undotted, $\mathbf 4$), $\bar\lambda_{\dot\alpha i}$ (dotted, $\bar{\mathbf 4}$), with the antifields $A^*_m$, $A^*_{ij}$, $\lambda^{*\alpha}{}_i$, $\bar\lambda^{*\dot\alpha i}$, $c^*$ and the degrees $1,0,0,0,0,-1,-1,-1,-1,-2$ for $c,A_m,A_{ij},\lambda_\alpha{}^i,\bar\lambda_{\dot\alpha i},A^*_m,A^*_{ij},\lambda^{*\alpha}{}_i,\bar\lambda^{*\dot\alpha i},c^*$; $(\lambda,\lambda^*)=(\lambda_\alpha{}^i,\lambda^{*\alpha}{}_i)-(\bar\lambda_{\dot\alpha i},\bar\lambda^{*\dot\alpha i})$ and $(u,w)=\tfrac14u^{ij}w_{ij}$ on $W_{\mathbb R}^\perp$.
\item Composite kernels. The internal vertex along $\Delta$ is $\int\big((2E+2)A_m,A^*_m\big)=2\int x_n(\partial_nA_m,A^*_m)+2\int(A_m,A^*_m)$ with $E=x_n\partial_n$, along $P_b$ it is $\int(b_n\partial_nA_m,A^*_m)$; it sits between $\langle A_mA_n\rangle=2(-,-)^{-1}\delta_{mn}$ and $\langle A^*_m\,c\rangle=-(-,-)^{-1}ip_m$, each times $e^{-\varepsilon|p|^2}|p|^{-2}$ at $T=\infty$, and the insertion factors are $2$, $2i(l+q_c)_n$ and $ib_n(l+q_c)_n$. The special edge carries $e^{-\varepsilon|p|^2}-1$ times $(-,-)^{-1}$ and the inverse fibre pairing of $(A_m,A^*_m)$, $\tfrac14(A^{ij},A^*_{ij})$, $(\lambda,\lambda^*)$ and $(c,c^*)$: on $(\lambda_\alpha{}^i,\lambda^{*\beta}{}_j)$ the factor $-|l+q_\sigma|^2\delta_\alpha^\beta\delta^i_j(-,-)^{-1}$, on $(\bar\lambda_{\dot\alpha i},\bar\lambda^{*\dot\beta j})$ the opposite sign.
\item Bilinears. For $O=\tfrac1{16}\int\omega\,(w_1^{ij}A_{ij},w_2^{kl}A_{kl})$ the composite pair kernel along $\Delta$ is
\[
\Big[-\tfrac13\big((2E+2)A_{ij}\big)A_{kl}-\tfrac13A_{ij}\big((2E+2)A_{kl}\big)-\tfrac43A_{ij}A_{kl}\Big]\ \text{units},
\]
and along $P_b$ it is $-\tfrac13$ units times $b_m\partial_m(A_{ij}A_{kl})$, with the same contraction.
\item Two-site density. On adjacent pairs $(A_{ij},A_{kl})$, $(Z,\psi)$, $(\psi,Z)$, $(\psi,\psi)$ of a single trace the composite pair kernels give $-\tfrac43$ units (translation-invariant part along $\Delta$) and $-\tfrac16$ units times $b_m\partial_m$ (along $P_b$) on the same word, and $0$ on the exchanged word.
\item Free directions. The sector parts of $S^{(1)}(\Delta)$ are $\tfrac14\int\big((2E+2)A^{ij},A^*_{ij}\big)$, $-\int\big[\big((2E+3)\lambda_\alpha{}^i,\lambda^{*\alpha}{}_i\big)-\big((2E+3)\bar\lambda_{\dot\alpha i},\bar\lambda^{*\dot\alpha i}\big)\big]$ and $\int\big((2E+2)A_m,A^*_m\big)-\int\big(2E(c),c^*\big)$; those of $S^{(1)}(P_b)$ are $\tfrac14\int(b_m\partial_mA^{ij},A^*_{ij})$, $-\int\big[(b_m\partial_m\lambda_\alpha{}^i,\lambda^{*\alpha}{}_i)-(b_m\partial_m\bar\lambda_{\dot\alpha i},\bar\lambda^{*\dot\alpha i})\big]$ and $\int(b_m\partial_mA_n,A^*_n)-\int(b_m\partial_mc,c^*)$. The exact directions are $Q\int x_m(A_m,c^*)$ and $Q\int(\rho(x)\lambda^*,\lambda^*)=-2i\,Q\int x_m\big(\lambda^{*\alpha}{}_i,\sigma^m_{\alpha\dot\alpha}\bar\lambda^{*\dot\alpha i}\big)$. The further class of Proposition \ref{prop:c11-tc0} without the reflection is a multiple of $\int\big[(\lambda_\alpha{}^i,\lambda^{*\alpha}{}_i)+(\bar\lambda_{\dot\alpha i},\bar\lambda^{*\dot\alpha i})\big]$, the rotation $\lambda_\alpha{}^i\mapsto e^{i\varphi}\lambda_\alpha{}^i$, $\bar\lambda_{\dot\alpha i}\mapsto e^{-i\varphi}\bar\lambda_{\dot\alpha i}$ of the free gaugino.
\item Letter count. It counts $2$ for each letter $A_m$, $A_{ij}$, $\lambda_\alpha{}^i$, $\bar\lambda_{\dot\alpha i}$ and $1$ for each derivative $\partial_m$ on it; (b) reads: the composite pair kernels along $P_b$ sum to $-\tfrac13$ units times $b_m\partial_m$ on cohomology.
\item Families and spectra. $n_m=\sigma^m_{1\dot1}$, $D_n=n_m(\partial_m+[A_m,-])$, $\psi=\lambda_1{}^i$, and $F_{11}$ the self-dual field strength with both undotted indices $1$. In Srednicki's normalisation, $A_m=-igA^a_mT^a$, $A_I=-igA^a_IT^a$, $\lambda=-\tfrac{ig}{\sqrt2}\lambda^aT^a$, $\operatorname{tr}(T^aT^b)=\tfrac12\delta^{ab}$, the unit is $-\frac{g^2N}{32\pi^2}$ and $-\tfrac13$ units is $\frac{g^2N}{96\pi^2}$. On the families $D_1=-8$ units times the eigenvalues $E$ of $\sum_lH_{l,l+1}$, which is $\frac{g^2N}{4\pi^2}E=2\gamma$ with the one-loop anomalous dimension $\gamma=\frac{g^2N}{8\pi^2}E$: $\gamma=\frac{g^2N}{2\pi^2}h(S)$ on twist two, $\gamma=\frac{3g^2N}{4\pi^2}$ for $\operatorname{tr}(F_{11}F_{11})$ and the trace bilinear, the Konishi multiplet.
\item Odd generators. $Q_1(O)=\xi_\alpha{}^i\,Q_1(O)^\alpha{}_i+\bar\xi_{\dot\alpha i}\,\bar Q_1(O)^{\dot\alpha i}$ and $S_1(O)=\eta^\alpha{}_i\,S_1(O)_\alpha{}^i+\bar\eta^{\dot\alpha i}\,\bar S_1(O)_{\dot\alpha i}$ (Definition \ref{def:c8-components}); the Killing spinor of $\eta$ has the components $(\rho(x)\eta)_\alpha{}^i=-i\,x_m\sigma^m_{\alpha\dot\alpha}\bar\eta^{\dot\alpha i}$ and $(\rho(x)\eta)_{\dot\alpha i}=i\,x_m\eta^\alpha{}_i\sigma^m_{\alpha\dot\alpha}$ \eqref{eq:c8-pairing}, so on a family with $P_1\cong0$ the position part along $b$ of $S_1$ along $\eta$ is, on cohomology, $Q_1$ along $\xi_\alpha{}^i=-i\,b_m\sigma^m_{\alpha\dot\alpha}\bar\eta^{\dot\alpha i}$, $\bar\xi_{\dot\alpha i}=i\,b_m\eta^\alpha{}_i\sigma^m_{\alpha\dot\alpha}$. For $\xi_\alpha{}^i=u_\alpha z^i$, $(\rho(w)\xi)^\alpha{}_j=\varepsilon^{\alpha\beta}w_{ji}u_\beta z^i$, so $\rho(w)\xi=0$ exactly when $w_{ij}z^j=0$, and then $\iota_wA=\tfrac14w^{ij}A_{ij}$ is a combination of the components $A_{ij}z^j$: Proposition \ref{prop:c11-plane} is the vanishing of $Q_1$ on $\int\omega\,(m_1^iA_{ij}z^j,m_2^kA_{kl}z^l)$, $m_1,m_2\in W$.
\item Terms of $S^{(1)}(\xi)$. For $\xi=(\xi_\alpha{}^i,0)$, with $(\xi^j\lambda^i):=\xi^{\alpha j}\lambda_\alpha{}^i$ and $F^0_{mn}=\partial_mA_n-\partial_nA_m$: $T_1=-i\big(\bar\lambda_{\dot\alpha i}\bar\sigma^{m\dot\alpha\alpha}\xi_\alpha{}^i,A^*_m\big)$, $T_2=\tfrac12\varepsilon_{jikl}\big((\xi^j\lambda^i),A^{*kl}\big)$, $T_3=-\tfrac14\big(F^0_{mn}(\sigma^m\bar\sigma^n)_\alpha{}^\beta\xi_\beta{}^i,\lambda^{*\alpha}{}_i\big)$, $T_4=-\tfrac i2\big(\xi^{\alpha j}\sigma^m_{\alpha\dot\alpha}\partial_mA_{ij},\bar\lambda^{*\dot\alpha i}\big)$, and $T_5$ the three terms $-\tfrac14\big([A_m,A_n](\sigma^m\bar\sigma^n)_\alpha{}^\beta\xi_\beta{}^i,\lambda^{*\alpha}{}_i\big)$, $-\tfrac i2\big(\xi^{\alpha j}\sigma^m_{\alpha\dot\alpha}[A_m,A_{ij}],\bar\lambda^{*\dot\alpha i}\big)$, $-\tfrac14\big([A^{ik},A_{kj}]\xi_\alpha{}^j,\lambda^{*\alpha}{}_i\big)$; for $\xi_\alpha{}^i=u_\alpha z^i$ these are the terms of the remark after Theorem \ref{thm:c8-twist}. The terms with $\bar\xi_{\dot\alpha i}$ follow by $\lambda_\alpha{}^i\leftrightarrow\bar\lambda_{\dot\alpha i}$, $\lambda^{*\alpha}{}_i\leftrightarrow\bar\lambda^{*\dot\alpha i}$, $\sigma\leftrightarrow\bar\sigma$. The intrinsic term along $\eta$ is $s_0=\big(A^{kj}\varepsilon_{\alpha\beta}\eta^\beta{}_j,\lambda^{*\alpha}{}_k\big)-\big(A_{kj}\varepsilon_{\dot\beta\dot\alpha}\bar\eta^{\dot\alpha j},\bar\lambda^{*\dot\beta k}\big)$, and $S^{(1)}(K_b)=\int\big\{(v.A_m,A^*_m)+\tfrac14(v.A^{ij},A^*_{ij})-(v.\lambda_\alpha{}^i,\lambda^{*\alpha}{}_i)+(v.\bar\lambda_{\dot\alpha i},\bar\lambda^{*\dot\alpha i})-(v.c,c^*)\big\}$ with $v^m=x^2b^m-2(b\cdot x)x^m$, $v.A_m=v^n\partial_nA_m+A_n\partial_mv^n$, $v.A_{ij}=v^n\partial_nA_{ij}-2(b\cdot x)A_{ij}$, $v.\lambda_\alpha{}^i=v^n\partial_n\lambda_\alpha{}^i-\tfrac12(x_mb_n-x_nb_m)(\sigma^m\bar\sigma^n)_\alpha{}^\beta\lambda_\beta{}^i-3(b\cdot x)\lambda_\alpha{}^i$, $v.c=v^n\partial_nc$.
\item Arrows and lines. $\xi.$ maps $A_{kl}\mapsto\Gamma_{kl}(\xi,\lambda)=-2\varepsilon_{jikl}(\xi^j\lambda^i)$ and $A_m\mapsto\Gamma^m(\xi,\lambda)=i\big(\bar\lambda_{\dot\alpha i}\bar\sigma^{m\dot\alpha\alpha}\xi_\alpha{}^i+\bar\xi_{\dot\alpha i}\bar\sigma^{m\dot\alpha\alpha}\lambda_\alpha{}^i\big)$. On a plane wave $C_\xi$ maps $\lambda$ to $c$ by $\lambda\mapsto-p_m\big(\bar\lambda_{\dot\alpha i}\bar\sigma^{m\dot\alpha\alpha}\xi_\alpha{}^i+\bar\xi_{\dot\alpha i}\bar\sigma^{m\dot\alpha\alpha}\lambda_\alpha{}^i\big)$. The composite line along $\xi_\alpha{}^i$ joins a leg $\lambda_\alpha{}^i$ to a leg $c$ through $\langle\lambda_\alpha{}^i\bar\lambda_{\dot\beta j}\rangle$, the term $T_1$ and $\langle A^*_m\,c\rangle$; with $\sigma^k_{\alpha\dot\beta}\bar\sigma^{m\dot\beta\gamma}p_kp_m=-|p|^2\delta_\alpha^\gamma$ it is $-2\,\xi_\alpha{}^i\,e^{-2\varepsilon|p|^2}|p|^{-2}(-,-)^{-1}$ on $c\otimes\lambda_\alpha{}^i$ and $+2\,\xi_\alpha{}^i\,e^{-2\varepsilon|p|^2}|p|^{-2}(-,-)^{-1}$ on $\lambda_\alpha{}^i\otimes c$. Along $\eta$, $(\rho(e_m)\eta)_\alpha{}^i=-i\sigma^m_{\alpha\dot\alpha}\bar\eta^{\dot\alpha i}$ and $(\rho(e_m)\eta)_{\dot\alpha i}=i\eta^\alpha{}_i\sigma^m_{\alpha\dot\alpha}$; along $K_b$, $\operatorname{div}v=-8\,b\cdot x$, $\Sigma^{\mu\nu}$ acts on $A_m$ as $(\Sigma^{\mu\nu}A)_m=\delta_{m\mu}A_\nu-\delta_{m\nu}A_\mu$ and on $\lambda_\alpha{}^i$ through $-\tfrac14(\sigma^\mu\bar\sigma^\nu-\sigma^\nu\bar\sigma^\mu)$, and $C_K$ maps $A_m$ to $-4b_mA_m$ in the slot of $c$.
\item Bilinears along $\xi$. For $O=\tfrac1{16}\int\omega\,(w_1^{ij}A_{ij},w_2^{kl}A_{kl})$,
\[
\mathcal B_1=\tfrac14\int\omega\,\big[(w_1^{kl}A_{kl},w_{2\,ji}\xi^{\alpha i}\lambda_\alpha{}^j)+(1\leftrightarrow2)\big],\qquad \mathcal B_2=\tfrac14w_1^{ij}w_{2\,ij}\int\omega\,(A_{ji},\xi^{\alpha i}\lambda_\alpha{}^j),
\]
$\mathcal B_3=\tfrac14w_1^{ij}w_{2\,ij}\int\omega\,i\big(\bar\lambda_{\dot\alpha i},A_m\bar\sigma^{m\dot\alpha\alpha}\xi_\alpha{}^i\big)$ and $\mathcal B_4=\tfrac14w_1^{ij}w_{2\,ij}\int\omega\,(c,\xi_\alpha{}^i\lambda^{*\alpha}{}_i)$; Theorem \ref{thm:c11-bilinears} reads, in units $-\frac{g^2N}{32\pi^2}$,
\[
\begin{aligned}
&Q_1(O)\big|_{\text{two-body}}=-2\log\tfrac32\,\mathcal B_1+4\log\tfrac32\,(\mathcal B_2+\mathcal B_3),\\
&L_1^{(0)}(O)=-\big(2+4\log\tfrac32\big)\tfrac14w_1^{ij}w_{2\,ij}\int\omega\,(\partial_mc,A_m),
\end{aligned}
\]
and the cocycle finite part is $\psi(O)=-\big(1+2\log\tfrac32\big)\tfrac14w_1^{ij}w_{2\,ij}\int\omega\,(A_m,A_m)$ units, $\psi(\mathcal B_1)=-2\,\mathcal B_3$ units. On the twenty traceless bilinears the transformation along $Q^i_\alpha$ is $1+c_Q$ times the classical one, with $c_Q$ of Theorem \ref{thm:c11-bilinears} (ii).
\item Special supercharges and special conformal fields on the bilinears. The two-body part of $S_1(O)$ along $\eta^\alpha{}_i$ is that of $\bar Q_1(O)$ along $\bar\xi_{\dot\alpha i}=i\,x_m\eta^\alpha{}_i\sigma^m_{\alpha\dot\alpha}$, and along $\bar\eta^{\dot\alpha i}$ that of $Q_1(O)$ along $\xi_\alpha{}^i=-i\,x_m\sigma^m_{\alpha\dot\alpha}\bar\eta^{\dot\alpha i}$. On the part with two scalar legs, $K_1(O)=-\big(\tfrac13+\theta\big)\big(v^m\partial_mO-4(b\cdot x)O\big)$ units on the traceless bilinears and the same plus $48\,(b\cdot x)\,O$ units on the trace $\tfrac14\int\omega\,(A^{ij},A_{ij})$.
\item Twist. With $z=e_4$: $\gamma^I=A_{I4}$, $\beta_I=u'^\alpha\lambda_\alpha{}^I$ ($I=1,2,3$), $\partial_{\dot\alpha}c=\bar\lambda_{\dot\alpha4}$, $b=u'^\alpha F_{mn}(\sigma^m\bar\sigma^n)_\alpha{}^\beta u'_\beta$ or $(\xi,\lambda^*)=u_\alpha\lambda^{*\alpha}{}_4$, and $\partial_{\dot\alpha}=u'^\alpha\sigma^m_{\alpha\dot\alpha}D_m$, so that $\partial_2\gamma^J\partial_1\gamma^K-\partial_1\gamma^J\partial_2\gamma^K$ of Table \ref{tab:c11-twistentries} is a multiple of $\partial_{\dot\alpha}\gamma^J\partial^{\dot\alpha}\gamma^K$. The Konishi entry is the coefficient of $\int\omega\,(\bar\lambda_{\dot\alpha4},\bar\lambda^{\dot\alpha}{}_4)$ in $Q_1\big(\int\omega\,(u'^\alpha\lambda_\alpha{}^i,A_{i4})\big)$, $0$ at $\psi=0$; [BK, (3.66)] is $Q_1\big((\beta_I)^A(\gamma^J)^B\big)=\tfrac12\delta^J_If_{ACD}f_{BCE}\,\partial_{\dot\alpha}c^D\partial^{\dot\alpha}c^E$, and the entries $\beta\beta$, $b\gamma$, $\beta b$ of [BK, (3.67)--(3.69)] have the same factor $\tfrac12f_{ACD}f_{BCE}$ in both terms.
\item Action-side data along $\Sigma$ and $\Sigma^*$. For $\xi=(\xi_\alpha{}^i,0)$, with $\Gamma_V^m=i\,\bar\lambda_{\dot\alpha i}\bar\sigma^{m\dot\alpha\alpha}\xi_\alpha{}^i$ and $(\Gamma_\perp)_{kl}=2\varepsilon_{jikl}\,\xi_\alpha{}^j\varepsilon^{\alpha\beta}\lambda_\beta{}^i$, the functionals of Proposition \ref{prop:c11-oddsol} are $g_1=\int d^4x\,(\Gamma_V^n,-\partial^2A_n+\partial_n\partial_mA_m)$, $g_2=-\int d^4x\,(\Gamma_V^n,\partial_n\partial_mA_m)$, $g_3=-\tfrac14\int d^4x\,((\Gamma_\perp)_{kl},\partial^2A^{kl})$, $g_4=-\int d^4x\,(\partial^2c,\xi_\alpha{}^i\lambda^{*\alpha}{}_i)$ and $T_0=-\int d^4x\,(\partial_mA_m,\xi_\alpha{}^i\lambda^{*\alpha}{}_i)$, and the solution along $\xi$ is $\big[(\log\tfrac32-\tfrac{79}{12})(T_1+T_2)+\theta^\xi_V(T_1+T_3)+\theta^\xi_\perp(T_2+T_4)\big]$ units with the terms of the bullet above. For $\eta=(\eta^\alpha{}_i,0)$, $s_0=\int d^4x\,(A^{kj}\varepsilon_{\alpha\beta}\eta^\beta{}_j,\lambda^{*\alpha}{}_k)$ and $u_0=-i\int d^4x\,(A_m\eta^\alpha{}_i\sigma^m_{\alpha\dot\alpha},\bar\lambda^{*\dot\alpha i})$. Along $M_\omega$ and $K_b$ the products $\rho_\mu\rho_\nu$ of the spin terms are $\sigma^\mu\bar\sigma^\nu$ on the undotted and $\bar\sigma^\mu\sigma^\nu$ on the dotted components, and $(\iota_bA,\iota_xA^*)=(b_mA_m,x_nA^*_n)$.
\item Relations among the coefficients. \eqref{eq:c11-trel} reads $t_{Q}+t_{\bar Q}-t_P=t_Q+t_{S}-t_\Delta=\tfrac{27}4+2\log\tfrac32$ for the coefficients along $Q^i_\alpha$ ($\xi_\alpha{}^i$), $\bar Q_{\dot\alpha i}$ ($\bar\xi_{\dot\alpha i}$), $S^\alpha{}_i$ ($\eta^\alpha{}_i$) and $P_m$, with $t_\Delta=t_M=-\tfrac13$ and $t_R=0$; the outer direction moves $t_Q$, $t_{\bar S}$ by $+a$ and $t_{\bar Q}$, $t_S$ by $-a$.
\item Components on the bilinears. In units $-\frac{g^2N}{32\pi^2}$, $Q_1(O)=\big(\tfrac{79}{12}-3\log\tfrac32-t_Q\big)\mathcal B_1+4\log\tfrac32\,(\mathcal B_2+\mathcal B_3)$ along $\xi_\alpha{}^i$, $S_1(O)=\big(\tfrac{79}{12}-3\log\tfrac32-t_S\big)\{S^{(1)}(\eta),O\}$ plus the two-body terms $4\log\tfrac32\,(\mathcal B_2+\mathcal B_3)$ at the spinor $\bar\xi_{\dot\alpha i}=i\,x_m\eta^\alpha{}_i\sigma^m_{\alpha\dot\alpha}$ along $\eta^\alpha{}_i$, $D_1=0$ and $-48$ units, $P_1(O)=-(\tfrac13+t_P)\,b_m\partial_mO$, $M_1(O)=-(\tfrac13+t_M)\,\omega_{mn}x_n\partial_mO$, $R_1=0$, and $K_1(O)=-(\tfrac13+\theta^K_\perp)\big(v^m\partial_mO-4(b\cdot x)O\big)$ plus $48\,(b\cdot x)\,O$ on the trace and the brackets with $G^K_j$; the invariant combination of Corollary \ref{cor:c11-F0} (ii) is $\tfrac{27}4-8\log\tfrac32$ units.
\item Weinberg. On the Dirac column $(\lambda_\alpha{}^i,\bar\lambda^{\dot\alpha}{}_i)$, $\rho(v)=\gamma_W^\mu v_\mu$ gives Weinberg's form of each line, as in Remark \ref{rem:c10-twocomp}.
\end{itemize}
\end{remark}
